% =====================================================================
%  CHƯƠNG 4 — HIỆN THỰC HOÁ HỆ THỐNG VÀ CÁC TÍNH NĂNG THÔNG MINH
%
%  File chương theo cấu trúc template Overleaf (hướng Ứng dụng).
%  Cách dùng: đặt file này trong thư mục Chuong/ của template và thêm
%      % =====================================================================
%  CHƯƠNG 4 — HIỆN THỰC HOÁ HỆ THỐNG VÀ CÁC TÍNH NĂNG THÔNG MINH
%
%  File chương theo cấu trúc template Overleaf (hướng Ứng dụng).
%  Cách dùng: đặt file này trong thư mục Chuong/ của template và thêm
%      % =====================================================================
%  CHƯƠNG 4 — HIỆN THỰC HOÁ HỆ THỐNG VÀ CÁC TÍNH NĂNG THÔNG MINH
%
%  File chương theo cấu trúc template Overleaf (hướng Ứng dụng).
%  Cách dùng: đặt file này trong thư mục Chuong/ của template và thêm
%      % =====================================================================
%  CHƯƠNG 4 — HIỆN THỰC HOÁ HỆ THỐNG VÀ CÁC TÍNH NĂNG THÔNG MINH
%
%  File chương theo cấu trúc template Overleaf (hướng Ứng dụng).
%  Cách dùng: đặt file này trong thư mục Chuong/ của template và thêm
%      \include{Chuong/Chuong4}
%  vào main.tex (sau Chương 3).
%
%  GHI CHÚ VỀ HÌNH VẼ: dùng macro \figph{<tên-file>}{<tỉ-lệ-rộng>}{<chú-thích>}{<nhãn>}
%  (đã khai báo ở các chương trước; khai báo lại bằng \providecommand để an
%  toàn nếu biên dịch độc lập). Mã nguồn (code) trình bày bằng môi trường
%  verbatim như các chương trước.
%
%  KHÔNG hardcode số chương/mục/bảng/hình/công thức — luôn dùng \label + \ref / \cite.
% =====================================================================

% --- Macro tiện ích (idempotent — không định nghĩa lại nếu đã có) ------
\providecommand{\code}[1]{\texttt{\detokenize{#1}}} % in định danh snake_case an toàn
\providecommand{\figph}[4]{%
  \begin{figure}[htbp]
    \centering
    \IfFileExists{images/#1}%
      {\includegraphics[width=#2\textwidth,height=0.82\textheight,keepaspectratio]{#1}}%
      {\setlength{\fboxsep}{10pt}\fbox{\begin{minipage}[c][3.6cm][c]{#2\textwidth}\centering\itshape Sơ đồ minh hoạ --- chèn tệp ảnh \code{images/#1} vào thư mục \code{images/}.\end{minipage}}}%
    \caption{#3}%
    \label{#4}%
  \end{figure}}

\chapter{HIỆN THỰC HOÁ HỆ THỐNG VÀ CÁC TÍNH NĂNG THÔNG MINH}
\label{chuong:hien-thuc-hoa}

Hai chương trước đã trình bày kết quả phân tích, thiết kế
(Chương~\ref{chuong:phan-tich-thiet-ke}) và cơ sở lý thuyết, công nghệ
(Chương~\ref{chuong:co-so-ly-thuyet}). Chương này trình bày cách các thiết
kế đó được \textit{hiện thực hoá} thành mã nguồn chạy được, với trọng tâm là
các \textbf{thành phần thông minh và tự động hoá} làm nên giá trị cốt lõi của
hệ thống: động cơ học theo ôn tập ngắt quãng, cơ chế sinh câu hỏi và chấm
điểm thích nghi, các tính năng làm giàu nội dung và chấm điểm dựa trên trí
tuệ nhân tạo, cùng hạ tầng xử lý nền bất đồng bộ liên kết chúng lại.

Mục~\ref{sec:tong-quan-hien-thuc} mô tả tổ chức mã nguồn và các mối quan tâm
xuyên suốt. Mục~\ref{sec:hien-thuc-du-lieu} và
mục~\ref{sec:hien-thuc-auth} trình bày ngắn gọn hiện thực hoá tầng dữ liệu và
phân hệ xác thực. Trọng tâm của chương nằm ở
mục~\ref{sec:dong-co-hoc} (động cơ học thông minh) và
mục~\ref{sec:ai-tu-dong} (các tính năng AI và tự động hoá nội dung). Cuối
cùng, mục~\ref{sec:tien-do} trình bày phân hệ theo dõi tiến độ tự động,
mục~\ref{sec:hien-thuc-frontend} trình bày hiện thực hoá tầng giao diện người
dùng (frontend), mục~\ref{sec:kiem-thu} trình bày chiến lược kiểm thử và
mục~\ref{sec:ket-chuong-4} tóm tắt chương.

% =====================================================================
\section{Tổng quan quá trình hiện thực hoá}
\label{sec:tong-quan-hien-thuc}

Hệ thống được hiện thực hoá đúng theo ngăn xếp công nghệ đã chọn ở
Chương~\ref{chuong:co-so-ly-thuyet}: backend viết bằng \textbf{NestJS~11}
trên \textbf{TypeScript} (chế độ \textit{strict}), truy cập
\textbf{PostgreSQL} qua \textbf{TypeORM}, và dùng \textbf{Redis/BullMQ} cho
xử lý nền. Toàn bộ mã nguồn được biên dịch bằng trình biên dịch của NestJS và
khởi chạy dưới \textbf{hai tiến trình} chia sẻ chung cơ sở dữ liệu và hàng
đợi, đúng mô hình nhà sản xuất--người tiêu thụ đã nêu ở
mục~\ref{sec:kien-truc} của Chương~\ref{chuong:phan-tich-thiet-ke}:

\begin{itemize}
  \item \textbf{Tiến trình API} (\code{npm run start}): phục vụ các yêu cầu
        HTTP đồng bộ; với tác vụ nặng nó chỉ tạo bản ghi công việc rồi đẩy
        vào hàng đợi và trả về ngay.
  \item \textbf{Tiến trình Worker} (\code{npm run start:worker}): tiêu thụ
        hàng đợi để thực thi các tác vụ nền (làm giàu từ vựng, sinh âm thanh,
        chấm điểm luyện tập).
\end{itemize}

Mọi tham số vận hành (chuỗi kết nối CSDL, khoá bí mật, khoá API của dịch vụ
AI, hạn mức tần suất\ldots) đều được nạp từ \textbf{biến môi trường} qua
\code{@nestjs/config}, đáp ứng yêu cầu tính khả chuyển NFR-08 trong
Bảng~\ref{tab:nfr}: cùng một ảnh mã nguồn chạy được ở nhiều môi trường mà
không cần sửa code.

\subsection{Tổ chức mã nguồn theo mô-đun}
\label{subsec:to-chuc-ma-nguon}

Mỗi mô-đun nghiệp vụ ở Bảng~\ref{tab:modules} (Chương~\ref{chuong:phan-tich-thiet-ke})
được hiện thực thành một thư mục dưới \code{src/} tự đóng gói controller,
service, entity và DTO của mình. Bảng~\ref{tab:source-layout} ánh xạ các
mô-đun trọng yếu tới những tệp hiện thực then chốt --- là cơ sở để các mục
sau tham chiếu trực tiếp.

\begin{table}[h]
  \centering
  \small
  \renewcommand{\arraystretch}{1.3}
  \begin{tabular}{|p{2.6cm}|p{4.4cm}|p{6.6cm}|}
    \hline
    \textbf{Mô-đun} & \textbf{Thư mục mã nguồn} & \textbf{Tệp hiện thực then chốt} \\
    \hline
    \code{auth} & \code{src/auth} & \code{token.service.ts} (cấp/xoay token), \code{jwt.strategy.ts}, các \code{guards/}, \code{services/} (Google/Apple/GitHub). \\
    \hline
    \code{learn} & \code{src/learn} & \code{learn.service.ts}, \code{vocab-picker.service.ts}, \code{question-builder.service.ts}, \code{distractor.service.ts}, \code{answer-grader.service.ts}, \code{hmac-signer.service.ts}, \code{question-bands.ts}. \\
    \hline
    \code{progress} & \code{src/progress} & \code{srs.ts} (thuật toán SM-2), \code{progress.service.ts} (lập lịch, thống kê, chuỗi ngày học). \\
    \hline
    \code{vocabularies} & \code{src/vocabularies} & \code{enrichment/} (làm giàu bằng LLM, nhập hàng loạt), \code{audio/}, \code{images/}, \code{vocabulary-persistence.service.ts}. \\
    \hline
    \code{practice} & \code{src/practice} & \code{gemma-judge.ts} (chấm câu bằng LLM), \code{scoring.processor.ts}, \code{practice.service.ts}. \\
    \hline
    \code{pronunciation} & \code{src/pronunciation} & \code{pronunciation-scoring.client.ts} (proxy tới vi dịch vụ), \code{pronunciation.service.ts}. \\
    \hline
    \code{leaderboard} & \code{src/leaderboard} & \code{leaderboard.service.ts}. \\
    \hline
    \code{config} & \code{src/config} & Các tệp cấu hình theo miền: \code{learn.config.ts}, \code{gemma.config.ts}, \code{auth.config.ts}\ldots \\
    \hline
  \end{tabular}
  \caption{Ánh xạ mô-đun nghiệp vụ tới tệp hiện thực then chốt}
  \label{tab:source-layout}
\end{table}

\subsection{Các mối quan tâm xuyên suốt}
\label{subsec:cross-cutting}

Nhờ cơ chế \textit{tiêm phụ thuộc} và các thành phần khai báo của NestJS, các
mối quan tâm chung được hiện thực một lần và áp dụng toàn cục thay vì lặp lại
trong từng controller:

\begin{itemize}
  \item \textbf{Kiểm tra đầu vào:} một \code{ValidationPipe} toàn cục với
        \code{whitelist}, \code{forbidNonWhitelisted} và \code{transform}
        biến mỗi DTO (khai báo bằng \code{class-validator}) thành lớp bảo vệ:
        trường lạ bị loại bỏ, dữ liệu sai kiểu bị chặn bằng lỗi 400 trước khi
        chạm tới service.
  \item \textbf{Xác thực và phân quyền:} \code{JwtAuthGuard} (dựa trên
        \code{passport-jwt}, xem \code{jwt.strategy.ts}) trích và xác minh
        access token ở mọi endpoint cần đăng nhập; \code{RolesGuard} kết hợp
        decorator \code{@Roles('admin')} chặn các thao tác quản trị; quyền sở
        hữu tài nguyên được kiểm tra ngay trong service (mục~\ref{subsec:rbac}).
  \item \textbf{Phiên bản hoá API:} bật \textit{URI versioning} mặc định
        \code{v1}; controller khai báo \code{@Controller(\{ path, version: '1' \})}
        nên mọi tài nguyên nằm dưới tiền tố \code{/v1} (mục~\ref{subsec:rest-theory}).
  \item \textbf{Cấu hình theo miền:} mỗi nhóm tham số được gói trong một
        \textit{namespace} cấu hình riêng (ví dụ \code{learn.config.ts},
        \code{gemma.config.ts}), tách bạch và có giá trị mặc định an toàn cho
        môi trường phát triển.
\end{itemize}

% =====================================================================
\section{Hiện thực hoá tầng dữ liệu}
\label{sec:hien-thuc-du-lieu}

Tầng dữ liệu hiện thực 17 bảng trong Bảng~\ref{tab:db-overview} thành các lớp
\textit{thực thể} (entity) TypeORM. Mỗi thực thể khai báo khoá chính
\code{uuid}, tên cột \code{snake_case} và kiểu \code{timestamptz} cho các mốc
thời gian, đúng quy ước thiết kế ở Chương~\ref{chuong:phan-tich-thiet-ke}.
Một số kỹ thuật hiện thực đáng chú ý:

\begin{itemize}
  \item \textbf{Giao dịch bảo đảm tính nguyên tử:} các thao tác ghi nhiều
        bảng được gói trong \code{dataSource.transaction(...)}. Ví dụ tiêu
        biểu là việc nộp một lượt ôn (mục~\ref{subsec:impl-hmac}): bản ghi
        tiến độ và bản ghi nhật ký hoạt động được lưu trong cùng một giao
        dịch, nên một lượt ôn không thể "ghi nửa vời" --- đáp ứng NFR-04 trong
        Bảng~\ref{tab:nfr}.
  \item \textbf{Ràng buộc toàn vẹn và chống trùng lặp:} ngoài khoá ngoại và
        quy tắc xoá lan truyền, hệ thống dùng \textit{chỉ mục duy nhất bộ
        phận} (partial unique index) trên bộ ba \code{(language, lemma,
        part_of_speech)} áp dụng riêng cho từ thuộc kho hệ thống
        (\code{WHERE source = 'system'}); nhờ đó kho hệ thống không bao giờ có
        hai bản ghi trùng từ loại cho cùng một từ, trong khi từ cá nhân của
        người dùng vẫn không bị ràng buộc này.
  \item \textbf{Quản lý lược đồ bằng migration:} mọi thay đổi lược đồ được mô
        tả bằng các \textit{migration} có thể phiên bản hoá, được nhận diện và
        chạy qua công cụ dòng lệnh của TypeORM. Cách này cho phép tái lập lược
        đồ một cách xác định giữa các môi trường (đáp ứng NFR-06).
  \item \textbf{Phi chuẩn hoá có kiểm soát:} một số trường được phi chuẩn hoá
        để tăng tốc đọc, ví dụ \code{decks.vocab_count} (số từ trong bộ thẻ)
        được cập nhật khi thêm/bớt thành viên thay vì đếm mỗi lần truy vấn.
\end{itemize}

% =====================================================================
\section{Hiện thực hoá xác thực và phân quyền}
\label{sec:hien-thuc-auth}

Phân hệ \code{auth} hiện thực các kỹ thuật bảo mật đã trình bày về lý thuyết
ở mục~\ref{sec:xac-thuc-phan-quyen}. Lớp \code{TokenService} đảm nhiệm vòng
đời phiên đăng nhập với một số quyết định hiện thực đáng chú ý:

\begin{itemize}
  \item \textbf{Cặp token bất đối xứng về bản chất:} \textit{access token} là
        một \textbf{JWT} được ký, tự chứa và sống ngắn (mục~\ref{subsec:jwt}),
        trong khi \textit{refresh token} là một chuỗi ngẫu nhiên 48 byte
        \textit{mờ} (opaque) --- không mang thông tin, chỉ lưu \textbf{bản băm
        SHA-256} trong bảng \code{refresh_tokens}. Máy chủ không bao giờ lưu
        token gốc, nên dù cơ sở dữ liệu bị lộ thì token cũng không thể tái sử
        dụng.
  \item \textbf{Xoay vòng và phát hiện tái dùng:} mỗi lần làm mới phiên, token
        cũ bị đánh dấu thu hồi và một token mới được cấp. Nếu một token đã thu
        hồi (hoặc không tồn tại) được trình lại --- dấu hiệu rò rỉ --- hệ thống
        \textbf{thu hồi toàn bộ} token còn hiệu lực của người dùng đó và trả
        về lỗi 401, hạn chế thiệt hại khi token bị đánh cắp.
  \item \textbf{Băm mật khẩu và đăng nhập bên thứ ba:} mật khẩu được băm bằng
        \textbf{bcrypt} (mục~\ref{subsec:bcrypt}); với OAuth/OIDC
        (mục~\ref{subsec:oauth}), token danh tính do nhà cung cấp cấp được xác
        minh rồi liên kết tới một bản ghi trong \code{user_identities}, cho
        phép một tài khoản đăng nhập bằng nhiều phương thức.
  \item \textbf{Chống vét cạn:} endpoint đăng nhập được giới hạn tần suất bằng
        \code{@nestjs/throttler}, trả về lỗi 429 khi vượt ngưỡng.
\end{itemize}

% =====================================================================
\section{Động cơ học thông minh theo ôn tập ngắt quãng}
\label{sec:dong-co-hoc}

Đây là phân hệ \textit{thông minh} cốt lõi của hệ thống. Một phiên học không
phải là một danh sách câu hỏi tĩnh: với mỗi người học và mỗi từ, hệ thống
\textbf{thích nghi} ở ba cấp độ --- chọn \textit{từ nào} để học, sinh
\textit{dạng câu hỏi nào} phù hợp giai đoạn ghi nhớ, và quyết định
\textit{khi nào} ôn lại. Toàn bộ quy trình tuân theo luồng đã mô hình hoá ở
Hình~\ref{fig:activity-srs}: chọn từ $\rightarrow$ mở rộng thành chuỗi câu hỏi
đã ký $\rightarrow$ chấm điểm $\rightarrow$ (ở bước cuối) lập lịch lại.

\subsection{Lập lịch ôn tập SM-2 mở rộng bước học}
\label{subsec:impl-sm2}

Hạt nhân lập lịch là hàm thuần \code{applySm2} trong \code{src/progress/srs.ts}:
nó nhận trạng thái hiện tại của một cặp (người dùng, từ), một
\textit{điểm chất lượng} \( q \in \{0,\dots,5\} \) và danh sách bước học, rồi
trả về trạng thái mới cùng thời điểm ôn kế tiếp. Hàm hiện thực thuật toán
\textbf{SM-2} (mục~\ref{subsec:sm2-theory}) được \textit{mở rộng bằng các bước
học} theo phong cách Anki. Hệ số dễ nhớ \( EF \) được cập nhật ở \textit{mọi}
lượt ôn theo công thức~\eqref{eq:sm2-ef}, còn khoảng cách ngày tuân theo
công thức~\eqref{eq:sm2-interval}; phần mở rộng bổ sung một thang
\textit{bước học} tính bằng phút để củng cố từ mới ngay trong phiên trước khi
giao cho thang ngày:

\begin{verbatim}
// src/progress/srs.ts — cập nhật hệ số dễ nhớ (chặn dưới 1.3)
function updateEaseFactor(easeFactor, quality) {
  return Math.max(
    1.3,
    easeFactor + (0.1 - (5 - quality) * (0.08 + (5 - quality) * 0.02)),
  );
}
\end{verbatim}

Logic chuyển trạng thái được tóm tắt như sau (Hình~\ref{fig:srs-states}):

\begin{itemize}
  \item Thẻ đang ở \textit{trạng thái bước} (\code{learningStepIndex} khác
        \code{null}) dùng khoảng cách \textbf{phút} lấy từ danh sách bước (mặc
        định \([1, 10]\) phút). Trả lời đúng (\( q \ge 3 \)) đẩy lên một bước;
        qua bước cuối, thẻ \textit{tốt nghiệp} sang thang ngày. Trả lời sai
        (\( q < 3 \)) đặt lại về bước 0.
  \item Thẻ đã tốt nghiệp dùng thang ngày: \( I = 1 \) ngày (lần đầu), \( 6 \)
        ngày (lần hai), sau đó \( I \leftarrow \lceil I \times EF \rceil \).
        Một lần trả lời sai sẽ đưa thẻ \textit{rớt} về bước 0 (ôn lại trong
        phút) thay vì chờ tới hôm sau.
  \item Trạng thái hiển thị được suy ra từ số lần lặp và khoảng cách:
        \( I \ge 90 \) ngày $\rightarrow$ \code{mastered}; số lần ôn đúng liên
        tiếp \( \ge 3 \) $\rightarrow$ \code{review}; còn lại \code{learning}.
\end{itemize}

Bảng~\ref{tab:srs-params} liệt kê các tham số then chốt và giá trị mặc định
của chúng (đều có thể cấu hình qua biến môi trường trong \code{learn.config.ts}).

\begin{table}[h]
  \centering
  \small
  \renewcommand{\arraystretch}{1.3}
  \begin{tabular}{|p{5.2cm}|p{3.0cm}|p{6.0cm}|}
    \hline
    \textbf{Tham số} & \textbf{Giá trị mặc định} & \textbf{Ý nghĩa} \\
    \hline
    Hệ số dễ nhớ khởi tạo \( EF_0 \) & 2.5 & Giá trị ban đầu của \( EF \). \\
    \hline
    Chặn dưới \( EF_{\min} \) & 1.3 & Ngăn khoảng ôn co lại quá nhanh. \\
    \hline
    Bước học (phút) & \([1, 10]\) & Khoảng cách trong phiên cho thẻ mới/đang học/vừa rớt. \\
    \hline
    Khoảng ngày sau tốt nghiệp & 1 ngày, 6 ngày & Hai mốc đầu của thang ngày. \\
    \hline
    Ngưỡng thành thạo & \( \ge 90 \) ngày & Khoảng ôn đạt mốc này $\rightarrow$ \code{mastered}. \\
    \hline
    Cửa sổ đưa lại trong phiên & 15 phút & Nếu đến hạn trong cửa sổ này, câu hỏi kế được trả kèm ngay. \\
    \hline
  \end{tabular}
  \caption{Tham số của thuật toán lập lịch SM-2 mở rộng}
  \label{tab:srs-params}
\end{table}

% --- PlantUML gợi ý cho Hình máy trạng thái SRS ----------------------
% @startuml
% [*] --> new
% new --> learning : ghi danh (bước 0)
% learning --> learning : sai (q<3) → bước 0\nhoặc đúng nhưng chưa tốt nghiệp
% learning --> review : đúng liên tiếp ≥ 3 lần
% review --> review : đúng → I ← ⌈I×EF⌉
% review --> learning : sai → rớt về bước 0
% review --> mastered : I ≥ 90 ngày
% mastered --> learning : sai → rớt về bước 0
% mastered --> [*]
% @enduml
\figph{srs-state-machine.png}{0.8}{Máy trạng thái vòng đời thẻ học theo SM-2 mở rộng}{fig:srs-states}

Một quyết định thiết kế quan trọng: \code{applySm2} là một \textbf{hàm thuần}
(pure function) --- không phụ thuộc NestJS, cơ sở dữ liệu hay thời gian thực
(mốc \code{now} được truyền vào) --- nên có thể kiểm thử đơn vị triệt để
(mục~\ref{sec:kiem-thu}).

\subsection{Chọn từ học thích nghi theo trình độ}
\label{subsec:impl-picker}

Lớp \code{VocabPickerService} chọn tập từ cho mỗi phiên theo bốn chế độ
(hằng ngày, theo chủ đề, theo bộ thẻ, ôn tập). Nguyên tắc chung là
\textit{ưu tiên từ đã đến hạn} (sắp theo \code{next_review_at} tăng dần),
sau đó bù thêm \textit{từ mới} cho đủ hạn mức phiên. Điểm "thông minh" nằm ở
cách chọn từ mới: thay vì lọc cứng theo trình độ (dễ dẫn tới phiên rỗng khi
người học đã học hết từ ở đúng bậc), hệ thống coi trình độ là một
\textit{ưu tiên mềm} --- sắp các từ chưa học theo \textbf{khoảng cách CEFR}
tới trình độ người học (gần nhất trước), rồi mới tới thứ hạng tần suất:

\begin{verbatim}
// src/learn/vocab-picker.service.ts — ưu tiên từ gần trình độ người học
qb.orderBy(`ABS(${cefrOrdinalCase('v.cefr_level')} - ${userIdx})`,
           'ASC', 'NULLS LAST')
  .addOrderBy('v.frequency_rank', 'ASC', 'NULLS LAST')
  .limit(limit);
\end{verbatim}

Nhờ đó người học đã cạn từ ở bậc của mình (hoặc các từ còn lại không gắn bậc
CEFR) vẫn nhận được các từ phù hợp nhất hiện có thay vì một phiên rỗng. Khi
thực sự không còn từ nào, picker trả về một \textit{lý do rỗng} có ngữ nghĩa
(\code{no_due_cards}, \code{no_more_at_level}, \code{deck_exhausted},
\code{no_enrollment}) để máy khách hiển thị thông điệp phù hợp; chỉ với
\code{no_due_cards} hệ thống mới kèm theo \textit{thời điểm đến hạn kế tiếp}.
Các từ mới được chọn sẽ được \textit{tự động ghi danh} (auto-enroll) vào hàng
đợi học của người dùng trước khi sinh câu hỏi.

\subsection{Sinh câu hỏi thích nghi theo giai đoạn ghi nhớ}
\label{subsec:impl-questions}

Mỗi từ trong phiên được \code{QuestionBuilderService} mở rộng thành một
\textit{bài học} (lesson) gồm một chuỗi câu hỏi xếp từ dễ đến khó. Mức độ khó
được tổ chức thành ba \textbf{dải khó} (difficulty band) và \textbf{12 dạng
câu hỏi} (Bảng~\ref{tab:question-ladder}). Điểm thích nghi cốt lõi:
\textit{giai đoạn ghi nhớ của từ quyết định những dải nào được dùng}. Từ càng
thành thạo thì "sàn" độ khó càng được nâng lên, loại dần các dạng nhận biết
dễ và tập trung vào sản sinh ngôn ngữ:

\begin{table}[h]
  \centering
  \small
  \renewcommand{\arraystretch}{1.3}
  \begin{tabular}{|p{2.6cm}|p{3.4cm}|p{7.6cm}|}
    \hline
    \textbf{Trạng thái từ} & \textbf{Dải khó tối thiểu} & \textbf{Nhóm câu hỏi đại diện} \\
    \hline
    \code{new} & Nhận biết (NEW) & Thẻ học (flashcard), điền khuyết trắc nghiệm, chọn từ/nghĩa, nghe chọn, nhìn hình chọn từ. \\
    \hline
    \code{learning} / \code{review} & Gợi nhớ (REVIEW) & Điền khuyết tự gõ, nghe gõ lại (dictation), luyện phát âm --- người học phải tự tạo ra dạng từ. \\
    \hline
    \code{mastered} & Sản sinh (MASTER) & Phân biệt nghĩa (sense disambiguation) --- dạng khó nhất, đòi hỏi phân biệt giữa các nghĩa của từ đa nghĩa. \\
    \hline
  \end{tabular}
  \caption{Ánh xạ giai đoạn ghi nhớ của từ tới dải khó và nhóm câu hỏi}
  \label{tab:question-ladder}
\end{table}

Trên nền tảng đó, bộ dựng câu hỏi áp dụng thêm các cơ chế giữ cho bài học
\textit{ngắn, đa dạng và khả thi}:

\begin{itemize}
  \item \textbf{Lấy mẫu xác định theo từ:} với mỗi dải, số dạng câu hỏi
        "tuỳ chọn" được giới hạn (mặc định 2 dạng/dải) và được chọn bằng phép
        \textit{xáo trộn xác định} theo định danh của từ. Nhờ vậy mỗi từ luyện
        một tập dạng câu hỏi khác nhau (thay vì mọi bài đều chạy hết các
        dạng), nhưng kết quả vẫn \textit{tái lập được}. Dạng thẻ học luôn được
        giữ làm bước nhập môn.
  \item \textbf{Trần họ điền khuyết:} các dạng cùng "khoét" một từ trong câu
        bị giới hạn số lần (mặc định 2) để một câu ví dụ không bị khoét nhiều
        bước liên tiếp.
  \item \textbf{Lọc theo dữ liệu sẵn có:} mỗi dạng chỉ được dựng khi đủ dữ
        liệu (ví dụ dạng nghe cần có \code{audio_url}, dạng nhìn hình cần ảnh
        nghĩa, dạng dịch cần bản dịch ở ngôn ngữ đích). Dạng nào thiếu dữ liệu
        sẽ bị bỏ qua một cách an toàn.
\end{itemize}

\paragraph{Sinh phương án nhiễu thông minh.} Chất lượng của một câu trắc
nghiệm phụ thuộc vào \textit{phương án nhiễu} (distractor). \code{DistractorService}
không lấy nhiễu ngẫu nhiên mà chọn các từ "dễ nhầm": \textbf{cùng từ loại,
cùng ngôn ngữ, cùng bậc CEFR \( \pm 1 \)}, và \textit{ưu tiên các từ chia sẻ
ít nhất một chủ đề} với từ mục tiêu; chỉ khi không đủ mới nới lỏng điều kiện.
Với dạng phân biệt nghĩa, nhiễu còn được lấy chính từ \textit{các nghĩa khác}
của cùng từ (bẫy đa nghĩa), buộc người học thực sự phân biệt ngữ cảnh.

\subsection{Chấm điểm phi trạng thái có ký HMAC}
\label{subsec:impl-hmac}

Như đã nêu ở mục~\ref{subsec:hmac}, hệ thống chấm điểm phiên học theo cách
\textbf{phi trạng thái}: máy chủ không lưu phiên trong bộ nhớ mà "đóng dấu"
tính toàn vẹn vào chính mỗi câu hỏi. Khi sinh câu hỏi, \code{HmacSignerService}
ký một thông điệp gồm các trường định danh và --- quan trọng --- cả
\code{stepIndex}/\code{stepCount}, để máy khách không thể giả mạo việc câu
nào là câu cuối (câu duy nhất làm thay đổi lịch ôn):

\begin{verbatim}
// src/learn/hmac-signer.service.ts — thông điệp được ký
const message = [userId, vocabularyId, type, exampleId,
                 translationLang ?? '', String(stepIndex),
                 String(stepCount), nonce, String(issuedAtMs)].join('|');
return createHmac('sha256', this.secret).update(message).digest('hex');
\end{verbatim}

Khi người học nộp đáp án (Hình~\ref{fig:seq-answer}), \code{LearnService} thực
hiện ba việc:

\begin{enumerate}
  \item \textbf{Xác minh chữ ký:} kiểm tra chữ ký HMAC bằng phép so sánh
        \textit{thời gian hằng số} (\code{timingSafeEqual}) và kiểm tra hạn
        dùng (mặc định 30 phút); chữ ký sai/hết hạn $\rightarrow$ lỗi 401.
  \item \textbf{Chấm điểm tại chỗ:} \code{AnswerGraderService} \textit{suy
        lại} đáp án đúng từ dữ liệu từ vựng (dạng từ bị khoét trong câu, bản
        dịch của nghĩa, hay chính lemma) rồi so khớp. Điểm chất lượng SM-2
        được quy đổi theo loại câu: trả lời nhanh và đúng được 5, đúng nhưng
        chậm được 4, sai lệch một ký tự (đo bằng khoảng cách Levenshtein) được
        coi là "gần đúng" (3), còn lại 2; riêng thẻ học cho người dùng tự đánh
        giá theo thang bốn mức.
  \item \textbf{Lập lịch ở bước cuối:} mỗi bài học chỉ là \textit{một} sự kiện
        SRS --- chỉ câu hỏi cuối (\code{stepIndex === stepCount - 1}) mới gọi
        \code{submitReview} để chạy \code{applySm2} và cập nhật lịch; các bước
        trước chỉ chấm để phản hồi tức thì. Nhờ vậy một từ không thể "tốt
        nghiệp" chỉ trong một lượt ngồi học.
\end{enumerate}

Việc cập nhật tiến độ và ghi nhật ký hoạt động được thực hiện trong cùng một
giao dịch (mục~\ref{sec:hien-thuc-du-lieu}). Nếu lịch vừa được dời vào trong
\textit{cửa sổ đưa lại} (15 phút), phản hồi sẽ kèm sẵn câu hỏi kế đã ký để máy
khách hiển thị lại từ ngay trong phiên mà không cần gọi lại API tạo phiên.

% =====================================================================
\section{Các tính năng trí tuệ nhân tạo và tự động hoá nội dung}
\label{sec:ai-tu-dong}

Nhóm tính năng thứ hai làm nên tính "thông minh" của hệ thống là việc
\textbf{tự động tạo và làm giàu nội dung} bằng các dịch vụ AI, giải quyết
trực tiếp rào cản "tạo nội dung thủ công tốn công" đã nêu ở
Chương~\ref{chuong:tong-quan}. Tất cả các tác vụ này đều nặng và phụ thuộc
dịch vụ ngoài, nên được hiện thực trên một \textit{hạ tầng xử lý nền} chung.

\subsection{Hạ tầng xử lý nền bất đồng bộ}
\label{subsec:impl-worker}

Mọi tác vụ AI được tổ chức theo mẫu \textbf{nhà sản xuất--người tiêu thụ}
trên \textbf{BullMQ}~\cite{bullmq2024}: tiến trình API tạo bản ghi công việc,
đẩy vào hàng đợi Redis rồi trả về mã \code{202} kèm mã công việc; tiến trình
worker tiêu thụ và xử lý nền (xem Hình~\ref{fig:activity-ai} và
Hình~\ref{fig:activity-practice}). Bảng~\ref{tab:async-jobs} liệt kê các hàng
đợi đã hiện thực. Bốn cơ chế bảo đảm độ tin cậy (NFR-04) được áp dụng nhất
quán cho mọi worker:

\begin{itemize}
  \item \textbf{Tính idempotent:} mỗi worker đọc lại bản ghi và \textit{chỉ}
        xử lý khi trạng thái còn là \code{pending}; một công việc bị chạy lại
        sẽ thoát sớm thay vì tạo dữ liệu trùng.
  \item \textbf{Thử lại có khoảng lùi:} lỗi tạm thời (mạng, 429, một phản hồi
        không phân tích được) khiến BullMQ thử lại với \textit{backoff}; chỉ
        khi cạn số lần thử, sự kiện \code{onFailed} mới ghi trạng thái
        \code{failed}.
  \item \textbf{Tự giới hạn tần suất:} worker đặt \textit{limiter} (mặc định
        15 lượt/phút) và độ song song bằng 1 để không vượt hạn mức của khoá
        AI dùng chung (tầng miễn phí).
  \item \textbf{Hạn mức theo ngày:} các điểm cuối người dùng có hạn mức theo
        ngày (ví dụ 30 lượt chấm câu/người/ngày) để bảo vệ tài nguyên chung,
        trả về lỗi 429 khi vượt.
\end{itemize}

\begin{table}[h]
  \centering
  \small
  \renewcommand{\arraystretch}{1.3}
  \begin{tabular}{|p{3.0cm}|p{4.6cm}|p{6.4cm}|}
    \hline
    \textbf{Hàng đợi} & \textbf{Kích hoạt bởi} & \textbf{Công việc của worker} \\
    \hline
    Làm giàu từ vựng & Tạo nhanh / nhập hàng loạt & Tra từ điển + gọi LLM sinh nghĩa, ví dụ, bản dịch, CEFR; tạo bản ghi từ vựng. \\
    \hline
    Sinh âm thanh & Sau khi tạo/duyệt từ & Lấy bản ghi âm thật hoặc tổng hợp giọng nói, tải lên kho media. \\
    \hline
    Sinh hình ảnh & Quy trình làm giàu nghĩa & Tìm ảnh minh hoạ theo từ, tải lên kho media. \\
    \hline
    Chấm điểm luyện viết & Nộp bài luyện viết/nói & Gọi LLM chấm câu theo rubric, lưu điểm và nhận xét. \\
    \hline
  \end{tabular}
  \caption{Các hàng đợi xử lý nền và vai trò của worker}
  \label{tab:async-jobs}
\end{table}

\paragraph{Lưu ý về mô hình ngôn ngữ.} Các lời gọi LLM được thực hiện qua API
\code{generateContent} của nền tảng \textbf{Google AI Studio}, dùng họ mô
hình mở \textbf{Gemma} của Google DeepMind~\cite{gemma2025technical} làm nền
tảng tham chiếu. Vì mô hình được cấu hình qua biến môi trường
(\code{GEMMA_MODEL}), thiết kế là \textit{độc lập mô hình}; trong cấu hình
triển khai mặc định, do dòng \code{gemma-3-*} đã ngừng phục vụ trên tầng miễn
phí của AI Studio, hệ thống dùng một mô hình Gemini Flash nhẹ tương đương.
Một chi tiết hiện thực quan trọng để tiết kiệm hạn mức: các lời gọi đặt
\code{thinkingConfig.thinkingBudget = 0} nhằm tắt token "suy nghĩ" ẩn ---
nếu không, mô hình suy luận sẽ tiêu hết hạn mức đầu ra cho phần suy nghĩ và
cắt cụt JSON kết quả.

\subsection{Làm giàu dữ liệu từ điển tự động bằng LLM}
\label{subsec:impl-enrich}

Đây là tính năng AI trung tâm: từ một \textit{lemma} đơn lẻ, hệ thống tự
dựng nên một bản ghi từ vựng đầy đủ (use case ở Bảng~\ref{tab:uc-quickcreate}).
\code{EnrichmentProcessor} chọn một trong hai đường đi tuỳ dữ liệu sẵn có:

\begin{itemize}
  \item \textbf{Đường có từ điển (tiếng Anh):} truy vấn Free Dictionary
        API~\cite{freedictionaryapi} để lấy từ loại, định nghĩa, phiên âm
        \textbf{IPA} và từ đồng/trái nghĩa; sau đó LLM chỉ cần bổ sung một
        \textit{gloss} ngắn, \textit{tối thiểu hai câu ví dụ} cho mỗi nghĩa,
        một bản dịch ngắn và ước lượng bậc CEFR. Cách kết hợp này tận dụng dữ
        liệu từ điển chính xác cho phần "khó tin tưởng AI" (định nghĩa, phiên
        âm) và chỉ giao cho AI phần nó làm tốt (ví dụ tự nhiên, bản dịch).
  \item \textbf{Đường chỉ dùng LLM (phi tiếng Anh hoặc từ điển không có):}
        LLM sinh toàn bộ cấu trúc nghĩa theo từng từ loại.
\end{itemize}

Vì Gemma trên AI Studio không hỗ trợ \code{responseSchema} hay vai trò
\textit{hệ thống} (system role), lược đồ JSON mong muốn được \textbf{nhúng
thẳng vào prompt} và phản hồi được \textbf{phân tích phòng thủ}. Bộ phân tích
bóc bỏ dấu rào ```\code{```json}```, dung thứ phần văn xuôi thừa, cắt và giới
hạn độ dài các trường, và --- quan trọng nhất --- \textit{từ chối} một CEFR
sai hoặc thiếu hay một nghĩa có dưới hai ví dụ bằng cách ném lỗi, để worker
\textit{thử lại} thay vì lưu dữ liệu rác. Đoạn dưới minh hoạ lược đồ được yêu
cầu và một phần kiểm tra phòng thủ:

\begin{verbatim}
// Lược đồ JSON nhúng trong prompt (rút gọn)
{ "cefr": "<A1|A2|B1|B2|C1|C2>",
  "senses": [ { "gloss": "<nhãn ngắn>",
                "examples": ["<câu>", "<câu>"],
                "translation": "<bản dịch ngắn>" } ] }

// src/vocabularies/enrichment/gemma-enricher.ts — kiểm tra phòng thủ
if (!VALID_CEFR.has(cefr))
  throw new Error(`enricher returned invalid cefr: ${cefr}`);
if (examples.length < 2)
  throw new Error(`sense ${i + 1} returned fewer than 2 examples`);
\end{verbatim}

\paragraph{Vòng kiểm soát chất lượng.} Kết quả do AI sinh ra luôn được kiểm soát theo
\textit{quyền sở hữu}, đúng tinh thần mục~\ref{subsec:llm-theory}: công việc do
\textit{quản trị viên} tạo sẽ sinh ra \textbf{bản nháp} thuộc kho hệ thống ở
trạng thái \code{is_approved = false} --- không xuất hiện trong phiên học của
ai cho tới khi được duyệt; còn từ do \textit{người học} tạo là từ cá nhân,
riêng tư và tự duyệt. Nhờ \textit{chỉ mục duy nhất bộ phận}
(mục~\ref{sec:hien-thuc-du-lieu}), worker bỏ qua mọi từ loại đã tồn tại nên
không bao giờ ghi đè dữ liệu sẵn có. Các từ được duyệt/tự duyệt tiếp tục kích
hoạt hàng đợi sinh âm thanh.

\subsection{Nhập từ vựng hàng loạt tự động}
\label{subsec:impl-bulk}

Tính năng nhập hàng loạt mở rộng cơ chế làm giàu ở quy mô lớn. Hàm
\code{extractCandidates} (\code{enrichment/import/}) nhận đầu vào ở nhiều định
dạng --- văn bản dán trực tiếp, CSV, bảng tính \textbf{XLSX} hay tệp
\textbf{PDF} --- quy về một khối văn bản rồi đưa qua bộ \textit{tách token}
chung. Bộ tách hỗ trợ hai chế độ: chế độ \textit{danh sách} (mỗi dòng/ô là một
mục) và chế độ \textit{văn xuôi} (tách từ, chuẩn hoá, loại bỏ \textit{từ dừng});
kết quả là danh sách \textit{ứng viên lemma}. Sau khi loại các từ đã có trong
kho, mỗi lemma còn lại được \textbf{phát tán} thành một công việc làm giàu
độc lập (chia sẻ một \code{batch_id}, và có thể kèm \code{target_deck_id} để
tự thêm từ mới vào một bộ thẻ đích sau khi tạo xong). Người dùng theo dõi
tiến độ cả lô qua mã lô thay vì phải chờ đồng bộ.

\subsection{Sinh âm thanh và hình ảnh minh hoạ tự động}
\label{subsec:impl-media}

Hai loại media được sinh tự động qua các pipeline có \textit{phương án dự
phòng} nhiều tầng, kết quả đều được lưu trên kho media \textbf{Cloudinary}~\cite{cloudinary2024}:

\begin{itemize}
  \item \textbf{Âm thanh phát âm:} ưu tiên lấy \textit{bản ghi âm người thật}
        từ Free Dictionary API; nếu không có, tổng hợp bằng \textbf{Microsoft
        Edge TTS} (giọng nơ-ron miễn phí, không cần khoá). Hàm trả về cả nguồn
        đã dùng (\code{dictionary} hay \code{tts}), giúp ưu tiên chất lượng cho
        các từ phổ biến và chỉ dùng TTS cho phần đuôi dài.
  \item \textbf{Hình ảnh minh hoạ nghĩa:} tìm ảnh phù hợp qua kho ảnh
        \textbf{Pexels}~\cite{pexels2024} rồi sao lưu lên Cloudinary. Với các
        từ trừu tượng không có ảnh phù hợp, pipeline trả về \code{null} và để
        trống cho quản trị viên bổ sung --- tự động hoá phần "đuôi dài" của
        các từ cụ thể mà không tạo ra ảnh sai lệch.
\end{itemize}

\subsection{Chấm điểm câu luyện tập bằng LLM}
\label{subsec:impl-judge}

Tính năng luyện viết câu (use case ở Bảng~\ref{tab:uc-practice}) dùng LLM làm
\textit{giám khảo}. Khi người học nộp một câu chứa từ mục tiêu,
\code{PracticeService} kiểm tra hạn mức ngày, tạo bản ghi \code{pending}, đẩy
vào hàng đợi và trả \code{202} ngay; nếu việc đẩy hàng đợi thất bại, bản ghi
được \textit{hoàn tác} để không "kẹt" ở trạng thái chờ và không tiêu phí hạn
mức. Worker \code{ScoringProcessor} gọi hàm \code{scoreSentence}
(\code{gemma-judge.ts}) với một prompt yêu cầu chấm theo \textit{rubric} có
cấu trúc (Bảng~\ref{tab:rubric-fields}); nhiệt độ sinh đặt thấp (0.2) để kết
quả ổn định. Điểm đặc thù: mô hình được yêu cầu ước lượng
\textbf{bậc CEFR mà chính câu đó thể hiện} (không phải năng lực tổng thể của
người học), cùng một câu gợi ý sửa.

\begin{table}[h]
  \centering
  \small
  \renewcommand{\arraystretch}{1.3}
  \begin{tabular}{|p{3.6cm}|p{3.0cm}|p{7.4cm}|}
    \hline
    \textbf{Trường rubric} & \textbf{Kiểu} & \textbf{Ý nghĩa} \\
    \hline
    \code{overall} & số nguyên 0--100 & Điểm tổng của câu. \\
    \hline
    \code{usesTargetWord} & boolean & Câu có dùng đúng từ mục tiêu (kể cả dạng biến cách). \\
    \hline
    \code{correctUsage} & boolean & Từ được dùng đúng một trong các nghĩa dự kiến. \\
    \hline
    \code{criteria} & 4 tiêu chí 0--5 & Ngữ pháp, dùng từ, độ tự nhiên, độ liên quan. \\
    \hline
    \code{cefr} & A1--C2 & Bậc CEFR mà câu thể hiện. \\
    \hline
    \code{feedback} & chuỗi & Một, hai câu nhận xét cho người học. \\
    \hline
    \code{correctedSentence} & chuỗi, tuỳ chọn & Phiên bản câu được cải thiện (nếu cần). \\
    \hline
  \end{tabular}
  \caption{Cấu trúc rubric chấm điểm câu luyện tập do LLM trả về}
  \label{tab:rubric-fields}
\end{table}

Phản hồi của mô hình được phân tích phòng thủ giống như ở phần làm giàu: các
khoảng số bị kẹp về biên hợp lệ, CEFR sai bị từ chối để worker thử lại. Người
học thăm dò kết quả qua mã bài luyện.

\subsection{Chấm điểm phát âm qua vi dịch vụ}
\label{subsec:impl-pron}

Chức năng luyện phát âm (Hình~\ref{fig:seq-pron}) được hiện thực bằng cách
tách phần xử lý tín hiệu/học máy thành một \textbf{vi dịch vụ} riêng viết bằng
Python (\textbf{FastAPI})~\cite{fastapi2024}, còn backend chỉ đóng vai
\textit{proxy mỏng}. \code{PronunciationScoringClient} đóng gói bản ghi âm
của người học thành một yêu cầu \code{multipart/form-data} (trường
\code{audio} và \code{word}), gọi điểm cuối \code{/score} với
\textit{thời gian chờ} ràng buộc, rồi ánh xạ lỗi của dịch vụ ngoài về mã lỗi
HTTP có ngữ nghĩa: \code{422} (âm thanh quá ngắn / không chấm được)
$\rightarrow$ \code{400}, còn không kết nối được / lỗi máy chủ $\rightarrow$
\code{503} (suy giảm có kiểm soát, NFR-04). Kết quả --- điểm tổng và điểm từng
âm vị --- được lưu vào \code{pronunciation_attempts} làm lịch sử luyện tập.
Việc cô lập này giữ cho tiến trình backend chính không phải mang các phụ thuộc
nặng về xử lý âm thanh.

% =====================================================================
\section{Theo dõi tiến độ và tạo động lực tự động}
\label{sec:tien-do}

Phân hệ tiến độ biến nhật ký học thành các chỉ số động lực một cách tự động
(các yêu cầu FR-19, FR-20). Nguồn dữ liệu là bảng \textit{chỉ-thêm}
\code{learning_activity} --- mỗi lượt ôn ghi đúng một dòng, được ghi cùng giao
dịch với cập nhật SRS (mục~\ref{subsec:impl-hmac}) nên các con số luôn nhất
quán với nhau. \code{ProgressService} tính:

\begin{itemize}
  \item \textbf{Thống kê trang chủ:} số từ theo từng trạng thái
        (\code{new}/\code{learning}/\code{review}/\code{mastered}), số thẻ đang
        đến hạn, số lượt ôn hôm nay và thời điểm đến hạn kế tiếp.
  \item \textbf{Chuỗi ngày học (streak):} số ngày liên tiếp (theo lịch UTC) có
        ít nhất một lượt ôn, tính lùi từ ngày ôn gần nhất; chuỗi chỉ còn hiệu
        lực nếu ngày gần nhất là hôm nay hoặc hôm qua.
  \item \textbf{Biểu đồ hoạt động (heatmap):} gộp số lượt ôn theo từng ngày
        \textit{theo múi giờ của người dùng} bằng một truy vấn SQL dùng
        \code{AT TIME ZONE}, để các ô lịch khớp với nửa đêm địa phương của họ.
\end{itemize}

\textbf{Bảng xếp hạng} được \code{LeaderboardService} tính trực tiếp: xếp hạng
người học theo \textit{số từ đã thành thạo} (toàn thời gian), chỉ tính những
tài khoản người học thật, đang hoạt động và \textit{không chọn ẩn}
(\code{leaderboard_opt_out}); thứ hạng của bản thân người gọi được suy ra từ
cùng một tập đã sắp xếp để hai con số không bao giờ mâu thuẫn. Bảng xếp hạng
theo \textit{số từ mới} được dành cho giai đoạn phát triển tiếp theo.

% =====================================================================
\section{Hiện thực hoá tầng giao diện người dùng}
\label{sec:hien-thuc-frontend}

Toàn bộ các API trình bày ở những mục trên được tiêu thụ bởi một ứng dụng giao
diện người dùng \textbf{độc lập}, hiện thực hoá thiết kế giao diện đã phác thảo
ở mục~\ref{sec:thiet-ke-giao-dien} (Chương~\ref{chuong:phan-tich-thiet-ke}).
Ứng dụng được xây dựng bằng \textbf{Next.js~16} theo mô hình \textit{App
Router} với \textbf{React~19} và \textbf{TypeScript} (mục~\ref{sec:cong-nghe-frontend}).
Đặc trưng kiến trúc xuyên suốt là \textit{ưu tiên phía máy chủ}: phần lớn màn
hình là Server Components lấy dữ liệu ngay trên máy chủ, còn mọi thao tác ghi
đi qua \textit{Server Actions}; nhờ đó access token và refresh token không bao
giờ lộ ra mã JavaScript phía trình duyệt. Bảng~\ref{tab:frontend-layout} ánh xạ
các phần chính của ứng dụng tới tệp hiện thực then chốt.

\begin{table}[h]
  \centering
  \small
  \renewcommand{\arraystretch}{1.3}
  \begin{tabular}{|p{3.6cm}|p{4.2cm}|p{6.0cm}|}
    \hline
    \textbf{Phần} & \textbf{Thư mục / tệp} & \textbf{Vai trò} \\
    \hline
    Nhóm tuyến & \code{src/app/(auth)}, \code{(app)}, \code{(admin)} & Phân tách khu vực xác thực, ứng dụng học và quản trị bằng \textit{route group}. \\
    \hline
    Tầng truy cập API & \code{src/lib/api.ts} & Client HTTP chạy phía máy chủ: \code{apiRequest} và \code{authedRequest}. \\
    \hline
    Quản lý phiên & \code{src/lib/auth/} & Lưu/đọc token trong cookie httpOnly, xoay token, đọc hồ sơ người dùng. \\
    \hline
    Kiểm tra dữ liệu & \code{src/lib/validations/} & Lược đồ \code{zod} dùng chung cho biểu mẫu và Server Action. \\
    \hline
    Động cơ phiên học & \code{src/app/(app)/learn/} & Bộ rút gọn \code{session-machine.ts} và các thành phần câu hỏi. \\
    \hline
  \end{tabular}
  \caption{Ánh xạ các phần frontend tới tệp hiện thực then chốt}
  \label{tab:frontend-layout}
\end{table}

\subsection{Cấu trúc App Router và phân nhóm tuyến}
\label{subsec:fe-approuter}

Cây tuyến được tổ chức bằng \textit{nhóm tuyến} (route group --- thư mục đặt
trong ngoặc đơn, không xuất hiện trong URL) để mỗi khu vực có bố cục và lớp bảo
vệ riêng: \code{(auth)} cho đăng nhập/đăng ký/xác minh email, \code{(app)} cho
trải nghiệm học của người dùng (trang chủ, học, khám phá, luyện tập, cộng đồng,
hồ sơ) và \code{(admin)} cho khu vực quản trị từ vựng. Mỗi phân đoạn tận dụng
các tệp đặc biệt của App Router --- \code{layout.tsx} (bố cục dùng chung),
\code{loading.tsx} (trạng thái chờ qua \textit{Suspense}), \code{error.tsx}
(ranh giới lỗi cấp phân đoạn) và \code{not-found.tsx} --- nên một lỗi cục bộ
không làm trắng toàn bộ ứng dụng. Các thành phần chỉ phục vụ một tuyến được đặt
ngay trong thư mục tuyến đó, chỉ nâng lên \code{src/components/} khi dùng lại ở
nhiều nơi.

\subsection{Tầng truy cập API phía máy chủ}
\label{subsec:fe-api}

Mọi lời gọi tới backend đi qua một client HTTP \textit{chỉ chạy phía máy chủ}
trong \code{src/lib/api.ts}. Địa chỉ backend (\code{API_BASE_URL}) được đọc từ
biến môi trường phía máy chủ, cố tình không gắn tiền tố \code{NEXT_PUBLIC_}
để không rò rỉ ra phía máy khách. Module phơi bày hai điểm vào:

\begin{itemize}
  \item \code{apiRequest} --- gọi \textit{không cần xác thực} (đăng nhập, đăng
        ký, làm mới token); tự đặt \code{Content-Type} JSON hoặc để runtime tự
        đặt biên \textit{multipart} khi thân là \code{FormData} (tải tệp lên).
  \item \code{authedRequest} --- đính kèm access token và, khi gặp
        \code{401}, \textit{tự động} làm mới refresh token đúng một lần rồi thử
        lại; chỉ gọi từ Server Action hoặc Route Handler vì nó có thể ghi lại
        cookie token đã xoay.
\end{itemize}

Cả hai chuẩn hoá thân lỗi của NestJS thành một kiểu \code{ApiResult} thống nhất
(cờ \code{ok}, mã trạng thái, dữ liệu, thân lỗi), nhờ đó tầng gọi xử lý lỗi
theo \textit{giá trị trả về} thay vì ngoại lệ --- khớp với triết lý phân biệt
lỗi \textit{dự kiến} và \textit{không dự kiến} của App Router.

\subsection{Quản lý phiên và xác thực phía máy khách}
\label{subsec:fe-session}

Token phiên được lưu trong \textbf{cookie httpOnly} với vòng đời bám đúng hợp
đồng backend (mục~\ref{sec:hien-thuc-auth}): access token 15~phút, refresh
token 30~ngày và được xoay sau mỗi lần gọi \code{/v1/auth/refresh}. Vì các
helper phiên import \code{next/headers}, chúng \textit{chỉ biên dịch được phía
máy chủ} --- một lớp bảo vệ ở mức biên dịch ngăn token lọt vào Client Component.
Khi \code{authedRequest} gặp \code{401}, nó âm thầm đổi refresh token lấy cặp
token mới, ghi lại cookie và thử lại đúng một lần; nếu vẫn thất bại, phiên bị
xoá và người dùng được chuyển về trang đăng nhập. Sau khi xác thực thành công,
Server Action điều hướng người dùng theo vai trò và trạng thái khai báo hồ sơ:
quản trị viên tới khu vực \code{/admin}, người học đã khai báo tới
\code{/dashboard}, người học chưa khai báo tới luồng \code{/onboarding}.

\subsection{Động cơ phiên học phía máy khách}
\label{subsec:fe-learn}

Phần tương tác nhất của ứng dụng là vòng lặp phiên học. Hàng đợi câu hỏi được
điều phối bằng một \textbf{bộ rút gọn (reducer) thuần} trong
\code{session-machine.ts}, tách hẳn khỏi React và mạng nên kiểm thử đơn vị được
độc lập. Nó nắm danh sách câu hỏi còn lại (câu đang hiển thị luôn là
\code{queue[0]}), các bộ đếm đúng/đã trả lời và trạng thái kết thúc
(\code{active}, \code{empty}, \code{done}, \code{expired}, \code{error}). Khi
một thẻ cần lặp lại trong cùng phiên, nó được \textit{nối lại vào cuối} hàng
đợi thay vì hẹn giờ theo đồng hồ thực. Bộ rút gọn này được ghép với hai Server
Action --- \code{startSessionAction} và \code{submitAnswerAction} --- gọi tới
\code{/v1/me/learn/session} và \code{/v1/me/learn/answer}. Vì máy chủ \textit{ký
HMAC} từng câu hỏi (mục~\ref{subsec:impl-hmac}), một \code{401} còn lại sau khi
đã làm mới token được hiểu là chữ ký câu hỏi đã hết hạn hoặc bị giả mạo, và giao
diện chuyển sang trạng thái \code{expired} mời người học bắt đầu phiên mới. Mỗi
\textit{chế độ} câu hỏi (thẻ ghi nhớ, điền khuyết trắc nghiệm/gõ chữ, nghe chép
chính tả, chọn theo hình/âm thanh, luyện phát âm\ldots) là một thành phần riêng
dưới \code{learn/questions/}, nhận dữ liệu đã ký từ máy chủ và chỉ chịu trách
nhiệm thu thập câu trả lời để nộp lại.

% =====================================================================
\section{Kiểm thử hệ thống}
\label{sec:kiem-thu}

Việc kiểm thử được hỗ trợ trực tiếp bởi một lựa chọn thiết kế lặp lại xuyên
suốt chương: \textbf{tách logic thông minh thành các hàm thuần, độc lập
framework}. Các hàm như \code{applySm2} (SM-2), bộ ký HMAC, bộ chấm đáp án,
bộ dựng prompt và bộ phân tích phản hồi LLM đều không phụ thuộc NestJS hay cơ
sở dữ liệu, nên có thể kiểm thử đơn vị nhanh và tất định. Chiến lược kiểm thử
gồm hai mức (Bảng~\ref{tab:test-targets}):

\begin{itemize}
  \item \textbf{Kiểm thử đơn vị} bằng \textbf{Jest}~\cite{jest2024}: tập trung
        vào các thuật toán và bộ phân tích --- nơi lỗi tinh vi dễ lọt và hậu
        quả lớn (lập lịch sai, lưu dữ liệu AI rác, chấm điểm sai).
  \item \textbf{Kiểm thử tích hợp (end-to-end)} bằng \textbf{Supertest}: gọi
        thật vào các điểm cuối API để kiểm tra luồng xác thực, kiểm tra đầu
        vào và mã trạng thái trả về.
\end{itemize}

\begin{table}[h]
  \centering
  \small
  \renewcommand{\arraystretch}{1.3}
  \begin{tabular}{|p{5.0cm}|p{9.4cm}|}
    \hline
    \textbf{Thành phần} & \textbf{Trọng tâm kiểm thử} \\
    \hline
    \code{srs.ts} & Chuyển trạng thái SM-2, bước học, tốt nghiệp, rớt bước, ngưỡng thành thạo. \\
    \hline
    \code{hmac-signer.service.ts} & Ký/xác minh, chống giả mạo, hết hạn, so sánh thời gian hằng số. \\
    \hline
    \code{answer-grader.service.ts} & Quy đổi điểm chất lượng cho từng loại câu hỏi. \\
    \hline
    \code{question-builder.service.ts} & Dựng bài học đúng dải khó, tôn trọng các trần và lọc dữ liệu. \\
    \hline
    \code{gemma-judge.ts}, \code{gemma-enricher.ts} & Phân tích phòng thủ JSON, kẹp số, từ chối CEFR sai. \\
    \hline
    \code{import/tokenize.ts}, \code{normalize.ts}, \code{xlsx-parser.ts} & Tách lemma, chuẩn hoá, đọc bảng tính cho nhập hàng loạt. \\
    \hline
  \end{tabular}
  \caption{Một số thành phần tiêu biểu và trọng tâm kiểm thử}
  \label{tab:test-targets}
\end{table}

Chất lượng mã nguồn còn được kiểm soát ở mức tĩnh bằng trình kiểm tra kiểu của
TypeScript (chế độ \textit{strict}), \textbf{ESLint} và \textbf{Prettier},
chạy trước mỗi lần commit (đáp ứng NFR-06).

% =====================================================================
\section{Kết chương}
\label{sec:ket-chuong-4}

Chương này đã trình bày quá trình hiện thực hoá hệ thống, với trọng tâm là các
thành phần thông minh và tự động hoá. Ở phía \textit{động cơ học}, hệ thống
thích nghi với từng người học qua thuật toán SM-2 mở rộng bước học
(\code{srs.ts}), cơ chế chọn từ theo khoảng cách CEFR, bộ sinh câu hỏi nhiều
tầng theo giai đoạn ghi nhớ với phương án nhiễu "dễ nhầm", và cơ chế chấm điểm
phi trạng thái an toàn bằng chữ ký HMAC. Ở phía \textit{nội dung}, các tính
năng AI --- làm giàu từ điển bằng LLM kết hợp từ điển, nhập hàng loạt từ nhiều
định dạng, sinh âm thanh và hình ảnh tự động, chấm điểm câu và phát âm --- đều
chạy trên một hạ tầng xử lý nền bất đồng bộ có tính idempotent, biết thử lại,
tự giới hạn tần suất và suy giảm có kiểm soát. Ở phía \textit{máy khách}, một
ứng dụng Next.js theo App Router tiêu thụ toàn bộ các API trên qua một tầng
truy cập phía máy chủ, quản lý phiên bằng cookie httpOnly với cơ chế xoay token
tự động, và điều phối phiên học bằng một bộ rút gọn thuần tách khỏi React. Tất
cả được củng cố bằng các chỉ số tiến độ tự động và một chiến lược kiểm thử dựa
trên việc tách logic thành các hàm thuần. Những kết quả hiện thực hoá này là cơ
sở để chương tiếp theo đánh giá hệ thống và rút ra kết luận.

%  vào main.tex (sau Chương 3).
%
%  GHI CHÚ VỀ HÌNH VẼ: dùng macro \figph{<tên-file>}{<tỉ-lệ-rộng>}{<chú-thích>}{<nhãn>}
%  (đã khai báo ở các chương trước; khai báo lại bằng \providecommand để an
%  toàn nếu biên dịch độc lập). Mã nguồn (code) trình bày bằng môi trường
%  verbatim như các chương trước.
%
%  KHÔNG hardcode số chương/mục/bảng/hình/công thức — luôn dùng \label + \ref / \cite.
% =====================================================================

% --- Macro tiện ích (idempotent — không định nghĩa lại nếu đã có) ------
\providecommand{\code}[1]{\texttt{\detokenize{#1}}} % in định danh snake_case an toàn
\providecommand{\figph}[4]{%
  \begin{figure}[htbp]
    \centering
    \IfFileExists{images/#1}%
      {\includegraphics[width=#2\textwidth,height=0.82\textheight,keepaspectratio]{#1}}%
      {\setlength{\fboxsep}{10pt}\fbox{\begin{minipage}[c][3.6cm][c]{#2\textwidth}\centering\itshape Sơ đồ minh hoạ --- chèn tệp ảnh \code{images/#1} vào thư mục \code{images/}.\end{minipage}}}%
    \caption{#3}%
    \label{#4}%
  \end{figure}}

\chapter{HIỆN THỰC HOÁ HỆ THỐNG VÀ CÁC TÍNH NĂNG THÔNG MINH}
\label{chuong:hien-thuc-hoa}

Hai chương trước đã trình bày kết quả phân tích, thiết kế
(Chương~\ref{chuong:phan-tich-thiet-ke}) và cơ sở lý thuyết, công nghệ
(Chương~\ref{chuong:co-so-ly-thuyet}). Chương này trình bày cách các thiết
kế đó được \textit{hiện thực hoá} thành mã nguồn chạy được, với trọng tâm là
các \textbf{thành phần thông minh và tự động hoá} làm nên giá trị cốt lõi của
hệ thống: động cơ học theo ôn tập ngắt quãng, cơ chế sinh câu hỏi và chấm
điểm thích nghi, các tính năng làm giàu nội dung và chấm điểm dựa trên trí
tuệ nhân tạo, cùng hạ tầng xử lý nền bất đồng bộ liên kết chúng lại.

Mục~\ref{sec:tong-quan-hien-thuc} mô tả tổ chức mã nguồn và các mối quan tâm
xuyên suốt. Mục~\ref{sec:hien-thuc-du-lieu} và
mục~\ref{sec:hien-thuc-auth} trình bày ngắn gọn hiện thực hoá tầng dữ liệu và
phân hệ xác thực. Trọng tâm của chương nằm ở
mục~\ref{sec:dong-co-hoc} (động cơ học thông minh) và
mục~\ref{sec:ai-tu-dong} (các tính năng AI và tự động hoá nội dung). Cuối
cùng, mục~\ref{sec:tien-do} trình bày phân hệ theo dõi tiến độ tự động,
mục~\ref{sec:hien-thuc-frontend} trình bày hiện thực hoá tầng giao diện người
dùng (frontend), mục~\ref{sec:kiem-thu} trình bày chiến lược kiểm thử và
mục~\ref{sec:ket-chuong-4} tóm tắt chương.

% =====================================================================
\section{Tổng quan quá trình hiện thực hoá}
\label{sec:tong-quan-hien-thuc}

Hệ thống được hiện thực hoá đúng theo ngăn xếp công nghệ đã chọn ở
Chương~\ref{chuong:co-so-ly-thuyet}: backend viết bằng \textbf{NestJS~11}
trên \textbf{TypeScript} (chế độ \textit{strict}), truy cập
\textbf{PostgreSQL} qua \textbf{TypeORM}, và dùng \textbf{Redis/BullMQ} cho
xử lý nền. Toàn bộ mã nguồn được biên dịch bằng trình biên dịch của NestJS và
khởi chạy dưới \textbf{hai tiến trình} chia sẻ chung cơ sở dữ liệu và hàng
đợi, đúng mô hình nhà sản xuất--người tiêu thụ đã nêu ở
mục~\ref{sec:kien-truc} của Chương~\ref{chuong:phan-tich-thiet-ke}:

\begin{itemize}
  \item \textbf{Tiến trình API} (\code{npm run start}): phục vụ các yêu cầu
        HTTP đồng bộ; với tác vụ nặng nó chỉ tạo bản ghi công việc rồi đẩy
        vào hàng đợi và trả về ngay.
  \item \textbf{Tiến trình Worker} (\code{npm run start:worker}): tiêu thụ
        hàng đợi để thực thi các tác vụ nền (làm giàu từ vựng, sinh âm thanh,
        chấm điểm luyện tập).
\end{itemize}

Mọi tham số vận hành (chuỗi kết nối CSDL, khoá bí mật, khoá API của dịch vụ
AI, hạn mức tần suất\ldots) đều được nạp từ \textbf{biến môi trường} qua
\code{@nestjs/config}, đáp ứng yêu cầu tính khả chuyển NFR-08 trong
Bảng~\ref{tab:nfr}: cùng một ảnh mã nguồn chạy được ở nhiều môi trường mà
không cần sửa code.

\subsection{Tổ chức mã nguồn theo mô-đun}
\label{subsec:to-chuc-ma-nguon}

Mỗi mô-đun nghiệp vụ ở Bảng~\ref{tab:modules} (Chương~\ref{chuong:phan-tich-thiet-ke})
được hiện thực thành một thư mục dưới \code{src/} tự đóng gói controller,
service, entity và DTO của mình. Bảng~\ref{tab:source-layout} ánh xạ các
mô-đun trọng yếu tới những tệp hiện thực then chốt --- là cơ sở để các mục
sau tham chiếu trực tiếp.

\begin{table}[h]
  \centering
  \small
  \renewcommand{\arraystretch}{1.3}
  \begin{tabular}{|p{2.6cm}|p{4.4cm}|p{6.6cm}|}
    \hline
    \textbf{Mô-đun} & \textbf{Thư mục mã nguồn} & \textbf{Tệp hiện thực then chốt} \\
    \hline
    \code{auth} & \code{src/auth} & \code{token.service.ts} (cấp/xoay token), \code{jwt.strategy.ts}, các \code{guards/}, \code{services/} (Google/Apple/GitHub). \\
    \hline
    \code{learn} & \code{src/learn} & \code{learn.service.ts}, \code{vocab-picker.service.ts}, \code{question-builder.service.ts}, \code{distractor.service.ts}, \code{answer-grader.service.ts}, \code{hmac-signer.service.ts}, \code{question-bands.ts}. \\
    \hline
    \code{progress} & \code{src/progress} & \code{srs.ts} (thuật toán SM-2), \code{progress.service.ts} (lập lịch, thống kê, chuỗi ngày học). \\
    \hline
    \code{vocabularies} & \code{src/vocabularies} & \code{enrichment/} (làm giàu bằng LLM, nhập hàng loạt), \code{audio/}, \code{images/}, \code{vocabulary-persistence.service.ts}. \\
    \hline
    \code{practice} & \code{src/practice} & \code{gemma-judge.ts} (chấm câu bằng LLM), \code{scoring.processor.ts}, \code{practice.service.ts}. \\
    \hline
    \code{pronunciation} & \code{src/pronunciation} & \code{pronunciation-scoring.client.ts} (proxy tới vi dịch vụ), \code{pronunciation.service.ts}. \\
    \hline
    \code{leaderboard} & \code{src/leaderboard} & \code{leaderboard.service.ts}. \\
    \hline
    \code{config} & \code{src/config} & Các tệp cấu hình theo miền: \code{learn.config.ts}, \code{gemma.config.ts}, \code{auth.config.ts}\ldots \\
    \hline
  \end{tabular}
  \caption{Ánh xạ mô-đun nghiệp vụ tới tệp hiện thực then chốt}
  \label{tab:source-layout}
\end{table}

\subsection{Các mối quan tâm xuyên suốt}
\label{subsec:cross-cutting}

Nhờ cơ chế \textit{tiêm phụ thuộc} và các thành phần khai báo của NestJS, các
mối quan tâm chung được hiện thực một lần và áp dụng toàn cục thay vì lặp lại
trong từng controller:

\begin{itemize}
  \item \textbf{Kiểm tra đầu vào:} một \code{ValidationPipe} toàn cục với
        \code{whitelist}, \code{forbidNonWhitelisted} và \code{transform}
        biến mỗi DTO (khai báo bằng \code{class-validator}) thành lớp bảo vệ:
        trường lạ bị loại bỏ, dữ liệu sai kiểu bị chặn bằng lỗi 400 trước khi
        chạm tới service.
  \item \textbf{Xác thực và phân quyền:} \code{JwtAuthGuard} (dựa trên
        \code{passport-jwt}, xem \code{jwt.strategy.ts}) trích và xác minh
        access token ở mọi endpoint cần đăng nhập; \code{RolesGuard} kết hợp
        decorator \code{@Roles('admin')} chặn các thao tác quản trị; quyền sở
        hữu tài nguyên được kiểm tra ngay trong service (mục~\ref{subsec:rbac}).
  \item \textbf{Phiên bản hoá API:} bật \textit{URI versioning} mặc định
        \code{v1}; controller khai báo \code{@Controller(\{ path, version: '1' \})}
        nên mọi tài nguyên nằm dưới tiền tố \code{/v1} (mục~\ref{subsec:rest-theory}).
  \item \textbf{Cấu hình theo miền:} mỗi nhóm tham số được gói trong một
        \textit{namespace} cấu hình riêng (ví dụ \code{learn.config.ts},
        \code{gemma.config.ts}), tách bạch và có giá trị mặc định an toàn cho
        môi trường phát triển.
\end{itemize}

% =====================================================================
\section{Hiện thực hoá tầng dữ liệu}
\label{sec:hien-thuc-du-lieu}

Tầng dữ liệu hiện thực 17 bảng trong Bảng~\ref{tab:db-overview} thành các lớp
\textit{thực thể} (entity) TypeORM. Mỗi thực thể khai báo khoá chính
\code{uuid}, tên cột \code{snake_case} và kiểu \code{timestamptz} cho các mốc
thời gian, đúng quy ước thiết kế ở Chương~\ref{chuong:phan-tich-thiet-ke}.
Một số kỹ thuật hiện thực đáng chú ý:

\begin{itemize}
  \item \textbf{Giao dịch bảo đảm tính nguyên tử:} các thao tác ghi nhiều
        bảng được gói trong \code{dataSource.transaction(...)}. Ví dụ tiêu
        biểu là việc nộp một lượt ôn (mục~\ref{subsec:impl-hmac}): bản ghi
        tiến độ và bản ghi nhật ký hoạt động được lưu trong cùng một giao
        dịch, nên một lượt ôn không thể "ghi nửa vời" --- đáp ứng NFR-04 trong
        Bảng~\ref{tab:nfr}.
  \item \textbf{Ràng buộc toàn vẹn và chống trùng lặp:} ngoài khoá ngoại và
        quy tắc xoá lan truyền, hệ thống dùng \textit{chỉ mục duy nhất bộ
        phận} (partial unique index) trên bộ ba \code{(language, lemma,
        part_of_speech)} áp dụng riêng cho từ thuộc kho hệ thống
        (\code{WHERE source = 'system'}); nhờ đó kho hệ thống không bao giờ có
        hai bản ghi trùng từ loại cho cùng một từ, trong khi từ cá nhân của
        người dùng vẫn không bị ràng buộc này.
  \item \textbf{Quản lý lược đồ bằng migration:} mọi thay đổi lược đồ được mô
        tả bằng các \textit{migration} có thể phiên bản hoá, được nhận diện và
        chạy qua công cụ dòng lệnh của TypeORM. Cách này cho phép tái lập lược
        đồ một cách xác định giữa các môi trường (đáp ứng NFR-06).
  \item \textbf{Phi chuẩn hoá có kiểm soát:} một số trường được phi chuẩn hoá
        để tăng tốc đọc, ví dụ \code{decks.vocab_count} (số từ trong bộ thẻ)
        được cập nhật khi thêm/bớt thành viên thay vì đếm mỗi lần truy vấn.
\end{itemize}

% =====================================================================
\section{Hiện thực hoá xác thực và phân quyền}
\label{sec:hien-thuc-auth}

Phân hệ \code{auth} hiện thực các kỹ thuật bảo mật đã trình bày về lý thuyết
ở mục~\ref{sec:xac-thuc-phan-quyen}. Lớp \code{TokenService} đảm nhiệm vòng
đời phiên đăng nhập với một số quyết định hiện thực đáng chú ý:

\begin{itemize}
  \item \textbf{Cặp token bất đối xứng về bản chất:} \textit{access token} là
        một \textbf{JWT} được ký, tự chứa và sống ngắn (mục~\ref{subsec:jwt}),
        trong khi \textit{refresh token} là một chuỗi ngẫu nhiên 48 byte
        \textit{mờ} (opaque) --- không mang thông tin, chỉ lưu \textbf{bản băm
        SHA-256} trong bảng \code{refresh_tokens}. Máy chủ không bao giờ lưu
        token gốc, nên dù cơ sở dữ liệu bị lộ thì token cũng không thể tái sử
        dụng.
  \item \textbf{Xoay vòng và phát hiện tái dùng:} mỗi lần làm mới phiên, token
        cũ bị đánh dấu thu hồi và một token mới được cấp. Nếu một token đã thu
        hồi (hoặc không tồn tại) được trình lại --- dấu hiệu rò rỉ --- hệ thống
        \textbf{thu hồi toàn bộ} token còn hiệu lực của người dùng đó và trả
        về lỗi 401, hạn chế thiệt hại khi token bị đánh cắp.
  \item \textbf{Băm mật khẩu và đăng nhập bên thứ ba:} mật khẩu được băm bằng
        \textbf{bcrypt} (mục~\ref{subsec:bcrypt}); với OAuth/OIDC
        (mục~\ref{subsec:oauth}), token danh tính do nhà cung cấp cấp được xác
        minh rồi liên kết tới một bản ghi trong \code{user_identities}, cho
        phép một tài khoản đăng nhập bằng nhiều phương thức.
  \item \textbf{Chống vét cạn:} endpoint đăng nhập được giới hạn tần suất bằng
        \code{@nestjs/throttler}, trả về lỗi 429 khi vượt ngưỡng.
\end{itemize}

% =====================================================================
\section{Động cơ học thông minh theo ôn tập ngắt quãng}
\label{sec:dong-co-hoc}

Đây là phân hệ \textit{thông minh} cốt lõi của hệ thống. Một phiên học không
phải là một danh sách câu hỏi tĩnh: với mỗi người học và mỗi từ, hệ thống
\textbf{thích nghi} ở ba cấp độ --- chọn \textit{từ nào} để học, sinh
\textit{dạng câu hỏi nào} phù hợp giai đoạn ghi nhớ, và quyết định
\textit{khi nào} ôn lại. Toàn bộ quy trình tuân theo luồng đã mô hình hoá ở
Hình~\ref{fig:activity-srs}: chọn từ $\rightarrow$ mở rộng thành chuỗi câu hỏi
đã ký $\rightarrow$ chấm điểm $\rightarrow$ (ở bước cuối) lập lịch lại.

\subsection{Lập lịch ôn tập SM-2 mở rộng bước học}
\label{subsec:impl-sm2}

Hạt nhân lập lịch là hàm thuần \code{applySm2} trong \code{src/progress/srs.ts}:
nó nhận trạng thái hiện tại của một cặp (người dùng, từ), một
\textit{điểm chất lượng} \( q \in \{0,\dots,5\} \) và danh sách bước học, rồi
trả về trạng thái mới cùng thời điểm ôn kế tiếp. Hàm hiện thực thuật toán
\textbf{SM-2} (mục~\ref{subsec:sm2-theory}) được \textit{mở rộng bằng các bước
học} theo phong cách Anki. Hệ số dễ nhớ \( EF \) được cập nhật ở \textit{mọi}
lượt ôn theo công thức~\eqref{eq:sm2-ef}, còn khoảng cách ngày tuân theo
công thức~\eqref{eq:sm2-interval}; phần mở rộng bổ sung một thang
\textit{bước học} tính bằng phút để củng cố từ mới ngay trong phiên trước khi
giao cho thang ngày:

\begin{verbatim}
// src/progress/srs.ts — cập nhật hệ số dễ nhớ (chặn dưới 1.3)
function updateEaseFactor(easeFactor, quality) {
  return Math.max(
    1.3,
    easeFactor + (0.1 - (5 - quality) * (0.08 + (5 - quality) * 0.02)),
  );
}
\end{verbatim}

Logic chuyển trạng thái được tóm tắt như sau (Hình~\ref{fig:srs-states}):

\begin{itemize}
  \item Thẻ đang ở \textit{trạng thái bước} (\code{learningStepIndex} khác
        \code{null}) dùng khoảng cách \textbf{phút} lấy từ danh sách bước (mặc
        định \([1, 10]\) phút). Trả lời đúng (\( q \ge 3 \)) đẩy lên một bước;
        qua bước cuối, thẻ \textit{tốt nghiệp} sang thang ngày. Trả lời sai
        (\( q < 3 \)) đặt lại về bước 0.
  \item Thẻ đã tốt nghiệp dùng thang ngày: \( I = 1 \) ngày (lần đầu), \( 6 \)
        ngày (lần hai), sau đó \( I \leftarrow \lceil I \times EF \rceil \).
        Một lần trả lời sai sẽ đưa thẻ \textit{rớt} về bước 0 (ôn lại trong
        phút) thay vì chờ tới hôm sau.
  \item Trạng thái hiển thị được suy ra từ số lần lặp và khoảng cách:
        \( I \ge 90 \) ngày $\rightarrow$ \code{mastered}; số lần ôn đúng liên
        tiếp \( \ge 3 \) $\rightarrow$ \code{review}; còn lại \code{learning}.
\end{itemize}

Bảng~\ref{tab:srs-params} liệt kê các tham số then chốt và giá trị mặc định
của chúng (đều có thể cấu hình qua biến môi trường trong \code{learn.config.ts}).

\begin{table}[h]
  \centering
  \small
  \renewcommand{\arraystretch}{1.3}
  \begin{tabular}{|p{5.2cm}|p{3.0cm}|p{6.0cm}|}
    \hline
    \textbf{Tham số} & \textbf{Giá trị mặc định} & \textbf{Ý nghĩa} \\
    \hline
    Hệ số dễ nhớ khởi tạo \( EF_0 \) & 2.5 & Giá trị ban đầu của \( EF \). \\
    \hline
    Chặn dưới \( EF_{\min} \) & 1.3 & Ngăn khoảng ôn co lại quá nhanh. \\
    \hline
    Bước học (phút) & \([1, 10]\) & Khoảng cách trong phiên cho thẻ mới/đang học/vừa rớt. \\
    \hline
    Khoảng ngày sau tốt nghiệp & 1 ngày, 6 ngày & Hai mốc đầu của thang ngày. \\
    \hline
    Ngưỡng thành thạo & \( \ge 90 \) ngày & Khoảng ôn đạt mốc này $\rightarrow$ \code{mastered}. \\
    \hline
    Cửa sổ đưa lại trong phiên & 15 phút & Nếu đến hạn trong cửa sổ này, câu hỏi kế được trả kèm ngay. \\
    \hline
  \end{tabular}
  \caption{Tham số của thuật toán lập lịch SM-2 mở rộng}
  \label{tab:srs-params}
\end{table}

% --- PlantUML gợi ý cho Hình máy trạng thái SRS ----------------------
% @startuml
% [*] --> new
% new --> learning : ghi danh (bước 0)
% learning --> learning : sai (q<3) → bước 0\nhoặc đúng nhưng chưa tốt nghiệp
% learning --> review : đúng liên tiếp ≥ 3 lần
% review --> review : đúng → I ← ⌈I×EF⌉
% review --> learning : sai → rớt về bước 0
% review --> mastered : I ≥ 90 ngày
% mastered --> learning : sai → rớt về bước 0
% mastered --> [*]
% @enduml
\figph{srs-state-machine.png}{0.8}{Máy trạng thái vòng đời thẻ học theo SM-2 mở rộng}{fig:srs-states}

Một quyết định thiết kế quan trọng: \code{applySm2} là một \textbf{hàm thuần}
(pure function) --- không phụ thuộc NestJS, cơ sở dữ liệu hay thời gian thực
(mốc \code{now} được truyền vào) --- nên có thể kiểm thử đơn vị triệt để
(mục~\ref{sec:kiem-thu}).

\subsection{Chọn từ học thích nghi theo trình độ}
\label{subsec:impl-picker}

Lớp \code{VocabPickerService} chọn tập từ cho mỗi phiên theo bốn chế độ
(hằng ngày, theo chủ đề, theo bộ thẻ, ôn tập). Nguyên tắc chung là
\textit{ưu tiên từ đã đến hạn} (sắp theo \code{next_review_at} tăng dần),
sau đó bù thêm \textit{từ mới} cho đủ hạn mức phiên. Điểm "thông minh" nằm ở
cách chọn từ mới: thay vì lọc cứng theo trình độ (dễ dẫn tới phiên rỗng khi
người học đã học hết từ ở đúng bậc), hệ thống coi trình độ là một
\textit{ưu tiên mềm} --- sắp các từ chưa học theo \textbf{khoảng cách CEFR}
tới trình độ người học (gần nhất trước), rồi mới tới thứ hạng tần suất:

\begin{verbatim}
// src/learn/vocab-picker.service.ts — ưu tiên từ gần trình độ người học
qb.orderBy(`ABS(${cefrOrdinalCase('v.cefr_level')} - ${userIdx})`,
           'ASC', 'NULLS LAST')
  .addOrderBy('v.frequency_rank', 'ASC', 'NULLS LAST')
  .limit(limit);
\end{verbatim}

Nhờ đó người học đã cạn từ ở bậc của mình (hoặc các từ còn lại không gắn bậc
CEFR) vẫn nhận được các từ phù hợp nhất hiện có thay vì một phiên rỗng. Khi
thực sự không còn từ nào, picker trả về một \textit{lý do rỗng} có ngữ nghĩa
(\code{no_due_cards}, \code{no_more_at_level}, \code{deck_exhausted},
\code{no_enrollment}) để máy khách hiển thị thông điệp phù hợp; chỉ với
\code{no_due_cards} hệ thống mới kèm theo \textit{thời điểm đến hạn kế tiếp}.
Các từ mới được chọn sẽ được \textit{tự động ghi danh} (auto-enroll) vào hàng
đợi học của người dùng trước khi sinh câu hỏi.

\subsection{Sinh câu hỏi thích nghi theo giai đoạn ghi nhớ}
\label{subsec:impl-questions}

Mỗi từ trong phiên được \code{QuestionBuilderService} mở rộng thành một
\textit{bài học} (lesson) gồm một chuỗi câu hỏi xếp từ dễ đến khó. Mức độ khó
được tổ chức thành ba \textbf{dải khó} (difficulty band) và \textbf{12 dạng
câu hỏi} (Bảng~\ref{tab:question-ladder}). Điểm thích nghi cốt lõi:
\textit{giai đoạn ghi nhớ của từ quyết định những dải nào được dùng}. Từ càng
thành thạo thì "sàn" độ khó càng được nâng lên, loại dần các dạng nhận biết
dễ và tập trung vào sản sinh ngôn ngữ:

\begin{table}[h]
  \centering
  \small
  \renewcommand{\arraystretch}{1.3}
  \begin{tabular}{|p{2.6cm}|p{3.4cm}|p{7.6cm}|}
    \hline
    \textbf{Trạng thái từ} & \textbf{Dải khó tối thiểu} & \textbf{Nhóm câu hỏi đại diện} \\
    \hline
    \code{new} & Nhận biết (NEW) & Thẻ học (flashcard), điền khuyết trắc nghiệm, chọn từ/nghĩa, nghe chọn, nhìn hình chọn từ. \\
    \hline
    \code{learning} / \code{review} & Gợi nhớ (REVIEW) & Điền khuyết tự gõ, nghe gõ lại (dictation), luyện phát âm --- người học phải tự tạo ra dạng từ. \\
    \hline
    \code{mastered} & Sản sinh (MASTER) & Phân biệt nghĩa (sense disambiguation) --- dạng khó nhất, đòi hỏi phân biệt giữa các nghĩa của từ đa nghĩa. \\
    \hline
  \end{tabular}
  \caption{Ánh xạ giai đoạn ghi nhớ của từ tới dải khó và nhóm câu hỏi}
  \label{tab:question-ladder}
\end{table}

Trên nền tảng đó, bộ dựng câu hỏi áp dụng thêm các cơ chế giữ cho bài học
\textit{ngắn, đa dạng và khả thi}:

\begin{itemize}
  \item \textbf{Lấy mẫu xác định theo từ:} với mỗi dải, số dạng câu hỏi
        "tuỳ chọn" được giới hạn (mặc định 2 dạng/dải) và được chọn bằng phép
        \textit{xáo trộn xác định} theo định danh của từ. Nhờ vậy mỗi từ luyện
        một tập dạng câu hỏi khác nhau (thay vì mọi bài đều chạy hết các
        dạng), nhưng kết quả vẫn \textit{tái lập được}. Dạng thẻ học luôn được
        giữ làm bước nhập môn.
  \item \textbf{Trần họ điền khuyết:} các dạng cùng "khoét" một từ trong câu
        bị giới hạn số lần (mặc định 2) để một câu ví dụ không bị khoét nhiều
        bước liên tiếp.
  \item \textbf{Lọc theo dữ liệu sẵn có:} mỗi dạng chỉ được dựng khi đủ dữ
        liệu (ví dụ dạng nghe cần có \code{audio_url}, dạng nhìn hình cần ảnh
        nghĩa, dạng dịch cần bản dịch ở ngôn ngữ đích). Dạng nào thiếu dữ liệu
        sẽ bị bỏ qua một cách an toàn.
\end{itemize}

\paragraph{Sinh phương án nhiễu thông minh.} Chất lượng của một câu trắc
nghiệm phụ thuộc vào \textit{phương án nhiễu} (distractor). \code{DistractorService}
không lấy nhiễu ngẫu nhiên mà chọn các từ "dễ nhầm": \textbf{cùng từ loại,
cùng ngôn ngữ, cùng bậc CEFR \( \pm 1 \)}, và \textit{ưu tiên các từ chia sẻ
ít nhất một chủ đề} với từ mục tiêu; chỉ khi không đủ mới nới lỏng điều kiện.
Với dạng phân biệt nghĩa, nhiễu còn được lấy chính từ \textit{các nghĩa khác}
của cùng từ (bẫy đa nghĩa), buộc người học thực sự phân biệt ngữ cảnh.

\subsection{Chấm điểm phi trạng thái có ký HMAC}
\label{subsec:impl-hmac}

Như đã nêu ở mục~\ref{subsec:hmac}, hệ thống chấm điểm phiên học theo cách
\textbf{phi trạng thái}: máy chủ không lưu phiên trong bộ nhớ mà "đóng dấu"
tính toàn vẹn vào chính mỗi câu hỏi. Khi sinh câu hỏi, \code{HmacSignerService}
ký một thông điệp gồm các trường định danh và --- quan trọng --- cả
\code{stepIndex}/\code{stepCount}, để máy khách không thể giả mạo việc câu
nào là câu cuối (câu duy nhất làm thay đổi lịch ôn):

\begin{verbatim}
// src/learn/hmac-signer.service.ts — thông điệp được ký
const message = [userId, vocabularyId, type, exampleId,
                 translationLang ?? '', String(stepIndex),
                 String(stepCount), nonce, String(issuedAtMs)].join('|');
return createHmac('sha256', this.secret).update(message).digest('hex');
\end{verbatim}

Khi người học nộp đáp án (Hình~\ref{fig:seq-answer}), \code{LearnService} thực
hiện ba việc:

\begin{enumerate}
  \item \textbf{Xác minh chữ ký:} kiểm tra chữ ký HMAC bằng phép so sánh
        \textit{thời gian hằng số} (\code{timingSafeEqual}) và kiểm tra hạn
        dùng (mặc định 30 phút); chữ ký sai/hết hạn $\rightarrow$ lỗi 401.
  \item \textbf{Chấm điểm tại chỗ:} \code{AnswerGraderService} \textit{suy
        lại} đáp án đúng từ dữ liệu từ vựng (dạng từ bị khoét trong câu, bản
        dịch của nghĩa, hay chính lemma) rồi so khớp. Điểm chất lượng SM-2
        được quy đổi theo loại câu: trả lời nhanh và đúng được 5, đúng nhưng
        chậm được 4, sai lệch một ký tự (đo bằng khoảng cách Levenshtein) được
        coi là "gần đúng" (3), còn lại 2; riêng thẻ học cho người dùng tự đánh
        giá theo thang bốn mức.
  \item \textbf{Lập lịch ở bước cuối:} mỗi bài học chỉ là \textit{một} sự kiện
        SRS --- chỉ câu hỏi cuối (\code{stepIndex === stepCount - 1}) mới gọi
        \code{submitReview} để chạy \code{applySm2} và cập nhật lịch; các bước
        trước chỉ chấm để phản hồi tức thì. Nhờ vậy một từ không thể "tốt
        nghiệp" chỉ trong một lượt ngồi học.
\end{enumerate}

Việc cập nhật tiến độ và ghi nhật ký hoạt động được thực hiện trong cùng một
giao dịch (mục~\ref{sec:hien-thuc-du-lieu}). Nếu lịch vừa được dời vào trong
\textit{cửa sổ đưa lại} (15 phút), phản hồi sẽ kèm sẵn câu hỏi kế đã ký để máy
khách hiển thị lại từ ngay trong phiên mà không cần gọi lại API tạo phiên.

% =====================================================================
\section{Các tính năng trí tuệ nhân tạo và tự động hoá nội dung}
\label{sec:ai-tu-dong}

Nhóm tính năng thứ hai làm nên tính "thông minh" của hệ thống là việc
\textbf{tự động tạo và làm giàu nội dung} bằng các dịch vụ AI, giải quyết
trực tiếp rào cản "tạo nội dung thủ công tốn công" đã nêu ở
Chương~\ref{chuong:tong-quan}. Tất cả các tác vụ này đều nặng và phụ thuộc
dịch vụ ngoài, nên được hiện thực trên một \textit{hạ tầng xử lý nền} chung.

\subsection{Hạ tầng xử lý nền bất đồng bộ}
\label{subsec:impl-worker}

Mọi tác vụ AI được tổ chức theo mẫu \textbf{nhà sản xuất--người tiêu thụ}
trên \textbf{BullMQ}~\cite{bullmq2024}: tiến trình API tạo bản ghi công việc,
đẩy vào hàng đợi Redis rồi trả về mã \code{202} kèm mã công việc; tiến trình
worker tiêu thụ và xử lý nền (xem Hình~\ref{fig:activity-ai} và
Hình~\ref{fig:activity-practice}). Bảng~\ref{tab:async-jobs} liệt kê các hàng
đợi đã hiện thực. Bốn cơ chế bảo đảm độ tin cậy (NFR-04) được áp dụng nhất
quán cho mọi worker:

\begin{itemize}
  \item \textbf{Tính idempotent:} mỗi worker đọc lại bản ghi và \textit{chỉ}
        xử lý khi trạng thái còn là \code{pending}; một công việc bị chạy lại
        sẽ thoát sớm thay vì tạo dữ liệu trùng.
  \item \textbf{Thử lại có khoảng lùi:} lỗi tạm thời (mạng, 429, một phản hồi
        không phân tích được) khiến BullMQ thử lại với \textit{backoff}; chỉ
        khi cạn số lần thử, sự kiện \code{onFailed} mới ghi trạng thái
        \code{failed}.
  \item \textbf{Tự giới hạn tần suất:} worker đặt \textit{limiter} (mặc định
        15 lượt/phút) và độ song song bằng 1 để không vượt hạn mức của khoá
        AI dùng chung (tầng miễn phí).
  \item \textbf{Hạn mức theo ngày:} các điểm cuối người dùng có hạn mức theo
        ngày (ví dụ 30 lượt chấm câu/người/ngày) để bảo vệ tài nguyên chung,
        trả về lỗi 429 khi vượt.
\end{itemize}

\begin{table}[h]
  \centering
  \small
  \renewcommand{\arraystretch}{1.3}
  \begin{tabular}{|p{3.0cm}|p{4.6cm}|p{6.4cm}|}
    \hline
    \textbf{Hàng đợi} & \textbf{Kích hoạt bởi} & \textbf{Công việc của worker} \\
    \hline
    Làm giàu từ vựng & Tạo nhanh / nhập hàng loạt & Tra từ điển + gọi LLM sinh nghĩa, ví dụ, bản dịch, CEFR; tạo bản ghi từ vựng. \\
    \hline
    Sinh âm thanh & Sau khi tạo/duyệt từ & Lấy bản ghi âm thật hoặc tổng hợp giọng nói, tải lên kho media. \\
    \hline
    Sinh hình ảnh & Quy trình làm giàu nghĩa & Tìm ảnh minh hoạ theo từ, tải lên kho media. \\
    \hline
    Chấm điểm luyện viết & Nộp bài luyện viết/nói & Gọi LLM chấm câu theo rubric, lưu điểm và nhận xét. \\
    \hline
  \end{tabular}
  \caption{Các hàng đợi xử lý nền và vai trò của worker}
  \label{tab:async-jobs}
\end{table}

\paragraph{Lưu ý về mô hình ngôn ngữ.} Các lời gọi LLM được thực hiện qua API
\code{generateContent} của nền tảng \textbf{Google AI Studio}, dùng họ mô
hình mở \textbf{Gemma} của Google DeepMind~\cite{gemma2025technical} làm nền
tảng tham chiếu. Vì mô hình được cấu hình qua biến môi trường
(\code{GEMMA_MODEL}), thiết kế là \textit{độc lập mô hình}; trong cấu hình
triển khai mặc định, do dòng \code{gemma-3-*} đã ngừng phục vụ trên tầng miễn
phí của AI Studio, hệ thống dùng một mô hình Gemini Flash nhẹ tương đương.
Một chi tiết hiện thực quan trọng để tiết kiệm hạn mức: các lời gọi đặt
\code{thinkingConfig.thinkingBudget = 0} nhằm tắt token "suy nghĩ" ẩn ---
nếu không, mô hình suy luận sẽ tiêu hết hạn mức đầu ra cho phần suy nghĩ và
cắt cụt JSON kết quả.

\subsection{Làm giàu dữ liệu từ điển tự động bằng LLM}
\label{subsec:impl-enrich}

Đây là tính năng AI trung tâm: từ một \textit{lemma} đơn lẻ, hệ thống tự
dựng nên một bản ghi từ vựng đầy đủ (use case ở Bảng~\ref{tab:uc-quickcreate}).
\code{EnrichmentProcessor} chọn một trong hai đường đi tuỳ dữ liệu sẵn có:

\begin{itemize}
  \item \textbf{Đường có từ điển (tiếng Anh):} truy vấn Free Dictionary
        API~\cite{freedictionaryapi} để lấy từ loại, định nghĩa, phiên âm
        \textbf{IPA} và từ đồng/trái nghĩa; sau đó LLM chỉ cần bổ sung một
        \textit{gloss} ngắn, \textit{tối thiểu hai câu ví dụ} cho mỗi nghĩa,
        một bản dịch ngắn và ước lượng bậc CEFR. Cách kết hợp này tận dụng dữ
        liệu từ điển chính xác cho phần "khó tin tưởng AI" (định nghĩa, phiên
        âm) và chỉ giao cho AI phần nó làm tốt (ví dụ tự nhiên, bản dịch).
  \item \textbf{Đường chỉ dùng LLM (phi tiếng Anh hoặc từ điển không có):}
        LLM sinh toàn bộ cấu trúc nghĩa theo từng từ loại.
\end{itemize}

Vì Gemma trên AI Studio không hỗ trợ \code{responseSchema} hay vai trò
\textit{hệ thống} (system role), lược đồ JSON mong muốn được \textbf{nhúng
thẳng vào prompt} và phản hồi được \textbf{phân tích phòng thủ}. Bộ phân tích
bóc bỏ dấu rào ```\code{```json}```, dung thứ phần văn xuôi thừa, cắt và giới
hạn độ dài các trường, và --- quan trọng nhất --- \textit{từ chối} một CEFR
sai hoặc thiếu hay một nghĩa có dưới hai ví dụ bằng cách ném lỗi, để worker
\textit{thử lại} thay vì lưu dữ liệu rác. Đoạn dưới minh hoạ lược đồ được yêu
cầu và một phần kiểm tra phòng thủ:

\begin{verbatim}
// Lược đồ JSON nhúng trong prompt (rút gọn)
{ "cefr": "<A1|A2|B1|B2|C1|C2>",
  "senses": [ { "gloss": "<nhãn ngắn>",
                "examples": ["<câu>", "<câu>"],
                "translation": "<bản dịch ngắn>" } ] }

// src/vocabularies/enrichment/gemma-enricher.ts — kiểm tra phòng thủ
if (!VALID_CEFR.has(cefr))
  throw new Error(`enricher returned invalid cefr: ${cefr}`);
if (examples.length < 2)
  throw new Error(`sense ${i + 1} returned fewer than 2 examples`);
\end{verbatim}

\paragraph{Vòng kiểm soát chất lượng.} Kết quả do AI sinh ra luôn được kiểm soát theo
\textit{quyền sở hữu}, đúng tinh thần mục~\ref{subsec:llm-theory}: công việc do
\textit{quản trị viên} tạo sẽ sinh ra \textbf{bản nháp} thuộc kho hệ thống ở
trạng thái \code{is_approved = false} --- không xuất hiện trong phiên học của
ai cho tới khi được duyệt; còn từ do \textit{người học} tạo là từ cá nhân,
riêng tư và tự duyệt. Nhờ \textit{chỉ mục duy nhất bộ phận}
(mục~\ref{sec:hien-thuc-du-lieu}), worker bỏ qua mọi từ loại đã tồn tại nên
không bao giờ ghi đè dữ liệu sẵn có. Các từ được duyệt/tự duyệt tiếp tục kích
hoạt hàng đợi sinh âm thanh.

\subsection{Nhập từ vựng hàng loạt tự động}
\label{subsec:impl-bulk}

Tính năng nhập hàng loạt mở rộng cơ chế làm giàu ở quy mô lớn. Hàm
\code{extractCandidates} (\code{enrichment/import/}) nhận đầu vào ở nhiều định
dạng --- văn bản dán trực tiếp, CSV, bảng tính \textbf{XLSX} hay tệp
\textbf{PDF} --- quy về một khối văn bản rồi đưa qua bộ \textit{tách token}
chung. Bộ tách hỗ trợ hai chế độ: chế độ \textit{danh sách} (mỗi dòng/ô là một
mục) và chế độ \textit{văn xuôi} (tách từ, chuẩn hoá, loại bỏ \textit{từ dừng});
kết quả là danh sách \textit{ứng viên lemma}. Sau khi loại các từ đã có trong
kho, mỗi lemma còn lại được \textbf{phát tán} thành một công việc làm giàu
độc lập (chia sẻ một \code{batch_id}, và có thể kèm \code{target_deck_id} để
tự thêm từ mới vào một bộ thẻ đích sau khi tạo xong). Người dùng theo dõi
tiến độ cả lô qua mã lô thay vì phải chờ đồng bộ.

\subsection{Sinh âm thanh và hình ảnh minh hoạ tự động}
\label{subsec:impl-media}

Hai loại media được sinh tự động qua các pipeline có \textit{phương án dự
phòng} nhiều tầng, kết quả đều được lưu trên kho media \textbf{Cloudinary}~\cite{cloudinary2024}:

\begin{itemize}
  \item \textbf{Âm thanh phát âm:} ưu tiên lấy \textit{bản ghi âm người thật}
        từ Free Dictionary API; nếu không có, tổng hợp bằng \textbf{Microsoft
        Edge TTS} (giọng nơ-ron miễn phí, không cần khoá). Hàm trả về cả nguồn
        đã dùng (\code{dictionary} hay \code{tts}), giúp ưu tiên chất lượng cho
        các từ phổ biến và chỉ dùng TTS cho phần đuôi dài.
  \item \textbf{Hình ảnh minh hoạ nghĩa:} tìm ảnh phù hợp qua kho ảnh
        \textbf{Pexels}~\cite{pexels2024} rồi sao lưu lên Cloudinary. Với các
        từ trừu tượng không có ảnh phù hợp, pipeline trả về \code{null} và để
        trống cho quản trị viên bổ sung --- tự động hoá phần "đuôi dài" của
        các từ cụ thể mà không tạo ra ảnh sai lệch.
\end{itemize}

\subsection{Chấm điểm câu luyện tập bằng LLM}
\label{subsec:impl-judge}

Tính năng luyện viết câu (use case ở Bảng~\ref{tab:uc-practice}) dùng LLM làm
\textit{giám khảo}. Khi người học nộp một câu chứa từ mục tiêu,
\code{PracticeService} kiểm tra hạn mức ngày, tạo bản ghi \code{pending}, đẩy
vào hàng đợi và trả \code{202} ngay; nếu việc đẩy hàng đợi thất bại, bản ghi
được \textit{hoàn tác} để không "kẹt" ở trạng thái chờ và không tiêu phí hạn
mức. Worker \code{ScoringProcessor} gọi hàm \code{scoreSentence}
(\code{gemma-judge.ts}) với một prompt yêu cầu chấm theo \textit{rubric} có
cấu trúc (Bảng~\ref{tab:rubric-fields}); nhiệt độ sinh đặt thấp (0.2) để kết
quả ổn định. Điểm đặc thù: mô hình được yêu cầu ước lượng
\textbf{bậc CEFR mà chính câu đó thể hiện} (không phải năng lực tổng thể của
người học), cùng một câu gợi ý sửa.

\begin{table}[h]
  \centering
  \small
  \renewcommand{\arraystretch}{1.3}
  \begin{tabular}{|p{3.6cm}|p{3.0cm}|p{7.4cm}|}
    \hline
    \textbf{Trường rubric} & \textbf{Kiểu} & \textbf{Ý nghĩa} \\
    \hline
    \code{overall} & số nguyên 0--100 & Điểm tổng của câu. \\
    \hline
    \code{usesTargetWord} & boolean & Câu có dùng đúng từ mục tiêu (kể cả dạng biến cách). \\
    \hline
    \code{correctUsage} & boolean & Từ được dùng đúng một trong các nghĩa dự kiến. \\
    \hline
    \code{criteria} & 4 tiêu chí 0--5 & Ngữ pháp, dùng từ, độ tự nhiên, độ liên quan. \\
    \hline
    \code{cefr} & A1--C2 & Bậc CEFR mà câu thể hiện. \\
    \hline
    \code{feedback} & chuỗi & Một, hai câu nhận xét cho người học. \\
    \hline
    \code{correctedSentence} & chuỗi, tuỳ chọn & Phiên bản câu được cải thiện (nếu cần). \\
    \hline
  \end{tabular}
  \caption{Cấu trúc rubric chấm điểm câu luyện tập do LLM trả về}
  \label{tab:rubric-fields}
\end{table}

Phản hồi của mô hình được phân tích phòng thủ giống như ở phần làm giàu: các
khoảng số bị kẹp về biên hợp lệ, CEFR sai bị từ chối để worker thử lại. Người
học thăm dò kết quả qua mã bài luyện.

\subsection{Chấm điểm phát âm qua vi dịch vụ}
\label{subsec:impl-pron}

Chức năng luyện phát âm (Hình~\ref{fig:seq-pron}) được hiện thực bằng cách
tách phần xử lý tín hiệu/học máy thành một \textbf{vi dịch vụ} riêng viết bằng
Python (\textbf{FastAPI})~\cite{fastapi2024}, còn backend chỉ đóng vai
\textit{proxy mỏng}. \code{PronunciationScoringClient} đóng gói bản ghi âm
của người học thành một yêu cầu \code{multipart/form-data} (trường
\code{audio} và \code{word}), gọi điểm cuối \code{/score} với
\textit{thời gian chờ} ràng buộc, rồi ánh xạ lỗi của dịch vụ ngoài về mã lỗi
HTTP có ngữ nghĩa: \code{422} (âm thanh quá ngắn / không chấm được)
$\rightarrow$ \code{400}, còn không kết nối được / lỗi máy chủ $\rightarrow$
\code{503} (suy giảm có kiểm soát, NFR-04). Kết quả --- điểm tổng và điểm từng
âm vị --- được lưu vào \code{pronunciation_attempts} làm lịch sử luyện tập.
Việc cô lập này giữ cho tiến trình backend chính không phải mang các phụ thuộc
nặng về xử lý âm thanh.

% =====================================================================
\section{Theo dõi tiến độ và tạo động lực tự động}
\label{sec:tien-do}

Phân hệ tiến độ biến nhật ký học thành các chỉ số động lực một cách tự động
(các yêu cầu FR-19, FR-20). Nguồn dữ liệu là bảng \textit{chỉ-thêm}
\code{learning_activity} --- mỗi lượt ôn ghi đúng một dòng, được ghi cùng giao
dịch với cập nhật SRS (mục~\ref{subsec:impl-hmac}) nên các con số luôn nhất
quán với nhau. \code{ProgressService} tính:

\begin{itemize}
  \item \textbf{Thống kê trang chủ:} số từ theo từng trạng thái
        (\code{new}/\code{learning}/\code{review}/\code{mastered}), số thẻ đang
        đến hạn, số lượt ôn hôm nay và thời điểm đến hạn kế tiếp.
  \item \textbf{Chuỗi ngày học (streak):} số ngày liên tiếp (theo lịch UTC) có
        ít nhất một lượt ôn, tính lùi từ ngày ôn gần nhất; chuỗi chỉ còn hiệu
        lực nếu ngày gần nhất là hôm nay hoặc hôm qua.
  \item \textbf{Biểu đồ hoạt động (heatmap):} gộp số lượt ôn theo từng ngày
        \textit{theo múi giờ của người dùng} bằng một truy vấn SQL dùng
        \code{AT TIME ZONE}, để các ô lịch khớp với nửa đêm địa phương của họ.
\end{itemize}

\textbf{Bảng xếp hạng} được \code{LeaderboardService} tính trực tiếp: xếp hạng
người học theo \textit{số từ đã thành thạo} (toàn thời gian), chỉ tính những
tài khoản người học thật, đang hoạt động và \textit{không chọn ẩn}
(\code{leaderboard_opt_out}); thứ hạng của bản thân người gọi được suy ra từ
cùng một tập đã sắp xếp để hai con số không bao giờ mâu thuẫn. Bảng xếp hạng
theo \textit{số từ mới} được dành cho giai đoạn phát triển tiếp theo.

% =====================================================================
\section{Hiện thực hoá tầng giao diện người dùng}
\label{sec:hien-thuc-frontend}

Toàn bộ các API trình bày ở những mục trên được tiêu thụ bởi một ứng dụng giao
diện người dùng \textbf{độc lập}, hiện thực hoá thiết kế giao diện đã phác thảo
ở mục~\ref{sec:thiet-ke-giao-dien} (Chương~\ref{chuong:phan-tich-thiet-ke}).
Ứng dụng được xây dựng bằng \textbf{Next.js~16} theo mô hình \textit{App
Router} với \textbf{React~19} và \textbf{TypeScript} (mục~\ref{sec:cong-nghe-frontend}).
Đặc trưng kiến trúc xuyên suốt là \textit{ưu tiên phía máy chủ}: phần lớn màn
hình là Server Components lấy dữ liệu ngay trên máy chủ, còn mọi thao tác ghi
đi qua \textit{Server Actions}; nhờ đó access token và refresh token không bao
giờ lộ ra mã JavaScript phía trình duyệt. Bảng~\ref{tab:frontend-layout} ánh xạ
các phần chính của ứng dụng tới tệp hiện thực then chốt.

\begin{table}[h]
  \centering
  \small
  \renewcommand{\arraystretch}{1.3}
  \begin{tabular}{|p{3.6cm}|p{4.2cm}|p{6.0cm}|}
    \hline
    \textbf{Phần} & \textbf{Thư mục / tệp} & \textbf{Vai trò} \\
    \hline
    Nhóm tuyến & \code{src/app/(auth)}, \code{(app)}, \code{(admin)} & Phân tách khu vực xác thực, ứng dụng học và quản trị bằng \textit{route group}. \\
    \hline
    Tầng truy cập API & \code{src/lib/api.ts} & Client HTTP chạy phía máy chủ: \code{apiRequest} và \code{authedRequest}. \\
    \hline
    Quản lý phiên & \code{src/lib/auth/} & Lưu/đọc token trong cookie httpOnly, xoay token, đọc hồ sơ người dùng. \\
    \hline
    Kiểm tra dữ liệu & \code{src/lib/validations/} & Lược đồ \code{zod} dùng chung cho biểu mẫu và Server Action. \\
    \hline
    Động cơ phiên học & \code{src/app/(app)/learn/} & Bộ rút gọn \code{session-machine.ts} và các thành phần câu hỏi. \\
    \hline
  \end{tabular}
  \caption{Ánh xạ các phần frontend tới tệp hiện thực then chốt}
  \label{tab:frontend-layout}
\end{table}

\subsection{Cấu trúc App Router và phân nhóm tuyến}
\label{subsec:fe-approuter}

Cây tuyến được tổ chức bằng \textit{nhóm tuyến} (route group --- thư mục đặt
trong ngoặc đơn, không xuất hiện trong URL) để mỗi khu vực có bố cục và lớp bảo
vệ riêng: \code{(auth)} cho đăng nhập/đăng ký/xác minh email, \code{(app)} cho
trải nghiệm học của người dùng (trang chủ, học, khám phá, luyện tập, cộng đồng,
hồ sơ) và \code{(admin)} cho khu vực quản trị từ vựng. Mỗi phân đoạn tận dụng
các tệp đặc biệt của App Router --- \code{layout.tsx} (bố cục dùng chung),
\code{loading.tsx} (trạng thái chờ qua \textit{Suspense}), \code{error.tsx}
(ranh giới lỗi cấp phân đoạn) và \code{not-found.tsx} --- nên một lỗi cục bộ
không làm trắng toàn bộ ứng dụng. Các thành phần chỉ phục vụ một tuyến được đặt
ngay trong thư mục tuyến đó, chỉ nâng lên \code{src/components/} khi dùng lại ở
nhiều nơi.

\subsection{Tầng truy cập API phía máy chủ}
\label{subsec:fe-api}

Mọi lời gọi tới backend đi qua một client HTTP \textit{chỉ chạy phía máy chủ}
trong \code{src/lib/api.ts}. Địa chỉ backend (\code{API_BASE_URL}) được đọc từ
biến môi trường phía máy chủ, cố tình không gắn tiền tố \code{NEXT_PUBLIC_}
để không rò rỉ ra phía máy khách. Module phơi bày hai điểm vào:

\begin{itemize}
  \item \code{apiRequest} --- gọi \textit{không cần xác thực} (đăng nhập, đăng
        ký, làm mới token); tự đặt \code{Content-Type} JSON hoặc để runtime tự
        đặt biên \textit{multipart} khi thân là \code{FormData} (tải tệp lên).
  \item \code{authedRequest} --- đính kèm access token và, khi gặp
        \code{401}, \textit{tự động} làm mới refresh token đúng một lần rồi thử
        lại; chỉ gọi từ Server Action hoặc Route Handler vì nó có thể ghi lại
        cookie token đã xoay.
\end{itemize}

Cả hai chuẩn hoá thân lỗi của NestJS thành một kiểu \code{ApiResult} thống nhất
(cờ \code{ok}, mã trạng thái, dữ liệu, thân lỗi), nhờ đó tầng gọi xử lý lỗi
theo \textit{giá trị trả về} thay vì ngoại lệ --- khớp với triết lý phân biệt
lỗi \textit{dự kiến} và \textit{không dự kiến} của App Router.

\subsection{Quản lý phiên và xác thực phía máy khách}
\label{subsec:fe-session}

Token phiên được lưu trong \textbf{cookie httpOnly} với vòng đời bám đúng hợp
đồng backend (mục~\ref{sec:hien-thuc-auth}): access token 15~phút, refresh
token 30~ngày và được xoay sau mỗi lần gọi \code{/v1/auth/refresh}. Vì các
helper phiên import \code{next/headers}, chúng \textit{chỉ biên dịch được phía
máy chủ} --- một lớp bảo vệ ở mức biên dịch ngăn token lọt vào Client Component.
Khi \code{authedRequest} gặp \code{401}, nó âm thầm đổi refresh token lấy cặp
token mới, ghi lại cookie và thử lại đúng một lần; nếu vẫn thất bại, phiên bị
xoá và người dùng được chuyển về trang đăng nhập. Sau khi xác thực thành công,
Server Action điều hướng người dùng theo vai trò và trạng thái khai báo hồ sơ:
quản trị viên tới khu vực \code{/admin}, người học đã khai báo tới
\code{/dashboard}, người học chưa khai báo tới luồng \code{/onboarding}.

\subsection{Động cơ phiên học phía máy khách}
\label{subsec:fe-learn}

Phần tương tác nhất của ứng dụng là vòng lặp phiên học. Hàng đợi câu hỏi được
điều phối bằng một \textbf{bộ rút gọn (reducer) thuần} trong
\code{session-machine.ts}, tách hẳn khỏi React và mạng nên kiểm thử đơn vị được
độc lập. Nó nắm danh sách câu hỏi còn lại (câu đang hiển thị luôn là
\code{queue[0]}), các bộ đếm đúng/đã trả lời và trạng thái kết thúc
(\code{active}, \code{empty}, \code{done}, \code{expired}, \code{error}). Khi
một thẻ cần lặp lại trong cùng phiên, nó được \textit{nối lại vào cuối} hàng
đợi thay vì hẹn giờ theo đồng hồ thực. Bộ rút gọn này được ghép với hai Server
Action --- \code{startSessionAction} và \code{submitAnswerAction} --- gọi tới
\code{/v1/me/learn/session} và \code{/v1/me/learn/answer}. Vì máy chủ \textit{ký
HMAC} từng câu hỏi (mục~\ref{subsec:impl-hmac}), một \code{401} còn lại sau khi
đã làm mới token được hiểu là chữ ký câu hỏi đã hết hạn hoặc bị giả mạo, và giao
diện chuyển sang trạng thái \code{expired} mời người học bắt đầu phiên mới. Mỗi
\textit{chế độ} câu hỏi (thẻ ghi nhớ, điền khuyết trắc nghiệm/gõ chữ, nghe chép
chính tả, chọn theo hình/âm thanh, luyện phát âm\ldots) là một thành phần riêng
dưới \code{learn/questions/}, nhận dữ liệu đã ký từ máy chủ và chỉ chịu trách
nhiệm thu thập câu trả lời để nộp lại.

% =====================================================================
\section{Kiểm thử hệ thống}
\label{sec:kiem-thu}

Việc kiểm thử được hỗ trợ trực tiếp bởi một lựa chọn thiết kế lặp lại xuyên
suốt chương: \textbf{tách logic thông minh thành các hàm thuần, độc lập
framework}. Các hàm như \code{applySm2} (SM-2), bộ ký HMAC, bộ chấm đáp án,
bộ dựng prompt và bộ phân tích phản hồi LLM đều không phụ thuộc NestJS hay cơ
sở dữ liệu, nên có thể kiểm thử đơn vị nhanh và tất định. Chiến lược kiểm thử
gồm hai mức (Bảng~\ref{tab:test-targets}):

\begin{itemize}
  \item \textbf{Kiểm thử đơn vị} bằng \textbf{Jest}~\cite{jest2024}: tập trung
        vào các thuật toán và bộ phân tích --- nơi lỗi tinh vi dễ lọt và hậu
        quả lớn (lập lịch sai, lưu dữ liệu AI rác, chấm điểm sai).
  \item \textbf{Kiểm thử tích hợp (end-to-end)} bằng \textbf{Supertest}: gọi
        thật vào các điểm cuối API để kiểm tra luồng xác thực, kiểm tra đầu
        vào và mã trạng thái trả về.
\end{itemize}

\begin{table}[h]
  \centering
  \small
  \renewcommand{\arraystretch}{1.3}
  \begin{tabular}{|p{5.0cm}|p{9.4cm}|}
    \hline
    \textbf{Thành phần} & \textbf{Trọng tâm kiểm thử} \\
    \hline
    \code{srs.ts} & Chuyển trạng thái SM-2, bước học, tốt nghiệp, rớt bước, ngưỡng thành thạo. \\
    \hline
    \code{hmac-signer.service.ts} & Ký/xác minh, chống giả mạo, hết hạn, so sánh thời gian hằng số. \\
    \hline
    \code{answer-grader.service.ts} & Quy đổi điểm chất lượng cho từng loại câu hỏi. \\
    \hline
    \code{question-builder.service.ts} & Dựng bài học đúng dải khó, tôn trọng các trần và lọc dữ liệu. \\
    \hline
    \code{gemma-judge.ts}, \code{gemma-enricher.ts} & Phân tích phòng thủ JSON, kẹp số, từ chối CEFR sai. \\
    \hline
    \code{import/tokenize.ts}, \code{normalize.ts}, \code{xlsx-parser.ts} & Tách lemma, chuẩn hoá, đọc bảng tính cho nhập hàng loạt. \\
    \hline
  \end{tabular}
  \caption{Một số thành phần tiêu biểu và trọng tâm kiểm thử}
  \label{tab:test-targets}
\end{table}

Chất lượng mã nguồn còn được kiểm soát ở mức tĩnh bằng trình kiểm tra kiểu của
TypeScript (chế độ \textit{strict}), \textbf{ESLint} và \textbf{Prettier},
chạy trước mỗi lần commit (đáp ứng NFR-06).

% =====================================================================
\section{Kết chương}
\label{sec:ket-chuong-4}

Chương này đã trình bày quá trình hiện thực hoá hệ thống, với trọng tâm là các
thành phần thông minh và tự động hoá. Ở phía \textit{động cơ học}, hệ thống
thích nghi với từng người học qua thuật toán SM-2 mở rộng bước học
(\code{srs.ts}), cơ chế chọn từ theo khoảng cách CEFR, bộ sinh câu hỏi nhiều
tầng theo giai đoạn ghi nhớ với phương án nhiễu "dễ nhầm", và cơ chế chấm điểm
phi trạng thái an toàn bằng chữ ký HMAC. Ở phía \textit{nội dung}, các tính
năng AI --- làm giàu từ điển bằng LLM kết hợp từ điển, nhập hàng loạt từ nhiều
định dạng, sinh âm thanh và hình ảnh tự động, chấm điểm câu và phát âm --- đều
chạy trên một hạ tầng xử lý nền bất đồng bộ có tính idempotent, biết thử lại,
tự giới hạn tần suất và suy giảm có kiểm soát. Ở phía \textit{máy khách}, một
ứng dụng Next.js theo App Router tiêu thụ toàn bộ các API trên qua một tầng
truy cập phía máy chủ, quản lý phiên bằng cookie httpOnly với cơ chế xoay token
tự động, và điều phối phiên học bằng một bộ rút gọn thuần tách khỏi React. Tất
cả được củng cố bằng các chỉ số tiến độ tự động và một chiến lược kiểm thử dựa
trên việc tách logic thành các hàm thuần. Những kết quả hiện thực hoá này là cơ
sở để chương tiếp theo đánh giá hệ thống và rút ra kết luận.

%  vào main.tex (sau Chương 3).
%
%  GHI CHÚ VỀ HÌNH VẼ: dùng macro \figph{<tên-file>}{<tỉ-lệ-rộng>}{<chú-thích>}{<nhãn>}
%  (đã khai báo ở các chương trước; khai báo lại bằng \providecommand để an
%  toàn nếu biên dịch độc lập). Mã nguồn (code) trình bày bằng môi trường
%  verbatim như các chương trước.
%
%  KHÔNG hardcode số chương/mục/bảng/hình/công thức — luôn dùng \label + \ref / \cite.
% =====================================================================

% --- Macro tiện ích (idempotent — không định nghĩa lại nếu đã có) ------
\providecommand{\code}[1]{\texttt{\detokenize{#1}}} % in định danh snake_case an toàn
\providecommand{\figph}[4]{%
  \begin{figure}[htbp]
    \centering
    \IfFileExists{images/#1}%
      {\includegraphics[width=#2\textwidth,height=0.82\textheight,keepaspectratio]{#1}}%
      {\setlength{\fboxsep}{10pt}\fbox{\begin{minipage}[c][3.6cm][c]{#2\textwidth}\centering\itshape Sơ đồ minh hoạ --- chèn tệp ảnh \code{images/#1} vào thư mục \code{images/}.\end{minipage}}}%
    \caption{#3}%
    \label{#4}%
  \end{figure}}

\chapter{HIỆN THỰC HOÁ HỆ THỐNG VÀ CÁC TÍNH NĂNG THÔNG MINH}
\label{chuong:hien-thuc-hoa}

Hai chương trước đã trình bày kết quả phân tích, thiết kế
(Chương~\ref{chuong:phan-tich-thiet-ke}) và cơ sở lý thuyết, công nghệ
(Chương~\ref{chuong:co-so-ly-thuyet}). Chương này trình bày cách các thiết
kế đó được \textit{hiện thực hoá} thành mã nguồn chạy được, với trọng tâm là
các \textbf{thành phần thông minh và tự động hoá} làm nên giá trị cốt lõi của
hệ thống: động cơ học theo ôn tập ngắt quãng, cơ chế sinh câu hỏi và chấm
điểm thích nghi, các tính năng làm giàu nội dung và chấm điểm dựa trên trí
tuệ nhân tạo, cùng hạ tầng xử lý nền bất đồng bộ liên kết chúng lại.

Mục~\ref{sec:tong-quan-hien-thuc} mô tả tổ chức mã nguồn và các mối quan tâm
xuyên suốt. Mục~\ref{sec:hien-thuc-du-lieu} và
mục~\ref{sec:hien-thuc-auth} trình bày ngắn gọn hiện thực hoá tầng dữ liệu và
phân hệ xác thực. Trọng tâm của chương nằm ở
mục~\ref{sec:dong-co-hoc} (động cơ học thông minh) và
mục~\ref{sec:ai-tu-dong} (các tính năng AI và tự động hoá nội dung). Cuối
cùng, mục~\ref{sec:tien-do} trình bày phân hệ theo dõi tiến độ tự động,
mục~\ref{sec:hien-thuc-frontend} trình bày hiện thực hoá tầng giao diện người
dùng (frontend), mục~\ref{sec:kiem-thu} trình bày chiến lược kiểm thử và
mục~\ref{sec:ket-chuong-4} tóm tắt chương.

% =====================================================================
\section{Tổng quan quá trình hiện thực hoá}
\label{sec:tong-quan-hien-thuc}

Hệ thống được hiện thực hoá đúng theo ngăn xếp công nghệ đã chọn ở
Chương~\ref{chuong:co-so-ly-thuyet}: backend viết bằng \textbf{NestJS~11}
trên \textbf{TypeScript} (chế độ \textit{strict}), truy cập
\textbf{PostgreSQL} qua \textbf{TypeORM}, và dùng \textbf{Redis/BullMQ} cho
xử lý nền. Toàn bộ mã nguồn được biên dịch bằng trình biên dịch của NestJS và
khởi chạy dưới \textbf{hai tiến trình} chia sẻ chung cơ sở dữ liệu và hàng
đợi, đúng mô hình nhà sản xuất--người tiêu thụ đã nêu ở
mục~\ref{sec:kien-truc} của Chương~\ref{chuong:phan-tich-thiet-ke}:

\begin{itemize}
  \item \textbf{Tiến trình API} (\code{npm run start}): phục vụ các yêu cầu
        HTTP đồng bộ; với tác vụ nặng nó chỉ tạo bản ghi công việc rồi đẩy
        vào hàng đợi và trả về ngay.
  \item \textbf{Tiến trình Worker} (\code{npm run start:worker}): tiêu thụ
        hàng đợi để thực thi các tác vụ nền (làm giàu từ vựng, sinh âm thanh,
        chấm điểm luyện tập).
\end{itemize}

Mọi tham số vận hành (chuỗi kết nối CSDL, khoá bí mật, khoá API của dịch vụ
AI, hạn mức tần suất\ldots) đều được nạp từ \textbf{biến môi trường} qua
\code{@nestjs/config}, đáp ứng yêu cầu tính khả chuyển NFR-08 trong
Bảng~\ref{tab:nfr}: cùng một ảnh mã nguồn chạy được ở nhiều môi trường mà
không cần sửa code.

\subsection{Tổ chức mã nguồn theo mô-đun}
\label{subsec:to-chuc-ma-nguon}

Mỗi mô-đun nghiệp vụ ở Bảng~\ref{tab:modules} (Chương~\ref{chuong:phan-tich-thiet-ke})
được hiện thực thành một thư mục dưới \code{src/} tự đóng gói controller,
service, entity và DTO của mình. Bảng~\ref{tab:source-layout} ánh xạ các
mô-đun trọng yếu tới những tệp hiện thực then chốt --- là cơ sở để các mục
sau tham chiếu trực tiếp.

\begin{table}[h]
  \centering
  \small
  \renewcommand{\arraystretch}{1.3}
  \begin{tabular}{|p{2.6cm}|p{4.4cm}|p{6.6cm}|}
    \hline
    \textbf{Mô-đun} & \textbf{Thư mục mã nguồn} & \textbf{Tệp hiện thực then chốt} \\
    \hline
    \code{auth} & \code{src/auth} & \code{token.service.ts} (cấp/xoay token), \code{jwt.strategy.ts}, các \code{guards/}, \code{services/} (Google/Apple/GitHub). \\
    \hline
    \code{learn} & \code{src/learn} & \code{learn.service.ts}, \code{vocab-picker.service.ts}, \code{question-builder.service.ts}, \code{distractor.service.ts}, \code{answer-grader.service.ts}, \code{hmac-signer.service.ts}, \code{question-bands.ts}. \\
    \hline
    \code{progress} & \code{src/progress} & \code{srs.ts} (thuật toán SM-2), \code{progress.service.ts} (lập lịch, thống kê, chuỗi ngày học). \\
    \hline
    \code{vocabularies} & \code{src/vocabularies} & \code{enrichment/} (làm giàu bằng LLM, nhập hàng loạt), \code{audio/}, \code{images/}, \code{vocabulary-persistence.service.ts}. \\
    \hline
    \code{practice} & \code{src/practice} & \code{gemma-judge.ts} (chấm câu bằng LLM), \code{scoring.processor.ts}, \code{practice.service.ts}. \\
    \hline
    \code{pronunciation} & \code{src/pronunciation} & \code{pronunciation-scoring.client.ts} (proxy tới vi dịch vụ), \code{pronunciation.service.ts}. \\
    \hline
    \code{leaderboard} & \code{src/leaderboard} & \code{leaderboard.service.ts}. \\
    \hline
    \code{config} & \code{src/config} & Các tệp cấu hình theo miền: \code{learn.config.ts}, \code{gemma.config.ts}, \code{auth.config.ts}\ldots \\
    \hline
  \end{tabular}
  \caption{Ánh xạ mô-đun nghiệp vụ tới tệp hiện thực then chốt}
  \label{tab:source-layout}
\end{table}

\subsection{Các mối quan tâm xuyên suốt}
\label{subsec:cross-cutting}

Nhờ cơ chế \textit{tiêm phụ thuộc} và các thành phần khai báo của NestJS, các
mối quan tâm chung được hiện thực một lần và áp dụng toàn cục thay vì lặp lại
trong từng controller:

\begin{itemize}
  \item \textbf{Kiểm tra đầu vào:} một \code{ValidationPipe} toàn cục với
        \code{whitelist}, \code{forbidNonWhitelisted} và \code{transform}
        biến mỗi DTO (khai báo bằng \code{class-validator}) thành lớp bảo vệ:
        trường lạ bị loại bỏ, dữ liệu sai kiểu bị chặn bằng lỗi 400 trước khi
        chạm tới service.
  \item \textbf{Xác thực và phân quyền:} \code{JwtAuthGuard} (dựa trên
        \code{passport-jwt}, xem \code{jwt.strategy.ts}) trích và xác minh
        access token ở mọi endpoint cần đăng nhập; \code{RolesGuard} kết hợp
        decorator \code{@Roles('admin')} chặn các thao tác quản trị; quyền sở
        hữu tài nguyên được kiểm tra ngay trong service (mục~\ref{subsec:rbac}).
  \item \textbf{Phiên bản hoá API:} bật \textit{URI versioning} mặc định
        \code{v1}; controller khai báo \code{@Controller(\{ path, version: '1' \})}
        nên mọi tài nguyên nằm dưới tiền tố \code{/v1} (mục~\ref{subsec:rest-theory}).
  \item \textbf{Cấu hình theo miền:} mỗi nhóm tham số được gói trong một
        \textit{namespace} cấu hình riêng (ví dụ \code{learn.config.ts},
        \code{gemma.config.ts}), tách bạch và có giá trị mặc định an toàn cho
        môi trường phát triển.
\end{itemize}

% =====================================================================
\section{Hiện thực hoá tầng dữ liệu}
\label{sec:hien-thuc-du-lieu}

Tầng dữ liệu hiện thực 17 bảng trong Bảng~\ref{tab:db-overview} thành các lớp
\textit{thực thể} (entity) TypeORM. Mỗi thực thể khai báo khoá chính
\code{uuid}, tên cột \code{snake_case} và kiểu \code{timestamptz} cho các mốc
thời gian, đúng quy ước thiết kế ở Chương~\ref{chuong:phan-tich-thiet-ke}.
Một số kỹ thuật hiện thực đáng chú ý:

\begin{itemize}
  \item \textbf{Giao dịch bảo đảm tính nguyên tử:} các thao tác ghi nhiều
        bảng được gói trong \code{dataSource.transaction(...)}. Ví dụ tiêu
        biểu là việc nộp một lượt ôn (mục~\ref{subsec:impl-hmac}): bản ghi
        tiến độ và bản ghi nhật ký hoạt động được lưu trong cùng một giao
        dịch, nên một lượt ôn không thể "ghi nửa vời" --- đáp ứng NFR-04 trong
        Bảng~\ref{tab:nfr}.
  \item \textbf{Ràng buộc toàn vẹn và chống trùng lặp:} ngoài khoá ngoại và
        quy tắc xoá lan truyền, hệ thống dùng \textit{chỉ mục duy nhất bộ
        phận} (partial unique index) trên bộ ba \code{(language, lemma,
        part_of_speech)} áp dụng riêng cho từ thuộc kho hệ thống
        (\code{WHERE source = 'system'}); nhờ đó kho hệ thống không bao giờ có
        hai bản ghi trùng từ loại cho cùng một từ, trong khi từ cá nhân của
        người dùng vẫn không bị ràng buộc này.
  \item \textbf{Quản lý lược đồ bằng migration:} mọi thay đổi lược đồ được mô
        tả bằng các \textit{migration} có thể phiên bản hoá, được nhận diện và
        chạy qua công cụ dòng lệnh của TypeORM. Cách này cho phép tái lập lược
        đồ một cách xác định giữa các môi trường (đáp ứng NFR-06).
  \item \textbf{Phi chuẩn hoá có kiểm soát:} một số trường được phi chuẩn hoá
        để tăng tốc đọc, ví dụ \code{decks.vocab_count} (số từ trong bộ thẻ)
        được cập nhật khi thêm/bớt thành viên thay vì đếm mỗi lần truy vấn.
\end{itemize}

% =====================================================================
\section{Hiện thực hoá xác thực và phân quyền}
\label{sec:hien-thuc-auth}

Phân hệ \code{auth} hiện thực các kỹ thuật bảo mật đã trình bày về lý thuyết
ở mục~\ref{sec:xac-thuc-phan-quyen}. Lớp \code{TokenService} đảm nhiệm vòng
đời phiên đăng nhập với một số quyết định hiện thực đáng chú ý:

\begin{itemize}
  \item \textbf{Cặp token bất đối xứng về bản chất:} \textit{access token} là
        một \textbf{JWT} được ký, tự chứa và sống ngắn (mục~\ref{subsec:jwt}),
        trong khi \textit{refresh token} là một chuỗi ngẫu nhiên 48 byte
        \textit{mờ} (opaque) --- không mang thông tin, chỉ lưu \textbf{bản băm
        SHA-256} trong bảng \code{refresh_tokens}. Máy chủ không bao giờ lưu
        token gốc, nên dù cơ sở dữ liệu bị lộ thì token cũng không thể tái sử
        dụng.
  \item \textbf{Xoay vòng và phát hiện tái dùng:} mỗi lần làm mới phiên, token
        cũ bị đánh dấu thu hồi và một token mới được cấp. Nếu một token đã thu
        hồi (hoặc không tồn tại) được trình lại --- dấu hiệu rò rỉ --- hệ thống
        \textbf{thu hồi toàn bộ} token còn hiệu lực của người dùng đó và trả
        về lỗi 401, hạn chế thiệt hại khi token bị đánh cắp.
  \item \textbf{Băm mật khẩu và đăng nhập bên thứ ba:} mật khẩu được băm bằng
        \textbf{bcrypt} (mục~\ref{subsec:bcrypt}); với OAuth/OIDC
        (mục~\ref{subsec:oauth}), token danh tính do nhà cung cấp cấp được xác
        minh rồi liên kết tới một bản ghi trong \code{user_identities}, cho
        phép một tài khoản đăng nhập bằng nhiều phương thức.
  \item \textbf{Chống vét cạn:} endpoint đăng nhập được giới hạn tần suất bằng
        \code{@nestjs/throttler}, trả về lỗi 429 khi vượt ngưỡng.
\end{itemize}

% =====================================================================
\section{Động cơ học thông minh theo ôn tập ngắt quãng}
\label{sec:dong-co-hoc}

Đây là phân hệ \textit{thông minh} cốt lõi của hệ thống. Một phiên học không
phải là một danh sách câu hỏi tĩnh: với mỗi người học và mỗi từ, hệ thống
\textbf{thích nghi} ở ba cấp độ --- chọn \textit{từ nào} để học, sinh
\textit{dạng câu hỏi nào} phù hợp giai đoạn ghi nhớ, và quyết định
\textit{khi nào} ôn lại. Toàn bộ quy trình tuân theo luồng đã mô hình hoá ở
Hình~\ref{fig:activity-srs}: chọn từ $\rightarrow$ mở rộng thành chuỗi câu hỏi
đã ký $\rightarrow$ chấm điểm $\rightarrow$ (ở bước cuối) lập lịch lại.

\subsection{Lập lịch ôn tập SM-2 mở rộng bước học}
\label{subsec:impl-sm2}

Hạt nhân lập lịch là hàm thuần \code{applySm2} trong \code{src/progress/srs.ts}:
nó nhận trạng thái hiện tại của một cặp (người dùng, từ), một
\textit{điểm chất lượng} \( q \in \{0,\dots,5\} \) và danh sách bước học, rồi
trả về trạng thái mới cùng thời điểm ôn kế tiếp. Hàm hiện thực thuật toán
\textbf{SM-2} (mục~\ref{subsec:sm2-theory}) được \textit{mở rộng bằng các bước
học} theo phong cách Anki. Hệ số dễ nhớ \( EF \) được cập nhật ở \textit{mọi}
lượt ôn theo công thức~\eqref{eq:sm2-ef}, còn khoảng cách ngày tuân theo
công thức~\eqref{eq:sm2-interval}; phần mở rộng bổ sung một thang
\textit{bước học} tính bằng phút để củng cố từ mới ngay trong phiên trước khi
giao cho thang ngày:

\begin{verbatim}
// src/progress/srs.ts — cập nhật hệ số dễ nhớ (chặn dưới 1.3)
function updateEaseFactor(easeFactor, quality) {
  return Math.max(
    1.3,
    easeFactor + (0.1 - (5 - quality) * (0.08 + (5 - quality) * 0.02)),
  );
}
\end{verbatim}

Logic chuyển trạng thái được tóm tắt như sau (Hình~\ref{fig:srs-states}):

\begin{itemize}
  \item Thẻ đang ở \textit{trạng thái bước} (\code{learningStepIndex} khác
        \code{null}) dùng khoảng cách \textbf{phút} lấy từ danh sách bước (mặc
        định \([1, 10]\) phút). Trả lời đúng (\( q \ge 3 \)) đẩy lên một bước;
        qua bước cuối, thẻ \textit{tốt nghiệp} sang thang ngày. Trả lời sai
        (\( q < 3 \)) đặt lại về bước 0.
  \item Thẻ đã tốt nghiệp dùng thang ngày: \( I = 1 \) ngày (lần đầu), \( 6 \)
        ngày (lần hai), sau đó \( I \leftarrow \lceil I \times EF \rceil \).
        Một lần trả lời sai sẽ đưa thẻ \textit{rớt} về bước 0 (ôn lại trong
        phút) thay vì chờ tới hôm sau.
  \item Trạng thái hiển thị được suy ra từ số lần lặp và khoảng cách:
        \( I \ge 90 \) ngày $\rightarrow$ \code{mastered}; số lần ôn đúng liên
        tiếp \( \ge 3 \) $\rightarrow$ \code{review}; còn lại \code{learning}.
\end{itemize}

Bảng~\ref{tab:srs-params} liệt kê các tham số then chốt và giá trị mặc định
của chúng (đều có thể cấu hình qua biến môi trường trong \code{learn.config.ts}).

\begin{table}[h]
  \centering
  \small
  \renewcommand{\arraystretch}{1.3}
  \begin{tabular}{|p{5.2cm}|p{3.0cm}|p{6.0cm}|}
    \hline
    \textbf{Tham số} & \textbf{Giá trị mặc định} & \textbf{Ý nghĩa} \\
    \hline
    Hệ số dễ nhớ khởi tạo \( EF_0 \) & 2.5 & Giá trị ban đầu của \( EF \). \\
    \hline
    Chặn dưới \( EF_{\min} \) & 1.3 & Ngăn khoảng ôn co lại quá nhanh. \\
    \hline
    Bước học (phút) & \([1, 10]\) & Khoảng cách trong phiên cho thẻ mới/đang học/vừa rớt. \\
    \hline
    Khoảng ngày sau tốt nghiệp & 1 ngày, 6 ngày & Hai mốc đầu của thang ngày. \\
    \hline
    Ngưỡng thành thạo & \( \ge 90 \) ngày & Khoảng ôn đạt mốc này $\rightarrow$ \code{mastered}. \\
    \hline
    Cửa sổ đưa lại trong phiên & 15 phút & Nếu đến hạn trong cửa sổ này, câu hỏi kế được trả kèm ngay. \\
    \hline
  \end{tabular}
  \caption{Tham số của thuật toán lập lịch SM-2 mở rộng}
  \label{tab:srs-params}
\end{table}

% --- PlantUML gợi ý cho Hình máy trạng thái SRS ----------------------
% @startuml
% [*] --> new
% new --> learning : ghi danh (bước 0)
% learning --> learning : sai (q<3) → bước 0\nhoặc đúng nhưng chưa tốt nghiệp
% learning --> review : đúng liên tiếp ≥ 3 lần
% review --> review : đúng → I ← ⌈I×EF⌉
% review --> learning : sai → rớt về bước 0
% review --> mastered : I ≥ 90 ngày
% mastered --> learning : sai → rớt về bước 0
% mastered --> [*]
% @enduml
\figph{srs-state-machine.png}{0.8}{Máy trạng thái vòng đời thẻ học theo SM-2 mở rộng}{fig:srs-states}

Một quyết định thiết kế quan trọng: \code{applySm2} là một \textbf{hàm thuần}
(pure function) --- không phụ thuộc NestJS, cơ sở dữ liệu hay thời gian thực
(mốc \code{now} được truyền vào) --- nên có thể kiểm thử đơn vị triệt để
(mục~\ref{sec:kiem-thu}).

\subsection{Chọn từ học thích nghi theo trình độ}
\label{subsec:impl-picker}

Lớp \code{VocabPickerService} chọn tập từ cho mỗi phiên theo bốn chế độ
(hằng ngày, theo chủ đề, theo bộ thẻ, ôn tập). Nguyên tắc chung là
\textit{ưu tiên từ đã đến hạn} (sắp theo \code{next_review_at} tăng dần),
sau đó bù thêm \textit{từ mới} cho đủ hạn mức phiên. Điểm "thông minh" nằm ở
cách chọn từ mới: thay vì lọc cứng theo trình độ (dễ dẫn tới phiên rỗng khi
người học đã học hết từ ở đúng bậc), hệ thống coi trình độ là một
\textit{ưu tiên mềm} --- sắp các từ chưa học theo \textbf{khoảng cách CEFR}
tới trình độ người học (gần nhất trước), rồi mới tới thứ hạng tần suất:

\begin{verbatim}
// src/learn/vocab-picker.service.ts — ưu tiên từ gần trình độ người học
qb.orderBy(`ABS(${cefrOrdinalCase('v.cefr_level')} - ${userIdx})`,
           'ASC', 'NULLS LAST')
  .addOrderBy('v.frequency_rank', 'ASC', 'NULLS LAST')
  .limit(limit);
\end{verbatim}

Nhờ đó người học đã cạn từ ở bậc của mình (hoặc các từ còn lại không gắn bậc
CEFR) vẫn nhận được các từ phù hợp nhất hiện có thay vì một phiên rỗng. Khi
thực sự không còn từ nào, picker trả về một \textit{lý do rỗng} có ngữ nghĩa
(\code{no_due_cards}, \code{no_more_at_level}, \code{deck_exhausted},
\code{no_enrollment}) để máy khách hiển thị thông điệp phù hợp; chỉ với
\code{no_due_cards} hệ thống mới kèm theo \textit{thời điểm đến hạn kế tiếp}.
Các từ mới được chọn sẽ được \textit{tự động ghi danh} (auto-enroll) vào hàng
đợi học của người dùng trước khi sinh câu hỏi.

\subsection{Sinh câu hỏi thích nghi theo giai đoạn ghi nhớ}
\label{subsec:impl-questions}

Mỗi từ trong phiên được \code{QuestionBuilderService} mở rộng thành một
\textit{bài học} (lesson) gồm một chuỗi câu hỏi xếp từ dễ đến khó. Mức độ khó
được tổ chức thành ba \textbf{dải khó} (difficulty band) và \textbf{12 dạng
câu hỏi} (Bảng~\ref{tab:question-ladder}). Điểm thích nghi cốt lõi:
\textit{giai đoạn ghi nhớ của từ quyết định những dải nào được dùng}. Từ càng
thành thạo thì "sàn" độ khó càng được nâng lên, loại dần các dạng nhận biết
dễ và tập trung vào sản sinh ngôn ngữ:

\begin{table}[h]
  \centering
  \small
  \renewcommand{\arraystretch}{1.3}
  \begin{tabular}{|p{2.6cm}|p{3.4cm}|p{7.6cm}|}
    \hline
    \textbf{Trạng thái từ} & \textbf{Dải khó tối thiểu} & \textbf{Nhóm câu hỏi đại diện} \\
    \hline
    \code{new} & Nhận biết (NEW) & Thẻ học (flashcard), điền khuyết trắc nghiệm, chọn từ/nghĩa, nghe chọn, nhìn hình chọn từ. \\
    \hline
    \code{learning} / \code{review} & Gợi nhớ (REVIEW) & Điền khuyết tự gõ, nghe gõ lại (dictation), luyện phát âm --- người học phải tự tạo ra dạng từ. \\
    \hline
    \code{mastered} & Sản sinh (MASTER) & Phân biệt nghĩa (sense disambiguation) --- dạng khó nhất, đòi hỏi phân biệt giữa các nghĩa của từ đa nghĩa. \\
    \hline
  \end{tabular}
  \caption{Ánh xạ giai đoạn ghi nhớ của từ tới dải khó và nhóm câu hỏi}
  \label{tab:question-ladder}
\end{table}

Trên nền tảng đó, bộ dựng câu hỏi áp dụng thêm các cơ chế giữ cho bài học
\textit{ngắn, đa dạng và khả thi}:

\begin{itemize}
  \item \textbf{Lấy mẫu xác định theo từ:} với mỗi dải, số dạng câu hỏi
        "tuỳ chọn" được giới hạn (mặc định 2 dạng/dải) và được chọn bằng phép
        \textit{xáo trộn xác định} theo định danh của từ. Nhờ vậy mỗi từ luyện
        một tập dạng câu hỏi khác nhau (thay vì mọi bài đều chạy hết các
        dạng), nhưng kết quả vẫn \textit{tái lập được}. Dạng thẻ học luôn được
        giữ làm bước nhập môn.
  \item \textbf{Trần họ điền khuyết:} các dạng cùng "khoét" một từ trong câu
        bị giới hạn số lần (mặc định 2) để một câu ví dụ không bị khoét nhiều
        bước liên tiếp.
  \item \textbf{Lọc theo dữ liệu sẵn có:} mỗi dạng chỉ được dựng khi đủ dữ
        liệu (ví dụ dạng nghe cần có \code{audio_url}, dạng nhìn hình cần ảnh
        nghĩa, dạng dịch cần bản dịch ở ngôn ngữ đích). Dạng nào thiếu dữ liệu
        sẽ bị bỏ qua một cách an toàn.
\end{itemize}

\paragraph{Sinh phương án nhiễu thông minh.} Chất lượng của một câu trắc
nghiệm phụ thuộc vào \textit{phương án nhiễu} (distractor). \code{DistractorService}
không lấy nhiễu ngẫu nhiên mà chọn các từ "dễ nhầm": \textbf{cùng từ loại,
cùng ngôn ngữ, cùng bậc CEFR \( \pm 1 \)}, và \textit{ưu tiên các từ chia sẻ
ít nhất một chủ đề} với từ mục tiêu; chỉ khi không đủ mới nới lỏng điều kiện.
Với dạng phân biệt nghĩa, nhiễu còn được lấy chính từ \textit{các nghĩa khác}
của cùng từ (bẫy đa nghĩa), buộc người học thực sự phân biệt ngữ cảnh.

\subsection{Chấm điểm phi trạng thái có ký HMAC}
\label{subsec:impl-hmac}

Như đã nêu ở mục~\ref{subsec:hmac}, hệ thống chấm điểm phiên học theo cách
\textbf{phi trạng thái}: máy chủ không lưu phiên trong bộ nhớ mà "đóng dấu"
tính toàn vẹn vào chính mỗi câu hỏi. Khi sinh câu hỏi, \code{HmacSignerService}
ký một thông điệp gồm các trường định danh và --- quan trọng --- cả
\code{stepIndex}/\code{stepCount}, để máy khách không thể giả mạo việc câu
nào là câu cuối (câu duy nhất làm thay đổi lịch ôn):

\begin{verbatim}
// src/learn/hmac-signer.service.ts — thông điệp được ký
const message = [userId, vocabularyId, type, exampleId,
                 translationLang ?? '', String(stepIndex),
                 String(stepCount), nonce, String(issuedAtMs)].join('|');
return createHmac('sha256', this.secret).update(message).digest('hex');
\end{verbatim}

Khi người học nộp đáp án (Hình~\ref{fig:seq-answer}), \code{LearnService} thực
hiện ba việc:

\begin{enumerate}
  \item \textbf{Xác minh chữ ký:} kiểm tra chữ ký HMAC bằng phép so sánh
        \textit{thời gian hằng số} (\code{timingSafeEqual}) và kiểm tra hạn
        dùng (mặc định 30 phút); chữ ký sai/hết hạn $\rightarrow$ lỗi 401.
  \item \textbf{Chấm điểm tại chỗ:} \code{AnswerGraderService} \textit{suy
        lại} đáp án đúng từ dữ liệu từ vựng (dạng từ bị khoét trong câu, bản
        dịch của nghĩa, hay chính lemma) rồi so khớp. Điểm chất lượng SM-2
        được quy đổi theo loại câu: trả lời nhanh và đúng được 5, đúng nhưng
        chậm được 4, sai lệch một ký tự (đo bằng khoảng cách Levenshtein) được
        coi là "gần đúng" (3), còn lại 2; riêng thẻ học cho người dùng tự đánh
        giá theo thang bốn mức.
  \item \textbf{Lập lịch ở bước cuối:} mỗi bài học chỉ là \textit{một} sự kiện
        SRS --- chỉ câu hỏi cuối (\code{stepIndex === stepCount - 1}) mới gọi
        \code{submitReview} để chạy \code{applySm2} và cập nhật lịch; các bước
        trước chỉ chấm để phản hồi tức thì. Nhờ vậy một từ không thể "tốt
        nghiệp" chỉ trong một lượt ngồi học.
\end{enumerate}

Việc cập nhật tiến độ và ghi nhật ký hoạt động được thực hiện trong cùng một
giao dịch (mục~\ref{sec:hien-thuc-du-lieu}). Nếu lịch vừa được dời vào trong
\textit{cửa sổ đưa lại} (15 phút), phản hồi sẽ kèm sẵn câu hỏi kế đã ký để máy
khách hiển thị lại từ ngay trong phiên mà không cần gọi lại API tạo phiên.

% =====================================================================
\section{Các tính năng trí tuệ nhân tạo và tự động hoá nội dung}
\label{sec:ai-tu-dong}

Nhóm tính năng thứ hai làm nên tính "thông minh" của hệ thống là việc
\textbf{tự động tạo và làm giàu nội dung} bằng các dịch vụ AI, giải quyết
trực tiếp rào cản "tạo nội dung thủ công tốn công" đã nêu ở
Chương~\ref{chuong:tong-quan}. Tất cả các tác vụ này đều nặng và phụ thuộc
dịch vụ ngoài, nên được hiện thực trên một \textit{hạ tầng xử lý nền} chung.

\subsection{Hạ tầng xử lý nền bất đồng bộ}
\label{subsec:impl-worker}

Mọi tác vụ AI được tổ chức theo mẫu \textbf{nhà sản xuất--người tiêu thụ}
trên \textbf{BullMQ}~\cite{bullmq2024}: tiến trình API tạo bản ghi công việc,
đẩy vào hàng đợi Redis rồi trả về mã \code{202} kèm mã công việc; tiến trình
worker tiêu thụ và xử lý nền (xem Hình~\ref{fig:activity-ai} và
Hình~\ref{fig:activity-practice}). Bảng~\ref{tab:async-jobs} liệt kê các hàng
đợi đã hiện thực. Bốn cơ chế bảo đảm độ tin cậy (NFR-04) được áp dụng nhất
quán cho mọi worker:

\begin{itemize}
  \item \textbf{Tính idempotent:} mỗi worker đọc lại bản ghi và \textit{chỉ}
        xử lý khi trạng thái còn là \code{pending}; một công việc bị chạy lại
        sẽ thoát sớm thay vì tạo dữ liệu trùng.
  \item \textbf{Thử lại có khoảng lùi:} lỗi tạm thời (mạng, 429, một phản hồi
        không phân tích được) khiến BullMQ thử lại với \textit{backoff}; chỉ
        khi cạn số lần thử, sự kiện \code{onFailed} mới ghi trạng thái
        \code{failed}.
  \item \textbf{Tự giới hạn tần suất:} worker đặt \textit{limiter} (mặc định
        15 lượt/phút) và độ song song bằng 1 để không vượt hạn mức của khoá
        AI dùng chung (tầng miễn phí).
  \item \textbf{Hạn mức theo ngày:} các điểm cuối người dùng có hạn mức theo
        ngày (ví dụ 30 lượt chấm câu/người/ngày) để bảo vệ tài nguyên chung,
        trả về lỗi 429 khi vượt.
\end{itemize}

\begin{table}[h]
  \centering
  \small
  \renewcommand{\arraystretch}{1.3}
  \begin{tabular}{|p{3.0cm}|p{4.6cm}|p{6.4cm}|}
    \hline
    \textbf{Hàng đợi} & \textbf{Kích hoạt bởi} & \textbf{Công việc của worker} \\
    \hline
    Làm giàu từ vựng & Tạo nhanh / nhập hàng loạt & Tra từ điển + gọi LLM sinh nghĩa, ví dụ, bản dịch, CEFR; tạo bản ghi từ vựng. \\
    \hline
    Sinh âm thanh & Sau khi tạo/duyệt từ & Lấy bản ghi âm thật hoặc tổng hợp giọng nói, tải lên kho media. \\
    \hline
    Sinh hình ảnh & Quy trình làm giàu nghĩa & Tìm ảnh minh hoạ theo từ, tải lên kho media. \\
    \hline
    Chấm điểm luyện viết & Nộp bài luyện viết/nói & Gọi LLM chấm câu theo rubric, lưu điểm và nhận xét. \\
    \hline
  \end{tabular}
  \caption{Các hàng đợi xử lý nền và vai trò của worker}
  \label{tab:async-jobs}
\end{table}

\paragraph{Lưu ý về mô hình ngôn ngữ.} Các lời gọi LLM được thực hiện qua API
\code{generateContent} của nền tảng \textbf{Google AI Studio}, dùng họ mô
hình mở \textbf{Gemma} của Google DeepMind~\cite{gemma2025technical} làm nền
tảng tham chiếu. Vì mô hình được cấu hình qua biến môi trường
(\code{GEMMA_MODEL}), thiết kế là \textit{độc lập mô hình}; trong cấu hình
triển khai mặc định, do dòng \code{gemma-3-*} đã ngừng phục vụ trên tầng miễn
phí của AI Studio, hệ thống dùng một mô hình Gemini Flash nhẹ tương đương.
Một chi tiết hiện thực quan trọng để tiết kiệm hạn mức: các lời gọi đặt
\code{thinkingConfig.thinkingBudget = 0} nhằm tắt token "suy nghĩ" ẩn ---
nếu không, mô hình suy luận sẽ tiêu hết hạn mức đầu ra cho phần suy nghĩ và
cắt cụt JSON kết quả.

\subsection{Làm giàu dữ liệu từ điển tự động bằng LLM}
\label{subsec:impl-enrich}

Đây là tính năng AI trung tâm: từ một \textit{lemma} đơn lẻ, hệ thống tự
dựng nên một bản ghi từ vựng đầy đủ (use case ở Bảng~\ref{tab:uc-quickcreate}).
\code{EnrichmentProcessor} chọn một trong hai đường đi tuỳ dữ liệu sẵn có:

\begin{itemize}
  \item \textbf{Đường có từ điển (tiếng Anh):} truy vấn Free Dictionary
        API~\cite{freedictionaryapi} để lấy từ loại, định nghĩa, phiên âm
        \textbf{IPA} và từ đồng/trái nghĩa; sau đó LLM chỉ cần bổ sung một
        \textit{gloss} ngắn, \textit{tối thiểu hai câu ví dụ} cho mỗi nghĩa,
        một bản dịch ngắn và ước lượng bậc CEFR. Cách kết hợp này tận dụng dữ
        liệu từ điển chính xác cho phần "khó tin tưởng AI" (định nghĩa, phiên
        âm) và chỉ giao cho AI phần nó làm tốt (ví dụ tự nhiên, bản dịch).
  \item \textbf{Đường chỉ dùng LLM (phi tiếng Anh hoặc từ điển không có):}
        LLM sinh toàn bộ cấu trúc nghĩa theo từng từ loại.
\end{itemize}

Vì Gemma trên AI Studio không hỗ trợ \code{responseSchema} hay vai trò
\textit{hệ thống} (system role), lược đồ JSON mong muốn được \textbf{nhúng
thẳng vào prompt} và phản hồi được \textbf{phân tích phòng thủ}. Bộ phân tích
bóc bỏ dấu rào ```\code{```json}```, dung thứ phần văn xuôi thừa, cắt và giới
hạn độ dài các trường, và --- quan trọng nhất --- \textit{từ chối} một CEFR
sai hoặc thiếu hay một nghĩa có dưới hai ví dụ bằng cách ném lỗi, để worker
\textit{thử lại} thay vì lưu dữ liệu rác. Đoạn dưới minh hoạ lược đồ được yêu
cầu và một phần kiểm tra phòng thủ:

\begin{verbatim}
// Lược đồ JSON nhúng trong prompt (rút gọn)
{ "cefr": "<A1|A2|B1|B2|C1|C2>",
  "senses": [ { "gloss": "<nhãn ngắn>",
                "examples": ["<câu>", "<câu>"],
                "translation": "<bản dịch ngắn>" } ] }

// src/vocabularies/enrichment/gemma-enricher.ts — kiểm tra phòng thủ
if (!VALID_CEFR.has(cefr))
  throw new Error(`enricher returned invalid cefr: ${cefr}`);
if (examples.length < 2)
  throw new Error(`sense ${i + 1} returned fewer than 2 examples`);
\end{verbatim}

\paragraph{Vòng kiểm soát chất lượng.} Kết quả do AI sinh ra luôn được kiểm soát theo
\textit{quyền sở hữu}, đúng tinh thần mục~\ref{subsec:llm-theory}: công việc do
\textit{quản trị viên} tạo sẽ sinh ra \textbf{bản nháp} thuộc kho hệ thống ở
trạng thái \code{is_approved = false} --- không xuất hiện trong phiên học của
ai cho tới khi được duyệt; còn từ do \textit{người học} tạo là từ cá nhân,
riêng tư và tự duyệt. Nhờ \textit{chỉ mục duy nhất bộ phận}
(mục~\ref{sec:hien-thuc-du-lieu}), worker bỏ qua mọi từ loại đã tồn tại nên
không bao giờ ghi đè dữ liệu sẵn có. Các từ được duyệt/tự duyệt tiếp tục kích
hoạt hàng đợi sinh âm thanh.

\subsection{Nhập từ vựng hàng loạt tự động}
\label{subsec:impl-bulk}

Tính năng nhập hàng loạt mở rộng cơ chế làm giàu ở quy mô lớn. Hàm
\code{extractCandidates} (\code{enrichment/import/}) nhận đầu vào ở nhiều định
dạng --- văn bản dán trực tiếp, CSV, bảng tính \textbf{XLSX} hay tệp
\textbf{PDF} --- quy về một khối văn bản rồi đưa qua bộ \textit{tách token}
chung. Bộ tách hỗ trợ hai chế độ: chế độ \textit{danh sách} (mỗi dòng/ô là một
mục) và chế độ \textit{văn xuôi} (tách từ, chuẩn hoá, loại bỏ \textit{từ dừng});
kết quả là danh sách \textit{ứng viên lemma}. Sau khi loại các từ đã có trong
kho, mỗi lemma còn lại được \textbf{phát tán} thành một công việc làm giàu
độc lập (chia sẻ một \code{batch_id}, và có thể kèm \code{target_deck_id} để
tự thêm từ mới vào một bộ thẻ đích sau khi tạo xong). Người dùng theo dõi
tiến độ cả lô qua mã lô thay vì phải chờ đồng bộ.

\subsection{Sinh âm thanh và hình ảnh minh hoạ tự động}
\label{subsec:impl-media}

Hai loại media được sinh tự động qua các pipeline có \textit{phương án dự
phòng} nhiều tầng, kết quả đều được lưu trên kho media \textbf{Cloudinary}~\cite{cloudinary2024}:

\begin{itemize}
  \item \textbf{Âm thanh phát âm:} ưu tiên lấy \textit{bản ghi âm người thật}
        từ Free Dictionary API; nếu không có, tổng hợp bằng \textbf{Microsoft
        Edge TTS} (giọng nơ-ron miễn phí, không cần khoá). Hàm trả về cả nguồn
        đã dùng (\code{dictionary} hay \code{tts}), giúp ưu tiên chất lượng cho
        các từ phổ biến và chỉ dùng TTS cho phần đuôi dài.
  \item \textbf{Hình ảnh minh hoạ nghĩa:} tìm ảnh phù hợp qua kho ảnh
        \textbf{Pexels}~\cite{pexels2024} rồi sao lưu lên Cloudinary. Với các
        từ trừu tượng không có ảnh phù hợp, pipeline trả về \code{null} và để
        trống cho quản trị viên bổ sung --- tự động hoá phần "đuôi dài" của
        các từ cụ thể mà không tạo ra ảnh sai lệch.
\end{itemize}

\subsection{Chấm điểm câu luyện tập bằng LLM}
\label{subsec:impl-judge}

Tính năng luyện viết câu (use case ở Bảng~\ref{tab:uc-practice}) dùng LLM làm
\textit{giám khảo}. Khi người học nộp một câu chứa từ mục tiêu,
\code{PracticeService} kiểm tra hạn mức ngày, tạo bản ghi \code{pending}, đẩy
vào hàng đợi và trả \code{202} ngay; nếu việc đẩy hàng đợi thất bại, bản ghi
được \textit{hoàn tác} để không "kẹt" ở trạng thái chờ và không tiêu phí hạn
mức. Worker \code{ScoringProcessor} gọi hàm \code{scoreSentence}
(\code{gemma-judge.ts}) với một prompt yêu cầu chấm theo \textit{rubric} có
cấu trúc (Bảng~\ref{tab:rubric-fields}); nhiệt độ sinh đặt thấp (0.2) để kết
quả ổn định. Điểm đặc thù: mô hình được yêu cầu ước lượng
\textbf{bậc CEFR mà chính câu đó thể hiện} (không phải năng lực tổng thể của
người học), cùng một câu gợi ý sửa.

\begin{table}[h]
  \centering
  \small
  \renewcommand{\arraystretch}{1.3}
  \begin{tabular}{|p{3.6cm}|p{3.0cm}|p{7.4cm}|}
    \hline
    \textbf{Trường rubric} & \textbf{Kiểu} & \textbf{Ý nghĩa} \\
    \hline
    \code{overall} & số nguyên 0--100 & Điểm tổng của câu. \\
    \hline
    \code{usesTargetWord} & boolean & Câu có dùng đúng từ mục tiêu (kể cả dạng biến cách). \\
    \hline
    \code{correctUsage} & boolean & Từ được dùng đúng một trong các nghĩa dự kiến. \\
    \hline
    \code{criteria} & 4 tiêu chí 0--5 & Ngữ pháp, dùng từ, độ tự nhiên, độ liên quan. \\
    \hline
    \code{cefr} & A1--C2 & Bậc CEFR mà câu thể hiện. \\
    \hline
    \code{feedback} & chuỗi & Một, hai câu nhận xét cho người học. \\
    \hline
    \code{correctedSentence} & chuỗi, tuỳ chọn & Phiên bản câu được cải thiện (nếu cần). \\
    \hline
  \end{tabular}
  \caption{Cấu trúc rubric chấm điểm câu luyện tập do LLM trả về}
  \label{tab:rubric-fields}
\end{table}

Phản hồi của mô hình được phân tích phòng thủ giống như ở phần làm giàu: các
khoảng số bị kẹp về biên hợp lệ, CEFR sai bị từ chối để worker thử lại. Người
học thăm dò kết quả qua mã bài luyện.

\subsection{Chấm điểm phát âm qua vi dịch vụ}
\label{subsec:impl-pron}

Chức năng luyện phát âm (Hình~\ref{fig:seq-pron}) được hiện thực bằng cách
tách phần xử lý tín hiệu/học máy thành một \textbf{vi dịch vụ} riêng viết bằng
Python (\textbf{FastAPI})~\cite{fastapi2024}, còn backend chỉ đóng vai
\textit{proxy mỏng}. \code{PronunciationScoringClient} đóng gói bản ghi âm
của người học thành một yêu cầu \code{multipart/form-data} (trường
\code{audio} và \code{word}), gọi điểm cuối \code{/score} với
\textit{thời gian chờ} ràng buộc, rồi ánh xạ lỗi của dịch vụ ngoài về mã lỗi
HTTP có ngữ nghĩa: \code{422} (âm thanh quá ngắn / không chấm được)
$\rightarrow$ \code{400}, còn không kết nối được / lỗi máy chủ $\rightarrow$
\code{503} (suy giảm có kiểm soát, NFR-04). Kết quả --- điểm tổng và điểm từng
âm vị --- được lưu vào \code{pronunciation_attempts} làm lịch sử luyện tập.
Việc cô lập này giữ cho tiến trình backend chính không phải mang các phụ thuộc
nặng về xử lý âm thanh.

% =====================================================================
\section{Theo dõi tiến độ và tạo động lực tự động}
\label{sec:tien-do}

Phân hệ tiến độ biến nhật ký học thành các chỉ số động lực một cách tự động
(các yêu cầu FR-19, FR-20). Nguồn dữ liệu là bảng \textit{chỉ-thêm}
\code{learning_activity} --- mỗi lượt ôn ghi đúng một dòng, được ghi cùng giao
dịch với cập nhật SRS (mục~\ref{subsec:impl-hmac}) nên các con số luôn nhất
quán với nhau. \code{ProgressService} tính:

\begin{itemize}
  \item \textbf{Thống kê trang chủ:} số từ theo từng trạng thái
        (\code{new}/\code{learning}/\code{review}/\code{mastered}), số thẻ đang
        đến hạn, số lượt ôn hôm nay và thời điểm đến hạn kế tiếp.
  \item \textbf{Chuỗi ngày học (streak):} số ngày liên tiếp (theo lịch UTC) có
        ít nhất một lượt ôn, tính lùi từ ngày ôn gần nhất; chuỗi chỉ còn hiệu
        lực nếu ngày gần nhất là hôm nay hoặc hôm qua.
  \item \textbf{Biểu đồ hoạt động (heatmap):} gộp số lượt ôn theo từng ngày
        \textit{theo múi giờ của người dùng} bằng một truy vấn SQL dùng
        \code{AT TIME ZONE}, để các ô lịch khớp với nửa đêm địa phương của họ.
\end{itemize}

\textbf{Bảng xếp hạng} được \code{LeaderboardService} tính trực tiếp: xếp hạng
người học theo \textit{số từ đã thành thạo} (toàn thời gian), chỉ tính những
tài khoản người học thật, đang hoạt động và \textit{không chọn ẩn}
(\code{leaderboard_opt_out}); thứ hạng của bản thân người gọi được suy ra từ
cùng một tập đã sắp xếp để hai con số không bao giờ mâu thuẫn. Bảng xếp hạng
theo \textit{số từ mới} được dành cho giai đoạn phát triển tiếp theo.

% =====================================================================
\section{Hiện thực hoá tầng giao diện người dùng}
\label{sec:hien-thuc-frontend}

Toàn bộ các API trình bày ở những mục trên được tiêu thụ bởi một ứng dụng giao
diện người dùng \textbf{độc lập}, hiện thực hoá thiết kế giao diện đã phác thảo
ở mục~\ref{sec:thiet-ke-giao-dien} (Chương~\ref{chuong:phan-tich-thiet-ke}).
Ứng dụng được xây dựng bằng \textbf{Next.js~16} theo mô hình \textit{App
Router} với \textbf{React~19} và \textbf{TypeScript} (mục~\ref{sec:cong-nghe-frontend}).
Đặc trưng kiến trúc xuyên suốt là \textit{ưu tiên phía máy chủ}: phần lớn màn
hình là Server Components lấy dữ liệu ngay trên máy chủ, còn mọi thao tác ghi
đi qua \textit{Server Actions}; nhờ đó access token và refresh token không bao
giờ lộ ra mã JavaScript phía trình duyệt. Bảng~\ref{tab:frontend-layout} ánh xạ
các phần chính của ứng dụng tới tệp hiện thực then chốt.

\begin{table}[h]
  \centering
  \small
  \renewcommand{\arraystretch}{1.3}
  \begin{tabular}{|p{3.6cm}|p{4.2cm}|p{6.0cm}|}
    \hline
    \textbf{Phần} & \textbf{Thư mục / tệp} & \textbf{Vai trò} \\
    \hline
    Nhóm tuyến & \code{src/app/(auth)}, \code{(app)}, \code{(admin)} & Phân tách khu vực xác thực, ứng dụng học và quản trị bằng \textit{route group}. \\
    \hline
    Tầng truy cập API & \code{src/lib/api.ts} & Client HTTP chạy phía máy chủ: \code{apiRequest} và \code{authedRequest}. \\
    \hline
    Quản lý phiên & \code{src/lib/auth/} & Lưu/đọc token trong cookie httpOnly, xoay token, đọc hồ sơ người dùng. \\
    \hline
    Kiểm tra dữ liệu & \code{src/lib/validations/} & Lược đồ \code{zod} dùng chung cho biểu mẫu và Server Action. \\
    \hline
    Động cơ phiên học & \code{src/app/(app)/learn/} & Bộ rút gọn \code{session-machine.ts} và các thành phần câu hỏi. \\
    \hline
  \end{tabular}
  \caption{Ánh xạ các phần frontend tới tệp hiện thực then chốt}
  \label{tab:frontend-layout}
\end{table}

\subsection{Cấu trúc App Router và phân nhóm tuyến}
\label{subsec:fe-approuter}

Cây tuyến được tổ chức bằng \textit{nhóm tuyến} (route group --- thư mục đặt
trong ngoặc đơn, không xuất hiện trong URL) để mỗi khu vực có bố cục và lớp bảo
vệ riêng: \code{(auth)} cho đăng nhập/đăng ký/xác minh email, \code{(app)} cho
trải nghiệm học của người dùng (trang chủ, học, khám phá, luyện tập, cộng đồng,
hồ sơ) và \code{(admin)} cho khu vực quản trị từ vựng. Mỗi phân đoạn tận dụng
các tệp đặc biệt của App Router --- \code{layout.tsx} (bố cục dùng chung),
\code{loading.tsx} (trạng thái chờ qua \textit{Suspense}), \code{error.tsx}
(ranh giới lỗi cấp phân đoạn) và \code{not-found.tsx} --- nên một lỗi cục bộ
không làm trắng toàn bộ ứng dụng. Các thành phần chỉ phục vụ một tuyến được đặt
ngay trong thư mục tuyến đó, chỉ nâng lên \code{src/components/} khi dùng lại ở
nhiều nơi.

\subsection{Tầng truy cập API phía máy chủ}
\label{subsec:fe-api}

Mọi lời gọi tới backend đi qua một client HTTP \textit{chỉ chạy phía máy chủ}
trong \code{src/lib/api.ts}. Địa chỉ backend (\code{API_BASE_URL}) được đọc từ
biến môi trường phía máy chủ, cố tình không gắn tiền tố \code{NEXT_PUBLIC_}
để không rò rỉ ra phía máy khách. Module phơi bày hai điểm vào:

\begin{itemize}
  \item \code{apiRequest} --- gọi \textit{không cần xác thực} (đăng nhập, đăng
        ký, làm mới token); tự đặt \code{Content-Type} JSON hoặc để runtime tự
        đặt biên \textit{multipart} khi thân là \code{FormData} (tải tệp lên).
  \item \code{authedRequest} --- đính kèm access token và, khi gặp
        \code{401}, \textit{tự động} làm mới refresh token đúng một lần rồi thử
        lại; chỉ gọi từ Server Action hoặc Route Handler vì nó có thể ghi lại
        cookie token đã xoay.
\end{itemize}

Cả hai chuẩn hoá thân lỗi của NestJS thành một kiểu \code{ApiResult} thống nhất
(cờ \code{ok}, mã trạng thái, dữ liệu, thân lỗi), nhờ đó tầng gọi xử lý lỗi
theo \textit{giá trị trả về} thay vì ngoại lệ --- khớp với triết lý phân biệt
lỗi \textit{dự kiến} và \textit{không dự kiến} của App Router.

\subsection{Quản lý phiên và xác thực phía máy khách}
\label{subsec:fe-session}

Token phiên được lưu trong \textbf{cookie httpOnly} với vòng đời bám đúng hợp
đồng backend (mục~\ref{sec:hien-thuc-auth}): access token 15~phút, refresh
token 30~ngày và được xoay sau mỗi lần gọi \code{/v1/auth/refresh}. Vì các
helper phiên import \code{next/headers}, chúng \textit{chỉ biên dịch được phía
máy chủ} --- một lớp bảo vệ ở mức biên dịch ngăn token lọt vào Client Component.
Khi \code{authedRequest} gặp \code{401}, nó âm thầm đổi refresh token lấy cặp
token mới, ghi lại cookie và thử lại đúng một lần; nếu vẫn thất bại, phiên bị
xoá và người dùng được chuyển về trang đăng nhập. Sau khi xác thực thành công,
Server Action điều hướng người dùng theo vai trò và trạng thái khai báo hồ sơ:
quản trị viên tới khu vực \code{/admin}, người học đã khai báo tới
\code{/dashboard}, người học chưa khai báo tới luồng \code{/onboarding}.

\subsection{Động cơ phiên học phía máy khách}
\label{subsec:fe-learn}

Phần tương tác nhất của ứng dụng là vòng lặp phiên học. Hàng đợi câu hỏi được
điều phối bằng một \textbf{bộ rút gọn (reducer) thuần} trong
\code{session-machine.ts}, tách hẳn khỏi React và mạng nên kiểm thử đơn vị được
độc lập. Nó nắm danh sách câu hỏi còn lại (câu đang hiển thị luôn là
\code{queue[0]}), các bộ đếm đúng/đã trả lời và trạng thái kết thúc
(\code{active}, \code{empty}, \code{done}, \code{expired}, \code{error}). Khi
một thẻ cần lặp lại trong cùng phiên, nó được \textit{nối lại vào cuối} hàng
đợi thay vì hẹn giờ theo đồng hồ thực. Bộ rút gọn này được ghép với hai Server
Action --- \code{startSessionAction} và \code{submitAnswerAction} --- gọi tới
\code{/v1/me/learn/session} và \code{/v1/me/learn/answer}. Vì máy chủ \textit{ký
HMAC} từng câu hỏi (mục~\ref{subsec:impl-hmac}), một \code{401} còn lại sau khi
đã làm mới token được hiểu là chữ ký câu hỏi đã hết hạn hoặc bị giả mạo, và giao
diện chuyển sang trạng thái \code{expired} mời người học bắt đầu phiên mới. Mỗi
\textit{chế độ} câu hỏi (thẻ ghi nhớ, điền khuyết trắc nghiệm/gõ chữ, nghe chép
chính tả, chọn theo hình/âm thanh, luyện phát âm\ldots) là một thành phần riêng
dưới \code{learn/questions/}, nhận dữ liệu đã ký từ máy chủ và chỉ chịu trách
nhiệm thu thập câu trả lời để nộp lại.

% =====================================================================
\section{Kiểm thử hệ thống}
\label{sec:kiem-thu}

Việc kiểm thử được hỗ trợ trực tiếp bởi một lựa chọn thiết kế lặp lại xuyên
suốt chương: \textbf{tách logic thông minh thành các hàm thuần, độc lập
framework}. Các hàm như \code{applySm2} (SM-2), bộ ký HMAC, bộ chấm đáp án,
bộ dựng prompt và bộ phân tích phản hồi LLM đều không phụ thuộc NestJS hay cơ
sở dữ liệu, nên có thể kiểm thử đơn vị nhanh và tất định. Chiến lược kiểm thử
gồm hai mức (Bảng~\ref{tab:test-targets}):

\begin{itemize}
  \item \textbf{Kiểm thử đơn vị} bằng \textbf{Jest}~\cite{jest2024}: tập trung
        vào các thuật toán và bộ phân tích --- nơi lỗi tinh vi dễ lọt và hậu
        quả lớn (lập lịch sai, lưu dữ liệu AI rác, chấm điểm sai).
  \item \textbf{Kiểm thử tích hợp (end-to-end)} bằng \textbf{Supertest}: gọi
        thật vào các điểm cuối API để kiểm tra luồng xác thực, kiểm tra đầu
        vào và mã trạng thái trả về.
\end{itemize}

\begin{table}[h]
  \centering
  \small
  \renewcommand{\arraystretch}{1.3}
  \begin{tabular}{|p{5.0cm}|p{9.4cm}|}
    \hline
    \textbf{Thành phần} & \textbf{Trọng tâm kiểm thử} \\
    \hline
    \code{srs.ts} & Chuyển trạng thái SM-2, bước học, tốt nghiệp, rớt bước, ngưỡng thành thạo. \\
    \hline
    \code{hmac-signer.service.ts} & Ký/xác minh, chống giả mạo, hết hạn, so sánh thời gian hằng số. \\
    \hline
    \code{answer-grader.service.ts} & Quy đổi điểm chất lượng cho từng loại câu hỏi. \\
    \hline
    \code{question-builder.service.ts} & Dựng bài học đúng dải khó, tôn trọng các trần và lọc dữ liệu. \\
    \hline
    \code{gemma-judge.ts}, \code{gemma-enricher.ts} & Phân tích phòng thủ JSON, kẹp số, từ chối CEFR sai. \\
    \hline
    \code{import/tokenize.ts}, \code{normalize.ts}, \code{xlsx-parser.ts} & Tách lemma, chuẩn hoá, đọc bảng tính cho nhập hàng loạt. \\
    \hline
  \end{tabular}
  \caption{Một số thành phần tiêu biểu và trọng tâm kiểm thử}
  \label{tab:test-targets}
\end{table}

Chất lượng mã nguồn còn được kiểm soát ở mức tĩnh bằng trình kiểm tra kiểu của
TypeScript (chế độ \textit{strict}), \textbf{ESLint} và \textbf{Prettier},
chạy trước mỗi lần commit (đáp ứng NFR-06).

% =====================================================================
\section{Kết chương}
\label{sec:ket-chuong-4}

Chương này đã trình bày quá trình hiện thực hoá hệ thống, với trọng tâm là các
thành phần thông minh và tự động hoá. Ở phía \textit{động cơ học}, hệ thống
thích nghi với từng người học qua thuật toán SM-2 mở rộng bước học
(\code{srs.ts}), cơ chế chọn từ theo khoảng cách CEFR, bộ sinh câu hỏi nhiều
tầng theo giai đoạn ghi nhớ với phương án nhiễu "dễ nhầm", và cơ chế chấm điểm
phi trạng thái an toàn bằng chữ ký HMAC. Ở phía \textit{nội dung}, các tính
năng AI --- làm giàu từ điển bằng LLM kết hợp từ điển, nhập hàng loạt từ nhiều
định dạng, sinh âm thanh và hình ảnh tự động, chấm điểm câu và phát âm --- đều
chạy trên một hạ tầng xử lý nền bất đồng bộ có tính idempotent, biết thử lại,
tự giới hạn tần suất và suy giảm có kiểm soát. Ở phía \textit{máy khách}, một
ứng dụng Next.js theo App Router tiêu thụ toàn bộ các API trên qua một tầng
truy cập phía máy chủ, quản lý phiên bằng cookie httpOnly với cơ chế xoay token
tự động, và điều phối phiên học bằng một bộ rút gọn thuần tách khỏi React. Tất
cả được củng cố bằng các chỉ số tiến độ tự động và một chiến lược kiểm thử dựa
trên việc tách logic thành các hàm thuần. Những kết quả hiện thực hoá này là cơ
sở để chương tiếp theo đánh giá hệ thống và rút ra kết luận.

%  vào main.tex (sau Chương 3).
%
%  GHI CHÚ VỀ HÌNH VẼ: dùng macro \figph{<tên-file>}{<tỉ-lệ-rộng>}{<chú-thích>}{<nhãn>}
%  (đã khai báo ở các chương trước; khai báo lại bằng \providecommand để an
%  toàn nếu biên dịch độc lập). Mã nguồn (code) trình bày bằng môi trường
%  verbatim như các chương trước.
%
%  KHÔNG hardcode số chương/mục/bảng/hình/công thức — luôn dùng \label + \ref / \cite.
% =====================================================================

% --- Macro tiện ích (idempotent — không định nghĩa lại nếu đã có) ------
\providecommand{\code}[1]{\texttt{\detokenize{#1}}} % in định danh snake_case an toàn
\providecommand{\figph}[4]{%
  \begin{figure}[htbp]
    \centering
    \IfFileExists{images/#1}%
      {\includegraphics[width=#2\textwidth,height=0.82\textheight,keepaspectratio]{#1}}%
      {\setlength{\fboxsep}{10pt}\fbox{\begin{minipage}[c][3.6cm][c]{#2\textwidth}\centering\itshape Sơ đồ minh hoạ --- chèn tệp ảnh \code{images/#1} vào thư mục \code{images/}.\end{minipage}}}%
    \caption{#3}%
    \label{#4}%
  \end{figure}}

\chapter{HIỆN THỰC HOÁ HỆ THỐNG VÀ CÁC TÍNH NĂNG THÔNG MINH}
\label{chuong:hien-thuc-hoa}

Hai chương trước đã trình bày kết quả phân tích, thiết kế
(Chương~\ref{chuong:phan-tich-thiet-ke}) và cơ sở lý thuyết, công nghệ
(Chương~\ref{chuong:co-so-ly-thuyet}). Chương này trình bày cách các thiết
kế đó được \textit{hiện thực hoá} thành mã nguồn chạy được, với trọng tâm là
các \textbf{thành phần thông minh và tự động hoá} làm nên giá trị cốt lõi của
hệ thống: động cơ học theo ôn tập ngắt quãng, cơ chế sinh câu hỏi và chấm
điểm thích nghi, các tính năng làm giàu nội dung và chấm điểm dựa trên trí
tuệ nhân tạo, cùng hạ tầng xử lý nền bất đồng bộ liên kết chúng lại.

Mục~\ref{sec:tong-quan-hien-thuc} mô tả tổ chức mã nguồn và các mối quan tâm
xuyên suốt. Mục~\ref{sec:hien-thuc-du-lieu} và
mục~\ref{sec:hien-thuc-auth} trình bày ngắn gọn hiện thực hoá tầng dữ liệu và
phân hệ xác thực. Trọng tâm của chương nằm ở
mục~\ref{sec:dong-co-hoc} (động cơ học thông minh) và
mục~\ref{sec:ai-tu-dong} (các tính năng AI và tự động hoá nội dung). Cuối
cùng, mục~\ref{sec:tien-do} trình bày phân hệ theo dõi tiến độ tự động,
mục~\ref{sec:hien-thuc-frontend} trình bày hiện thực hoá tầng giao diện người
dùng (frontend), mục~\ref{sec:kiem-thu} trình bày chiến lược kiểm thử và
mục~\ref{sec:ket-chuong-4} tóm tắt chương.

% =====================================================================
\section{Tổng quan quá trình hiện thực hoá}
\label{sec:tong-quan-hien-thuc}

Hệ thống được hiện thực hoá đúng theo ngăn xếp công nghệ đã chọn ở
Chương~\ref{chuong:co-so-ly-thuyet}: backend viết bằng \textbf{NestJS~11}
trên \textbf{TypeScript} (chế độ \textit{strict}), truy cập
\textbf{PostgreSQL} qua \textbf{TypeORM}, và dùng \textbf{Redis/BullMQ} cho
xử lý nền. Toàn bộ mã nguồn được biên dịch bằng trình biên dịch của NestJS và
khởi chạy dưới \textbf{hai tiến trình} chia sẻ chung cơ sở dữ liệu và hàng
đợi, đúng mô hình nhà sản xuất--người tiêu thụ đã nêu ở
mục~\ref{sec:kien-truc} của Chương~\ref{chuong:phan-tich-thiet-ke}:

\begin{itemize}
  \item \textbf{Tiến trình API} (\code{npm run start}): phục vụ các yêu cầu
        HTTP đồng bộ; với tác vụ nặng nó chỉ tạo bản ghi công việc rồi đẩy
        vào hàng đợi và trả về ngay.
  \item \textbf{Tiến trình Worker} (\code{npm run start:worker}): tiêu thụ
        hàng đợi để thực thi các tác vụ nền (làm giàu từ vựng, sinh âm thanh,
        chấm điểm luyện tập).
\end{itemize}

Mọi tham số vận hành (chuỗi kết nối CSDL, khoá bí mật, khoá API của dịch vụ
AI, hạn mức tần suất\ldots) đều được nạp từ \textbf{biến môi trường} qua
\code{@nestjs/config}, đáp ứng yêu cầu tính khả chuyển NFR-08 trong
Bảng~\ref{tab:nfr}: cùng một ảnh mã nguồn chạy được ở nhiều môi trường mà
không cần sửa code.

\subsection{Tổ chức mã nguồn theo mô-đun}
\label{subsec:to-chuc-ma-nguon}

Mỗi mô-đun nghiệp vụ ở Bảng~\ref{tab:modules} (Chương~\ref{chuong:phan-tich-thiet-ke})
được hiện thực thành một thư mục dưới \code{src/} tự đóng gói controller,
service, entity và DTO của mình. Bảng~\ref{tab:source-layout} ánh xạ các
mô-đun trọng yếu tới những tệp hiện thực then chốt --- là cơ sở để các mục
sau tham chiếu trực tiếp.

\begin{table}[h]
  \centering
  \small
  \renewcommand{\arraystretch}{1.3}
  \begin{tabular}{|p{2.6cm}|p{4.4cm}|p{6.6cm}|}
    \hline
    \textbf{Mô-đun} & \textbf{Thư mục mã nguồn} & \textbf{Tệp hiện thực then chốt} \\
    \hline
    \code{auth} & \code{src/auth} & \code{token.service.ts} (cấp/xoay token), \code{jwt.strategy.ts}, các \code{guards/}, \code{services/} (Google/Apple/GitHub). \\
    \hline
    \code{learn} & \code{src/learn} & \code{learn.service.ts}, \code{vocab-picker.service.ts}, \code{question-builder.service.ts}, \code{distractor.service.ts}, \code{answer-grader.service.ts}, \code{hmac-signer.service.ts}, \code{question-bands.ts}. \\
    \hline
    \code{progress} & \code{src/progress} & \code{srs.ts} (thuật toán SM-2), \code{progress.service.ts} (lập lịch, thống kê, chuỗi ngày học). \\
    \hline
    \code{vocabularies} & \code{src/vocabularies} & \code{enrichment/} (làm giàu bằng LLM, nhập hàng loạt), \code{audio/}, \code{images/}, \code{vocabulary-persistence.service.ts}. \\
    \hline
    \code{practice} & \code{src/practice} & \code{gemma-judge.ts} (chấm câu bằng LLM), \code{scoring.processor.ts}, \code{practice.service.ts}. \\
    \hline
    \code{pronunciation} & \code{src/pronunciation} & \code{pronunciation-scoring.client.ts} (proxy tới vi dịch vụ), \code{pronunciation.service.ts}. \\
    \hline
    \code{leaderboard} & \code{src/leaderboard} & \code{leaderboard.service.ts}. \\
    \hline
    \code{config} & \code{src/config} & Các tệp cấu hình theo miền: \code{learn.config.ts}, \code{gemma.config.ts}, \code{auth.config.ts}\ldots \\
    \hline
  \end{tabular}
  \caption{Ánh xạ mô-đun nghiệp vụ tới tệp hiện thực then chốt}
  \label{tab:source-layout}
\end{table}

\subsection{Các mối quan tâm xuyên suốt}
\label{subsec:cross-cutting}

Nhờ cơ chế \textit{tiêm phụ thuộc} và các thành phần khai báo của NestJS, các
mối quan tâm chung được hiện thực một lần và áp dụng toàn cục thay vì lặp lại
trong từng controller:

\begin{itemize}
  \item \textbf{Kiểm tra đầu vào:} một \code{ValidationPipe} toàn cục với
        \code{whitelist}, \code{forbidNonWhitelisted} và \code{transform}
        biến mỗi DTO (khai báo bằng \code{class-validator}) thành lớp bảo vệ:
        trường lạ bị loại bỏ, dữ liệu sai kiểu bị chặn bằng lỗi 400 trước khi
        chạm tới service.
  \item \textbf{Xác thực và phân quyền:} \code{JwtAuthGuard} (dựa trên
        \code{passport-jwt}, xem \code{jwt.strategy.ts}) trích và xác minh
        access token ở mọi endpoint cần đăng nhập; \code{RolesGuard} kết hợp
        decorator \code{@Roles('admin')} chặn các thao tác quản trị; quyền sở
        hữu tài nguyên được kiểm tra ngay trong service (mục~\ref{subsec:rbac}).
  \item \textbf{Phiên bản hoá API:} bật \textit{URI versioning} mặc định
        \code{v1}; controller khai báo \code{@Controller(\{ path, version: '1' \})}
        nên mọi tài nguyên nằm dưới tiền tố \code{/v1} (mục~\ref{subsec:rest-theory}).
  \item \textbf{Cấu hình theo miền:} mỗi nhóm tham số được gói trong một
        \textit{namespace} cấu hình riêng (ví dụ \code{learn.config.ts},
        \code{gemma.config.ts}), tách bạch và có giá trị mặc định an toàn cho
        môi trường phát triển.
\end{itemize}

% =====================================================================
\section{Hiện thực hoá tầng dữ liệu}
\label{sec:hien-thuc-du-lieu}

Tầng dữ liệu hiện thực 17 bảng trong Bảng~\ref{tab:db-overview} thành các lớp
\textit{thực thể} (entity) TypeORM. Mỗi thực thể khai báo khoá chính
\code{uuid}, tên cột \code{snake_case} và kiểu \code{timestamptz} cho các mốc
thời gian, đúng quy ước thiết kế ở Chương~\ref{chuong:phan-tich-thiet-ke}.
Một số kỹ thuật hiện thực đáng chú ý:

\begin{itemize}
  \item \textbf{Giao dịch bảo đảm tính nguyên tử:} các thao tác ghi nhiều
        bảng được gói trong \code{dataSource.transaction(...)}. Ví dụ tiêu
        biểu là việc nộp một lượt ôn (mục~\ref{subsec:impl-hmac}): bản ghi
        tiến độ và bản ghi nhật ký hoạt động được lưu trong cùng một giao
        dịch, nên một lượt ôn không thể "ghi nửa vời" --- đáp ứng NFR-04 trong
        Bảng~\ref{tab:nfr}.
  \item \textbf{Ràng buộc toàn vẹn và chống trùng lặp:} ngoài khoá ngoại và
        quy tắc xoá lan truyền, hệ thống dùng \textit{chỉ mục duy nhất bộ
        phận} (partial unique index) trên bộ ba \code{(language, lemma,
        part_of_speech)} áp dụng riêng cho từ thuộc kho hệ thống
        (\code{WHERE source = 'system'}); nhờ đó kho hệ thống không bao giờ có
        hai bản ghi trùng từ loại cho cùng một từ, trong khi từ cá nhân của
        người dùng vẫn không bị ràng buộc này.
  \item \textbf{Quản lý lược đồ bằng migration:} mọi thay đổi lược đồ được mô
        tả bằng các \textit{migration} có thể phiên bản hoá, được nhận diện và
        chạy qua công cụ dòng lệnh của TypeORM. Cách này cho phép tái lập lược
        đồ một cách xác định giữa các môi trường (đáp ứng NFR-06).
  \item \textbf{Phi chuẩn hoá có kiểm soát:} một số trường được phi chuẩn hoá
        để tăng tốc đọc, ví dụ \code{decks.vocab_count} (số từ trong bộ thẻ)
        được cập nhật khi thêm/bớt thành viên thay vì đếm mỗi lần truy vấn.
\end{itemize}

% =====================================================================
\section{Hiện thực hoá xác thực và phân quyền}
\label{sec:hien-thuc-auth}

Phân hệ \code{auth} hiện thực các kỹ thuật bảo mật đã trình bày về lý thuyết
ở mục~\ref{sec:xac-thuc-phan-quyen}. Lớp \code{TokenService} đảm nhiệm vòng
đời phiên đăng nhập với một số quyết định hiện thực đáng chú ý:

\begin{itemize}
  \item \textbf{Cặp token bất đối xứng về bản chất:} \textit{access token} là
        một \textbf{JWT} được ký, tự chứa và sống ngắn (mục~\ref{subsec:jwt}),
        trong khi \textit{refresh token} là một chuỗi ngẫu nhiên 48 byte
        \textit{mờ} (opaque) --- không mang thông tin, chỉ lưu \textbf{bản băm
        SHA-256} trong bảng \code{refresh_tokens}. Máy chủ không bao giờ lưu
        token gốc, nên dù cơ sở dữ liệu bị lộ thì token cũng không thể tái sử
        dụng.
  \item \textbf{Xoay vòng và phát hiện tái dùng:} mỗi lần làm mới phiên, token
        cũ bị đánh dấu thu hồi và một token mới được cấp. Nếu một token đã thu
        hồi (hoặc không tồn tại) được trình lại --- dấu hiệu rò rỉ --- hệ thống
        \textbf{thu hồi toàn bộ} token còn hiệu lực của người dùng đó và trả
        về lỗi 401, hạn chế thiệt hại khi token bị đánh cắp.
  \item \textbf{Băm mật khẩu và đăng nhập bên thứ ba:} mật khẩu được băm bằng
        \textbf{bcrypt} (mục~\ref{subsec:bcrypt}); với OAuth/OIDC
        (mục~\ref{subsec:oauth}), token danh tính do nhà cung cấp cấp được xác
        minh rồi liên kết tới một bản ghi trong \code{user_identities}, cho
        phép một tài khoản đăng nhập bằng nhiều phương thức.
  \item \textbf{Chống vét cạn:} endpoint đăng nhập được giới hạn tần suất bằng
        \code{@nestjs/throttler}, trả về lỗi 429 khi vượt ngưỡng.
\end{itemize}

% =====================================================================
\section{Động cơ học thông minh theo ôn tập ngắt quãng}
\label{sec:dong-co-hoc}

Đây là phân hệ \textit{thông minh} cốt lõi của hệ thống. Một phiên học không
phải là một danh sách câu hỏi tĩnh: với mỗi người học và mỗi từ, hệ thống
\textbf{thích nghi} ở ba cấp độ --- chọn \textit{từ nào} để học, sinh
\textit{dạng câu hỏi nào} phù hợp giai đoạn ghi nhớ, và quyết định
\textit{khi nào} ôn lại. Toàn bộ quy trình tuân theo luồng đã mô hình hoá ở
Hình~\ref{fig:activity-srs}: chọn từ $\rightarrow$ mở rộng thành chuỗi câu hỏi
đã ký $\rightarrow$ chấm điểm $\rightarrow$ (ở bước cuối) lập lịch lại.

\subsection{Lập lịch ôn tập SM-2 mở rộng bước học}
\label{subsec:impl-sm2}

Hạt nhân lập lịch là hàm thuần \code{applySm2} trong \code{src/progress/srs.ts}:
nó nhận trạng thái hiện tại của một cặp (người dùng, từ), một
\textit{điểm chất lượng} \( q \in \{0,\dots,5\} \) và danh sách bước học, rồi
trả về trạng thái mới cùng thời điểm ôn kế tiếp. Hàm hiện thực thuật toán
\textbf{SM-2} (mục~\ref{subsec:sm2-theory}) được \textit{mở rộng bằng các bước
học} theo phong cách Anki. Hệ số dễ nhớ \( EF \) được cập nhật ở \textit{mọi}
lượt ôn theo công thức~\eqref{eq:sm2-ef}, còn khoảng cách ngày tuân theo
công thức~\eqref{eq:sm2-interval}; phần mở rộng bổ sung một thang
\textit{bước học} tính bằng phút để củng cố từ mới ngay trong phiên trước khi
giao cho thang ngày:

\begin{verbatim}
// src/progress/srs.ts — cập nhật hệ số dễ nhớ (chặn dưới 1.3)
function updateEaseFactor(easeFactor, quality) {
  return Math.max(
    1.3,
    easeFactor + (0.1 - (5 - quality) * (0.08 + (5 - quality) * 0.02)),
  );
}
\end{verbatim}

Logic chuyển trạng thái được tóm tắt như sau (Hình~\ref{fig:srs-states}):

\begin{itemize}
  \item Thẻ đang ở \textit{trạng thái bước} (\code{learningStepIndex} khác
        \code{null}) dùng khoảng cách \textbf{phút} lấy từ danh sách bước (mặc
        định \([1, 10]\) phút). Trả lời đúng (\( q \ge 3 \)) đẩy lên một bước;
        qua bước cuối, thẻ \textit{tốt nghiệp} sang thang ngày. Trả lời sai
        (\( q < 3 \)) đặt lại về bước 0.
  \item Thẻ đã tốt nghiệp dùng thang ngày: \( I = 1 \) ngày (lần đầu), \( 6 \)
        ngày (lần hai), sau đó \( I \leftarrow \lceil I \times EF \rceil \).
        Một lần trả lời sai sẽ đưa thẻ \textit{rớt} về bước 0 (ôn lại trong
        phút) thay vì chờ tới hôm sau.
  \item Trạng thái hiển thị được suy ra từ số lần lặp và khoảng cách:
        \( I \ge 90 \) ngày $\rightarrow$ \code{mastered}; số lần ôn đúng liên
        tiếp \( \ge 3 \) $\rightarrow$ \code{review}; còn lại \code{learning}.
\end{itemize}

Bảng~\ref{tab:srs-params} liệt kê các tham số then chốt và giá trị mặc định
của chúng (đều có thể cấu hình qua biến môi trường trong \code{learn.config.ts}).

\begin{table}[h]
  \centering
  \small
  \renewcommand{\arraystretch}{1.3}
  \begin{tabular}{|p{5.2cm}|p{3.0cm}|p{6.0cm}|}
    \hline
    \textbf{Tham số} & \textbf{Giá trị mặc định} & \textbf{Ý nghĩa} \\
    \hline
    Hệ số dễ nhớ khởi tạo \( EF_0 \) & 2.5 & Giá trị ban đầu của \( EF \). \\
    \hline
    Chặn dưới \( EF_{\min} \) & 1.3 & Ngăn khoảng ôn co lại quá nhanh. \\
    \hline
    Bước học (phút) & \([1, 10]\) & Khoảng cách trong phiên cho thẻ mới/đang học/vừa rớt. \\
    \hline
    Khoảng ngày sau tốt nghiệp & 1 ngày, 6 ngày & Hai mốc đầu của thang ngày. \\
    \hline
    Ngưỡng thành thạo & \( \ge 90 \) ngày & Khoảng ôn đạt mốc này $\rightarrow$ \code{mastered}. \\
    \hline
    Cửa sổ đưa lại trong phiên & 15 phút & Nếu đến hạn trong cửa sổ này, câu hỏi kế được trả kèm ngay. \\
    \hline
  \end{tabular}
  \caption{Tham số của thuật toán lập lịch SM-2 mở rộng}
  \label{tab:srs-params}
\end{table}

% --- PlantUML gợi ý cho Hình máy trạng thái SRS ----------------------
% @startuml
% [*] --> new
% new --> learning : ghi danh (bước 0)
% learning --> learning : sai (q<3) → bước 0\nhoặc đúng nhưng chưa tốt nghiệp
% learning --> review : đúng liên tiếp ≥ 3 lần
% review --> review : đúng → I ← ⌈I×EF⌉
% review --> learning : sai → rớt về bước 0
% review --> mastered : I ≥ 90 ngày
% mastered --> learning : sai → rớt về bước 0
% mastered --> [*]
% @enduml
\figph{srs-state-machine.png}{0.8}{Máy trạng thái vòng đời thẻ học theo SM-2 mở rộng}{fig:srs-states}

Một quyết định thiết kế quan trọng: \code{applySm2} là một \textbf{hàm thuần}
(pure function) --- không phụ thuộc NestJS, cơ sở dữ liệu hay thời gian thực
(mốc \code{now} được truyền vào) --- nên có thể kiểm thử đơn vị triệt để
(mục~\ref{sec:kiem-thu}).

\subsection{Chọn từ học thích nghi theo trình độ}
\label{subsec:impl-picker}

Lớp \code{VocabPickerService} chọn tập từ cho mỗi phiên theo bốn chế độ
(hằng ngày, theo chủ đề, theo bộ thẻ, ôn tập). Nguyên tắc chung là
\textit{ưu tiên từ đã đến hạn} (sắp theo \code{next_review_at} tăng dần),
sau đó bù thêm \textit{từ mới} cho đủ hạn mức phiên. Điểm "thông minh" nằm ở
cách chọn từ mới: thay vì lọc cứng theo trình độ (dễ dẫn tới phiên rỗng khi
người học đã học hết từ ở đúng bậc), hệ thống coi trình độ là một
\textit{ưu tiên mềm} --- sắp các từ chưa học theo \textbf{khoảng cách CEFR}
tới trình độ người học (gần nhất trước), rồi mới tới thứ hạng tần suất:

\begin{verbatim}
// src/learn/vocab-picker.service.ts — ưu tiên từ gần trình độ người học
qb.orderBy(`ABS(${cefrOrdinalCase('v.cefr_level')} - ${userIdx})`,
           'ASC', 'NULLS LAST')
  .addOrderBy('v.frequency_rank', 'ASC', 'NULLS LAST')
  .limit(limit);
\end{verbatim}

Nhờ đó người học đã cạn từ ở bậc của mình (hoặc các từ còn lại không gắn bậc
CEFR) vẫn nhận được các từ phù hợp nhất hiện có thay vì một phiên rỗng. Khi
thực sự không còn từ nào, picker trả về một \textit{lý do rỗng} có ngữ nghĩa
(\code{no_due_cards}, \code{no_more_at_level}, \code{deck_exhausted},
\code{no_enrollment}) để máy khách hiển thị thông điệp phù hợp; chỉ với
\code{no_due_cards} hệ thống mới kèm theo \textit{thời điểm đến hạn kế tiếp}.
Các từ mới được chọn sẽ được \textit{tự động ghi danh} (auto-enroll) vào hàng
đợi học của người dùng trước khi sinh câu hỏi.

\subsection{Sinh câu hỏi thích nghi theo giai đoạn ghi nhớ}
\label{subsec:impl-questions}

Mỗi từ trong phiên được \code{QuestionBuilderService} mở rộng thành một
\textit{bài học} (lesson) gồm một chuỗi câu hỏi xếp từ dễ đến khó. Mức độ khó
được tổ chức thành ba \textbf{dải khó} (difficulty band) và \textbf{12 dạng
câu hỏi} (Bảng~\ref{tab:question-ladder}). Điểm thích nghi cốt lõi:
\textit{giai đoạn ghi nhớ của từ quyết định những dải nào được dùng}. Từ càng
thành thạo thì "sàn" độ khó càng được nâng lên, loại dần các dạng nhận biết
dễ và tập trung vào sản sinh ngôn ngữ:

\begin{table}[h]
  \centering
  \small
  \renewcommand{\arraystretch}{1.3}
  \begin{tabular}{|p{2.6cm}|p{3.4cm}|p{7.6cm}|}
    \hline
    \textbf{Trạng thái từ} & \textbf{Dải khó tối thiểu} & \textbf{Nhóm câu hỏi đại diện} \\
    \hline
    \code{new} & Nhận biết (NEW) & Thẻ học (flashcard), điền khuyết trắc nghiệm, chọn từ/nghĩa, nghe chọn, nhìn hình chọn từ. \\
    \hline
    \code{learning} / \code{review} & Gợi nhớ (REVIEW) & Điền khuyết tự gõ, nghe gõ lại (dictation), luyện phát âm --- người học phải tự tạo ra dạng từ. \\
    \hline
    \code{mastered} & Sản sinh (MASTER) & Phân biệt nghĩa (sense disambiguation) --- dạng khó nhất, đòi hỏi phân biệt giữa các nghĩa của từ đa nghĩa. \\
    \hline
  \end{tabular}
  \caption{Ánh xạ giai đoạn ghi nhớ của từ tới dải khó và nhóm câu hỏi}
  \label{tab:question-ladder}
\end{table}

Trên nền tảng đó, bộ dựng câu hỏi áp dụng thêm các cơ chế giữ cho bài học
\textit{ngắn, đa dạng và khả thi}:

\begin{itemize}
  \item \textbf{Lấy mẫu xác định theo từ:} với mỗi dải, số dạng câu hỏi
        "tuỳ chọn" được giới hạn (mặc định 2 dạng/dải) và được chọn bằng phép
        \textit{xáo trộn xác định} theo định danh của từ. Nhờ vậy mỗi từ luyện
        một tập dạng câu hỏi khác nhau (thay vì mọi bài đều chạy hết các
        dạng), nhưng kết quả vẫn \textit{tái lập được}. Dạng thẻ học luôn được
        giữ làm bước nhập môn.
  \item \textbf{Trần họ điền khuyết:} các dạng cùng "khoét" một từ trong câu
        bị giới hạn số lần (mặc định 2) để một câu ví dụ không bị khoét nhiều
        bước liên tiếp.
  \item \textbf{Lọc theo dữ liệu sẵn có:} mỗi dạng chỉ được dựng khi đủ dữ
        liệu (ví dụ dạng nghe cần có \code{audio_url}, dạng nhìn hình cần ảnh
        nghĩa, dạng dịch cần bản dịch ở ngôn ngữ đích). Dạng nào thiếu dữ liệu
        sẽ bị bỏ qua một cách an toàn.
\end{itemize}

\paragraph{Sinh phương án nhiễu thông minh.} Chất lượng của một câu trắc
nghiệm phụ thuộc vào \textit{phương án nhiễu} (distractor). \code{DistractorService}
không lấy nhiễu ngẫu nhiên mà chọn các từ "dễ nhầm": \textbf{cùng từ loại,
cùng ngôn ngữ, cùng bậc CEFR \( \pm 1 \)}, và \textit{ưu tiên các từ chia sẻ
ít nhất một chủ đề} với từ mục tiêu; chỉ khi không đủ mới nới lỏng điều kiện.
Với dạng phân biệt nghĩa, nhiễu còn được lấy chính từ \textit{các nghĩa khác}
của cùng từ (bẫy đa nghĩa), buộc người học thực sự phân biệt ngữ cảnh.

\subsection{Chấm điểm phi trạng thái có ký HMAC}
\label{subsec:impl-hmac}

Như đã nêu ở mục~\ref{subsec:hmac}, hệ thống chấm điểm phiên học theo cách
\textbf{phi trạng thái}: máy chủ không lưu phiên trong bộ nhớ mà "đóng dấu"
tính toàn vẹn vào chính mỗi câu hỏi. Khi sinh câu hỏi, \code{HmacSignerService}
ký một thông điệp gồm các trường định danh và --- quan trọng --- cả
\code{stepIndex}/\code{stepCount}, để máy khách không thể giả mạo việc câu
nào là câu cuối (câu duy nhất làm thay đổi lịch ôn):

\begin{verbatim}
// src/learn/hmac-signer.service.ts — thông điệp được ký
const message = [userId, vocabularyId, type, exampleId,
                 translationLang ?? '', String(stepIndex),
                 String(stepCount), nonce, String(issuedAtMs)].join('|');
return createHmac('sha256', this.secret).update(message).digest('hex');
\end{verbatim}

Khi người học nộp đáp án (Hình~\ref{fig:seq-answer}), \code{LearnService} thực
hiện ba việc:

\begin{enumerate}
  \item \textbf{Xác minh chữ ký:} kiểm tra chữ ký HMAC bằng phép so sánh
        \textit{thời gian hằng số} (\code{timingSafeEqual}) và kiểm tra hạn
        dùng (mặc định 30 phút); chữ ký sai/hết hạn $\rightarrow$ lỗi 401.
  \item \textbf{Chấm điểm tại chỗ:} \code{AnswerGraderService} \textit{suy
        lại} đáp án đúng từ dữ liệu từ vựng (dạng từ bị khoét trong câu, bản
        dịch của nghĩa, hay chính lemma) rồi so khớp. Điểm chất lượng SM-2
        được quy đổi theo loại câu: trả lời nhanh và đúng được 5, đúng nhưng
        chậm được 4, sai lệch một ký tự (đo bằng khoảng cách Levenshtein) được
        coi là "gần đúng" (3), còn lại 2; riêng thẻ học cho người dùng tự đánh
        giá theo thang bốn mức.
  \item \textbf{Lập lịch ở bước cuối:} mỗi bài học chỉ là \textit{một} sự kiện
        SRS --- chỉ câu hỏi cuối (\code{stepIndex === stepCount - 1}) mới gọi
        \code{submitReview} để chạy \code{applySm2} và cập nhật lịch; các bước
        trước chỉ chấm để phản hồi tức thì. Nhờ vậy một từ không thể "tốt
        nghiệp" chỉ trong một lượt ngồi học.
\end{enumerate}

Việc cập nhật tiến độ và ghi nhật ký hoạt động được thực hiện trong cùng một
giao dịch (mục~\ref{sec:hien-thuc-du-lieu}). Nếu lịch vừa được dời vào trong
\textit{cửa sổ đưa lại} (15 phút), phản hồi sẽ kèm sẵn câu hỏi kế đã ký để máy
khách hiển thị lại từ ngay trong phiên mà không cần gọi lại API tạo phiên.

% =====================================================================
\section{Các tính năng trí tuệ nhân tạo và tự động hoá nội dung}
\label{sec:ai-tu-dong}

Nhóm tính năng thứ hai làm nên tính "thông minh" của hệ thống là việc
\textbf{tự động tạo và làm giàu nội dung} bằng các dịch vụ AI, giải quyết
trực tiếp rào cản "tạo nội dung thủ công tốn công" đã nêu ở
Chương~\ref{chuong:tong-quan}. Tất cả các tác vụ này đều nặng và phụ thuộc
dịch vụ ngoài, nên được hiện thực trên một \textit{hạ tầng xử lý nền} chung.

\subsection{Hạ tầng xử lý nền bất đồng bộ}
\label{subsec:impl-worker}

Mọi tác vụ AI được tổ chức theo mẫu \textbf{nhà sản xuất--người tiêu thụ}
trên \textbf{BullMQ}~\cite{bullmq2024}: tiến trình API tạo bản ghi công việc,
đẩy vào hàng đợi Redis rồi trả về mã \code{202} kèm mã công việc; tiến trình
worker tiêu thụ và xử lý nền (xem Hình~\ref{fig:activity-ai} và
Hình~\ref{fig:activity-practice}). Bảng~\ref{tab:async-jobs} liệt kê các hàng
đợi đã hiện thực. Bốn cơ chế bảo đảm độ tin cậy (NFR-04) được áp dụng nhất
quán cho mọi worker:

\begin{itemize}
  \item \textbf{Tính idempotent:} mỗi worker đọc lại bản ghi và \textit{chỉ}
        xử lý khi trạng thái còn là \code{pending}; một công việc bị chạy lại
        sẽ thoát sớm thay vì tạo dữ liệu trùng.
  \item \textbf{Thử lại có khoảng lùi:} lỗi tạm thời (mạng, 429, một phản hồi
        không phân tích được) khiến BullMQ thử lại với \textit{backoff}; chỉ
        khi cạn số lần thử, sự kiện \code{onFailed} mới ghi trạng thái
        \code{failed}.
  \item \textbf{Tự giới hạn tần suất:} worker đặt \textit{limiter} (mặc định
        15 lượt/phút) và độ song song bằng 1 để không vượt hạn mức của khoá
        AI dùng chung (tầng miễn phí).
  \item \textbf{Hạn mức theo ngày:} các điểm cuối người dùng có hạn mức theo
        ngày (ví dụ 30 lượt chấm câu/người/ngày) để bảo vệ tài nguyên chung,
        trả về lỗi 429 khi vượt.
\end{itemize}

\begin{table}[h]
  \centering
  \small
  \renewcommand{\arraystretch}{1.3}
  \begin{tabular}{|p{3.0cm}|p{4.6cm}|p{6.4cm}|}
    \hline
    \textbf{Hàng đợi} & \textbf{Kích hoạt bởi} & \textbf{Công việc của worker} \\
    \hline
    Làm giàu từ vựng & Tạo nhanh / nhập hàng loạt & Tra từ điển + gọi LLM sinh nghĩa, ví dụ, bản dịch, CEFR; tạo bản ghi từ vựng. \\
    \hline
    Sinh âm thanh & Sau khi tạo/duyệt từ & Lấy bản ghi âm thật hoặc tổng hợp giọng nói, tải lên kho media. \\
    \hline
    Sinh hình ảnh & Quy trình làm giàu nghĩa & Tìm ảnh minh hoạ theo từ, tải lên kho media. \\
    \hline
    Chấm điểm luyện viết & Nộp bài luyện viết/nói & Gọi LLM chấm câu theo rubric, lưu điểm và nhận xét. \\
    \hline
  \end{tabular}
  \caption{Các hàng đợi xử lý nền và vai trò của worker}
  \label{tab:async-jobs}
\end{table}

\paragraph{Lưu ý về mô hình ngôn ngữ.} Các lời gọi LLM được thực hiện qua API
\code{generateContent} của nền tảng \textbf{Google AI Studio}, dùng họ mô
hình mở \textbf{Gemma} của Google DeepMind~\cite{gemma2025technical} làm nền
tảng tham chiếu. Vì mô hình được cấu hình qua biến môi trường
(\code{GEMMA_MODEL}), thiết kế là \textit{độc lập mô hình}; trong cấu hình
triển khai mặc định, do dòng \code{gemma-3-*} đã ngừng phục vụ trên tầng miễn
phí của AI Studio, hệ thống dùng một mô hình Gemini Flash nhẹ tương đương.
Một chi tiết hiện thực quan trọng để tiết kiệm hạn mức: các lời gọi đặt
\code{thinkingConfig.thinkingBudget = 0} nhằm tắt token "suy nghĩ" ẩn ---
nếu không, mô hình suy luận sẽ tiêu hết hạn mức đầu ra cho phần suy nghĩ và
cắt cụt JSON kết quả.

\subsection{Làm giàu dữ liệu từ điển tự động bằng LLM}
\label{subsec:impl-enrich}

Đây là tính năng AI trung tâm: từ một \textit{lemma} đơn lẻ, hệ thống tự
dựng nên một bản ghi từ vựng đầy đủ (use case ở Bảng~\ref{tab:uc-quickcreate}).
\code{EnrichmentProcessor} chọn một trong hai đường đi tuỳ dữ liệu sẵn có:

\begin{itemize}
  \item \textbf{Đường có từ điển (tiếng Anh):} truy vấn Free Dictionary
        API~\cite{freedictionaryapi} để lấy từ loại, định nghĩa, phiên âm
        \textbf{IPA} và từ đồng/trái nghĩa; sau đó LLM chỉ cần bổ sung một
        \textit{gloss} ngắn, \textit{tối thiểu hai câu ví dụ} cho mỗi nghĩa,
        một bản dịch ngắn và ước lượng bậc CEFR. Cách kết hợp này tận dụng dữ
        liệu từ điển chính xác cho phần "khó tin tưởng AI" (định nghĩa, phiên
        âm) và chỉ giao cho AI phần nó làm tốt (ví dụ tự nhiên, bản dịch).
  \item \textbf{Đường chỉ dùng LLM (phi tiếng Anh hoặc từ điển không có):}
        LLM sinh toàn bộ cấu trúc nghĩa theo từng từ loại.
\end{itemize}

Vì Gemma trên AI Studio không hỗ trợ \code{responseSchema} hay vai trò
\textit{hệ thống} (system role), lược đồ JSON mong muốn được \textbf{nhúng
thẳng vào prompt} và phản hồi được \textbf{phân tích phòng thủ}. Bộ phân tích
bóc bỏ dấu rào ```\code{```json}```, dung thứ phần văn xuôi thừa, cắt và giới
hạn độ dài các trường, và --- quan trọng nhất --- \textit{từ chối} một CEFR
sai hoặc thiếu hay một nghĩa có dưới hai ví dụ bằng cách ném lỗi, để worker
\textit{thử lại} thay vì lưu dữ liệu rác. Đoạn dưới minh hoạ lược đồ được yêu
cầu và một phần kiểm tra phòng thủ:

\begin{verbatim}
// Lược đồ JSON nhúng trong prompt (rút gọn)
{ "cefr": "<A1|A2|B1|B2|C1|C2>",
  "senses": [ { "gloss": "<nhãn ngắn>",
                "examples": ["<câu>", "<câu>"],
                "translation": "<bản dịch ngắn>" } ] }

// src/vocabularies/enrichment/gemma-enricher.ts — kiểm tra phòng thủ
if (!VALID_CEFR.has(cefr))
  throw new Error(`enricher returned invalid cefr: ${cefr}`);
if (examples.length < 2)
  throw new Error(`sense ${i + 1} returned fewer than 2 examples`);
\end{verbatim}

\paragraph{Vòng kiểm soát chất lượng.} Kết quả do AI sinh ra luôn được kiểm soát theo
\textit{quyền sở hữu}, đúng tinh thần mục~\ref{subsec:llm-theory}: công việc do
\textit{quản trị viên} tạo sẽ sinh ra \textbf{bản nháp} thuộc kho hệ thống ở
trạng thái \code{is_approved = false} --- không xuất hiện trong phiên học của
ai cho tới khi được duyệt; còn từ do \textit{người học} tạo là từ cá nhân,
riêng tư và tự duyệt. Nhờ \textit{chỉ mục duy nhất bộ phận}
(mục~\ref{sec:hien-thuc-du-lieu}), worker bỏ qua mọi từ loại đã tồn tại nên
không bao giờ ghi đè dữ liệu sẵn có. Các từ được duyệt/tự duyệt tiếp tục kích
hoạt hàng đợi sinh âm thanh.

\subsection{Nhập từ vựng hàng loạt tự động}
\label{subsec:impl-bulk}

Tính năng nhập hàng loạt mở rộng cơ chế làm giàu ở quy mô lớn. Hàm
\code{extractCandidates} (\code{enrichment/import/}) nhận đầu vào ở nhiều định
dạng --- văn bản dán trực tiếp, CSV, bảng tính \textbf{XLSX} hay tệp
\textbf{PDF} --- quy về một khối văn bản rồi đưa qua bộ \textit{tách token}
chung. Bộ tách hỗ trợ hai chế độ: chế độ \textit{danh sách} (mỗi dòng/ô là một
mục) và chế độ \textit{văn xuôi} (tách từ, chuẩn hoá, loại bỏ \textit{từ dừng});
kết quả là danh sách \textit{ứng viên lemma}. Sau khi loại các từ đã có trong
kho, mỗi lemma còn lại được \textbf{phát tán} thành một công việc làm giàu
độc lập (chia sẻ một \code{batch_id}, và có thể kèm \code{target_deck_id} để
tự thêm từ mới vào một bộ thẻ đích sau khi tạo xong). Người dùng theo dõi
tiến độ cả lô qua mã lô thay vì phải chờ đồng bộ.

\subsection{Sinh âm thanh và hình ảnh minh hoạ tự động}
\label{subsec:impl-media}

Hai loại media được sinh tự động qua các pipeline có \textit{phương án dự
phòng} nhiều tầng, kết quả đều được lưu trên kho media \textbf{Cloudinary}~\cite{cloudinary2024}:

\begin{itemize}
  \item \textbf{Âm thanh phát âm:} ưu tiên lấy \textit{bản ghi âm người thật}
        từ Free Dictionary API; nếu không có, tổng hợp bằng \textbf{Microsoft
        Edge TTS} (giọng nơ-ron miễn phí, không cần khoá). Hàm trả về cả nguồn
        đã dùng (\code{dictionary} hay \code{tts}), giúp ưu tiên chất lượng cho
        các từ phổ biến và chỉ dùng TTS cho phần đuôi dài.
  \item \textbf{Hình ảnh minh hoạ nghĩa:} tìm ảnh phù hợp qua kho ảnh
        \textbf{Pexels}~\cite{pexels2024} rồi sao lưu lên Cloudinary. Với các
        từ trừu tượng không có ảnh phù hợp, pipeline trả về \code{null} và để
        trống cho quản trị viên bổ sung --- tự động hoá phần "đuôi dài" của
        các từ cụ thể mà không tạo ra ảnh sai lệch.
\end{itemize}

\subsection{Chấm điểm câu luyện tập bằng LLM}
\label{subsec:impl-judge}

Tính năng luyện viết câu (use case ở Bảng~\ref{tab:uc-practice}) dùng LLM làm
\textit{giám khảo}. Khi người học nộp một câu chứa từ mục tiêu,
\code{PracticeService} kiểm tra hạn mức ngày, tạo bản ghi \code{pending}, đẩy
vào hàng đợi và trả \code{202} ngay; nếu việc đẩy hàng đợi thất bại, bản ghi
được \textit{hoàn tác} để không "kẹt" ở trạng thái chờ và không tiêu phí hạn
mức. Worker \code{ScoringProcessor} gọi hàm \code{scoreSentence}
(\code{gemma-judge.ts}) với một prompt yêu cầu chấm theo \textit{rubric} có
cấu trúc (Bảng~\ref{tab:rubric-fields}); nhiệt độ sinh đặt thấp (0.2) để kết
quả ổn định. Điểm đặc thù: mô hình được yêu cầu ước lượng
\textbf{bậc CEFR mà chính câu đó thể hiện} (không phải năng lực tổng thể của
người học), cùng một câu gợi ý sửa.

\begin{table}[h]
  \centering
  \small
  \renewcommand{\arraystretch}{1.3}
  \begin{tabular}{|p{3.6cm}|p{3.0cm}|p{7.4cm}|}
    \hline
    \textbf{Trường rubric} & \textbf{Kiểu} & \textbf{Ý nghĩa} \\
    \hline
    \code{overall} & số nguyên 0--100 & Điểm tổng của câu. \\
    \hline
    \code{usesTargetWord} & boolean & Câu có dùng đúng từ mục tiêu (kể cả dạng biến cách). \\
    \hline
    \code{correctUsage} & boolean & Từ được dùng đúng một trong các nghĩa dự kiến. \\
    \hline
    \code{criteria} & 4 tiêu chí 0--5 & Ngữ pháp, dùng từ, độ tự nhiên, độ liên quan. \\
    \hline
    \code{cefr} & A1--C2 & Bậc CEFR mà câu thể hiện. \\
    \hline
    \code{feedback} & chuỗi & Một, hai câu nhận xét cho người học. \\
    \hline
    \code{correctedSentence} & chuỗi, tuỳ chọn & Phiên bản câu được cải thiện (nếu cần). \\
    \hline
  \end{tabular}
  \caption{Cấu trúc rubric chấm điểm câu luyện tập do LLM trả về}
  \label{tab:rubric-fields}
\end{table}

Phản hồi của mô hình được phân tích phòng thủ giống như ở phần làm giàu: các
khoảng số bị kẹp về biên hợp lệ, CEFR sai bị từ chối để worker thử lại. Người
học thăm dò kết quả qua mã bài luyện.

\subsection{Chấm điểm phát âm qua vi dịch vụ}
\label{subsec:impl-pron}

Chức năng luyện phát âm (Hình~\ref{fig:seq-pron}) được hiện thực bằng cách
tách phần xử lý tín hiệu/học máy thành một \textbf{vi dịch vụ} riêng viết bằng
Python (\textbf{FastAPI})~\cite{fastapi2024}, còn backend chỉ đóng vai
\textit{proxy mỏng}. \code{PronunciationScoringClient} đóng gói bản ghi âm
của người học thành một yêu cầu \code{multipart/form-data} (trường
\code{audio} và \code{word}), gọi điểm cuối \code{/score} với
\textit{thời gian chờ} ràng buộc, rồi ánh xạ lỗi của dịch vụ ngoài về mã lỗi
HTTP có ngữ nghĩa: \code{422} (âm thanh quá ngắn / không chấm được)
$\rightarrow$ \code{400}, còn không kết nối được / lỗi máy chủ $\rightarrow$
\code{503} (suy giảm có kiểm soát, NFR-04). Kết quả --- điểm tổng và điểm từng
âm vị --- được lưu vào \code{pronunciation_attempts} làm lịch sử luyện tập.
Việc cô lập này giữ cho tiến trình backend chính không phải mang các phụ thuộc
nặng về xử lý âm thanh.

% =====================================================================
\section{Theo dõi tiến độ và tạo động lực tự động}
\label{sec:tien-do}

Phân hệ tiến độ biến nhật ký học thành các chỉ số động lực một cách tự động
(các yêu cầu FR-19, FR-20). Nguồn dữ liệu là bảng \textit{chỉ-thêm}
\code{learning_activity} --- mỗi lượt ôn ghi đúng một dòng, được ghi cùng giao
dịch với cập nhật SRS (mục~\ref{subsec:impl-hmac}) nên các con số luôn nhất
quán với nhau. \code{ProgressService} tính:

\begin{itemize}
  \item \textbf{Thống kê trang chủ:} số từ theo từng trạng thái
        (\code{new}/\code{learning}/\code{review}/\code{mastered}), số thẻ đang
        đến hạn, số lượt ôn hôm nay và thời điểm đến hạn kế tiếp.
  \item \textbf{Chuỗi ngày học (streak):} số ngày liên tiếp (theo lịch UTC) có
        ít nhất một lượt ôn, tính lùi từ ngày ôn gần nhất; chuỗi chỉ còn hiệu
        lực nếu ngày gần nhất là hôm nay hoặc hôm qua.
  \item \textbf{Biểu đồ hoạt động (heatmap):} gộp số lượt ôn theo từng ngày
        \textit{theo múi giờ của người dùng} bằng một truy vấn SQL dùng
        \code{AT TIME ZONE}, để các ô lịch khớp với nửa đêm địa phương của họ.
\end{itemize}

\textbf{Bảng xếp hạng} được \code{LeaderboardService} tính trực tiếp: xếp hạng
người học theo \textit{số từ đã thành thạo} (toàn thời gian), chỉ tính những
tài khoản người học thật, đang hoạt động và \textit{không chọn ẩn}
(\code{leaderboard_opt_out}); thứ hạng của bản thân người gọi được suy ra từ
cùng một tập đã sắp xếp để hai con số không bao giờ mâu thuẫn. Bảng xếp hạng
theo \textit{số từ mới} được dành cho giai đoạn phát triển tiếp theo.

% =====================================================================
\section{Hiện thực hoá tầng giao diện người dùng}
\label{sec:hien-thuc-frontend}

Toàn bộ các API trình bày ở những mục trên được tiêu thụ bởi một ứng dụng giao
diện người dùng \textbf{độc lập}, hiện thực hoá thiết kế giao diện đã phác thảo
ở mục~\ref{sec:thiet-ke-giao-dien} (Chương~\ref{chuong:phan-tich-thiet-ke}).
Ứng dụng được xây dựng bằng \textbf{Next.js~16} theo mô hình \textit{App
Router} với \textbf{React~19} và \textbf{TypeScript} (mục~\ref{sec:cong-nghe-frontend}).
Đặc trưng kiến trúc xuyên suốt là \textit{ưu tiên phía máy chủ}: phần lớn màn
hình là Server Components lấy dữ liệu ngay trên máy chủ, còn mọi thao tác ghi
đi qua \textit{Server Actions}; nhờ đó access token và refresh token không bao
giờ lộ ra mã JavaScript phía trình duyệt. Bảng~\ref{tab:frontend-layout} ánh xạ
các phần chính của ứng dụng tới tệp hiện thực then chốt.

\begin{table}[h]
  \centering
  \small
  \renewcommand{\arraystretch}{1.3}
  \begin{tabular}{|p{3.6cm}|p{4.2cm}|p{6.0cm}|}
    \hline
    \textbf{Phần} & \textbf{Thư mục / tệp} & \textbf{Vai trò} \\
    \hline
    Nhóm tuyến & \code{src/app/(auth)}, \code{(app)}, \code{(admin)} & Phân tách khu vực xác thực, ứng dụng học và quản trị bằng \textit{route group}. \\
    \hline
    Tầng truy cập API & \code{src/lib/api.ts} & Client HTTP chạy phía máy chủ: \code{apiRequest} và \code{authedRequest}. \\
    \hline
    Quản lý phiên & \code{src/lib/auth/} & Lưu/đọc token trong cookie httpOnly, xoay token, đọc hồ sơ người dùng. \\
    \hline
    Kiểm tra dữ liệu & \code{src/lib/validations/} & Lược đồ \code{zod} dùng chung cho biểu mẫu và Server Action. \\
    \hline
    Động cơ phiên học & \code{src/app/(app)/learn/} & Bộ rút gọn \code{session-machine.ts} và các thành phần câu hỏi. \\
    \hline
  \end{tabular}
  \caption{Ánh xạ các phần frontend tới tệp hiện thực then chốt}
  \label{tab:frontend-layout}
\end{table}

\subsection{Cấu trúc App Router và phân nhóm tuyến}
\label{subsec:fe-approuter}

Cây tuyến được tổ chức bằng \textit{nhóm tuyến} (route group --- thư mục đặt
trong ngoặc đơn, không xuất hiện trong URL) để mỗi khu vực có bố cục và lớp bảo
vệ riêng: \code{(auth)} cho đăng nhập/đăng ký/xác minh email, \code{(app)} cho
trải nghiệm học của người dùng (trang chủ, học, khám phá, luyện tập, cộng đồng,
hồ sơ) và \code{(admin)} cho khu vực quản trị từ vựng. Mỗi phân đoạn tận dụng
các tệp đặc biệt của App Router --- \code{layout.tsx} (bố cục dùng chung),
\code{loading.tsx} (trạng thái chờ qua \textit{Suspense}), \code{error.tsx}
(ranh giới lỗi cấp phân đoạn) và \code{not-found.tsx} --- nên một lỗi cục bộ
không làm trắng toàn bộ ứng dụng. Các thành phần chỉ phục vụ một tuyến được đặt
ngay trong thư mục tuyến đó, chỉ nâng lên \code{src/components/} khi dùng lại ở
nhiều nơi.

\subsection{Tầng truy cập API phía máy chủ}
\label{subsec:fe-api}

Mọi lời gọi tới backend đi qua một client HTTP \textit{chỉ chạy phía máy chủ}
trong \code{src/lib/api.ts}. Địa chỉ backend (\code{API_BASE_URL}) được đọc từ
biến môi trường phía máy chủ, cố tình không gắn tiền tố \code{NEXT_PUBLIC_}
để không rò rỉ ra phía máy khách. Module phơi bày hai điểm vào:

\begin{itemize}
  \item \code{apiRequest} --- gọi \textit{không cần xác thực} (đăng nhập, đăng
        ký, làm mới token); tự đặt \code{Content-Type} JSON hoặc để runtime tự
        đặt biên \textit{multipart} khi thân là \code{FormData} (tải tệp lên).
  \item \code{authedRequest} --- đính kèm access token và, khi gặp
        \code{401}, \textit{tự động} làm mới refresh token đúng một lần rồi thử
        lại; chỉ gọi từ Server Action hoặc Route Handler vì nó có thể ghi lại
        cookie token đã xoay.
\end{itemize}

Cả hai chuẩn hoá thân lỗi của NestJS thành một kiểu \code{ApiResult} thống nhất
(cờ \code{ok}, mã trạng thái, dữ liệu, thân lỗi), nhờ đó tầng gọi xử lý lỗi
theo \textit{giá trị trả về} thay vì ngoại lệ --- khớp với triết lý phân biệt
lỗi \textit{dự kiến} và \textit{không dự kiến} của App Router.

\subsection{Quản lý phiên và xác thực phía máy khách}
\label{subsec:fe-session}

Token phiên được lưu trong \textbf{cookie httpOnly} với vòng đời bám đúng hợp
đồng backend (mục~\ref{sec:hien-thuc-auth}): access token 15~phút, refresh
token 30~ngày và được xoay sau mỗi lần gọi \code{/v1/auth/refresh}. Vì các
helper phiên import \code{next/headers}, chúng \textit{chỉ biên dịch được phía
máy chủ} --- một lớp bảo vệ ở mức biên dịch ngăn token lọt vào Client Component.
Khi \code{authedRequest} gặp \code{401}, nó âm thầm đổi refresh token lấy cặp
token mới, ghi lại cookie và thử lại đúng một lần; nếu vẫn thất bại, phiên bị
xoá và người dùng được chuyển về trang đăng nhập. Sau khi xác thực thành công,
Server Action điều hướng người dùng theo vai trò và trạng thái khai báo hồ sơ:
quản trị viên tới khu vực \code{/admin}, người học đã khai báo tới
\code{/dashboard}, người học chưa khai báo tới luồng \code{/onboarding}.

\subsection{Động cơ phiên học phía máy khách}
\label{subsec:fe-learn}

Phần tương tác nhất của ứng dụng là vòng lặp phiên học. Hàng đợi câu hỏi được
điều phối bằng một \textbf{bộ rút gọn (reducer) thuần} trong
\code{session-machine.ts}, tách hẳn khỏi React và mạng nên kiểm thử đơn vị được
độc lập. Nó nắm danh sách câu hỏi còn lại (câu đang hiển thị luôn là
\code{queue[0]}), các bộ đếm đúng/đã trả lời và trạng thái kết thúc
(\code{active}, \code{empty}, \code{done}, \code{expired}, \code{error}). Khi
một thẻ cần lặp lại trong cùng phiên, nó được \textit{nối lại vào cuối} hàng
đợi thay vì hẹn giờ theo đồng hồ thực. Bộ rút gọn này được ghép với hai Server
Action --- \code{startSessionAction} và \code{submitAnswerAction} --- gọi tới
\code{/v1/me/learn/session} và \code{/v1/me/learn/answer}. Vì máy chủ \textit{ký
HMAC} từng câu hỏi (mục~\ref{subsec:impl-hmac}), một \code{401} còn lại sau khi
đã làm mới token được hiểu là chữ ký câu hỏi đã hết hạn hoặc bị giả mạo, và giao
diện chuyển sang trạng thái \code{expired} mời người học bắt đầu phiên mới. Mỗi
\textit{chế độ} câu hỏi (thẻ ghi nhớ, điền khuyết trắc nghiệm/gõ chữ, nghe chép
chính tả, chọn theo hình/âm thanh, luyện phát âm\ldots) là một thành phần riêng
dưới \code{learn/questions/}, nhận dữ liệu đã ký từ máy chủ và chỉ chịu trách
nhiệm thu thập câu trả lời để nộp lại.

% =====================================================================
\section{Kiểm thử hệ thống}
\label{sec:kiem-thu}

Việc kiểm thử được hỗ trợ trực tiếp bởi một lựa chọn thiết kế lặp lại xuyên
suốt chương: \textbf{tách logic thông minh thành các hàm thuần, độc lập
framework}. Các hàm như \code{applySm2} (SM-2), bộ ký HMAC, bộ chấm đáp án,
bộ dựng prompt và bộ phân tích phản hồi LLM đều không phụ thuộc NestJS hay cơ
sở dữ liệu, nên có thể kiểm thử đơn vị nhanh và tất định. Chiến lược kiểm thử
gồm hai mức (Bảng~\ref{tab:test-targets}):

\begin{itemize}
  \item \textbf{Kiểm thử đơn vị} bằng \textbf{Jest}~\cite{jest2024}: tập trung
        vào các thuật toán và bộ phân tích --- nơi lỗi tinh vi dễ lọt và hậu
        quả lớn (lập lịch sai, lưu dữ liệu AI rác, chấm điểm sai).
  \item \textbf{Kiểm thử tích hợp (end-to-end)} bằng \textbf{Supertest}: gọi
        thật vào các điểm cuối API để kiểm tra luồng xác thực, kiểm tra đầu
        vào và mã trạng thái trả về.
\end{itemize}

\begin{table}[h]
  \centering
  \small
  \renewcommand{\arraystretch}{1.3}
  \begin{tabular}{|p{5.0cm}|p{9.4cm}|}
    \hline
    \textbf{Thành phần} & \textbf{Trọng tâm kiểm thử} \\
    \hline
    \code{srs.ts} & Chuyển trạng thái SM-2, bước học, tốt nghiệp, rớt bước, ngưỡng thành thạo. \\
    \hline
    \code{hmac-signer.service.ts} & Ký/xác minh, chống giả mạo, hết hạn, so sánh thời gian hằng số. \\
    \hline
    \code{answer-grader.service.ts} & Quy đổi điểm chất lượng cho từng loại câu hỏi. \\
    \hline
    \code{question-builder.service.ts} & Dựng bài học đúng dải khó, tôn trọng các trần và lọc dữ liệu. \\
    \hline
    \code{gemma-judge.ts}, \code{gemma-enricher.ts} & Phân tích phòng thủ JSON, kẹp số, từ chối CEFR sai. \\
    \hline
    \code{import/tokenize.ts}, \code{normalize.ts}, \code{xlsx-parser.ts} & Tách lemma, chuẩn hoá, đọc bảng tính cho nhập hàng loạt. \\
    \hline
  \end{tabular}
  \caption{Một số thành phần tiêu biểu và trọng tâm kiểm thử}
  \label{tab:test-targets}
\end{table}

Chất lượng mã nguồn còn được kiểm soát ở mức tĩnh bằng trình kiểm tra kiểu của
TypeScript (chế độ \textit{strict}), \textbf{ESLint} và \textbf{Prettier},
chạy trước mỗi lần commit (đáp ứng NFR-06).

% =====================================================================
\section{Kết chương}
\label{sec:ket-chuong-4}

Chương này đã trình bày quá trình hiện thực hoá hệ thống, với trọng tâm là các
thành phần thông minh và tự động hoá. Ở phía \textit{động cơ học}, hệ thống
thích nghi với từng người học qua thuật toán SM-2 mở rộng bước học
(\code{srs.ts}), cơ chế chọn từ theo khoảng cách CEFR, bộ sinh câu hỏi nhiều
tầng theo giai đoạn ghi nhớ với phương án nhiễu "dễ nhầm", và cơ chế chấm điểm
phi trạng thái an toàn bằng chữ ký HMAC. Ở phía \textit{nội dung}, các tính
năng AI --- làm giàu từ điển bằng LLM kết hợp từ điển, nhập hàng loạt từ nhiều
định dạng, sinh âm thanh và hình ảnh tự động, chấm điểm câu và phát âm --- đều
chạy trên một hạ tầng xử lý nền bất đồng bộ có tính idempotent, biết thử lại,
tự giới hạn tần suất và suy giảm có kiểm soát. Ở phía \textit{máy khách}, một
ứng dụng Next.js theo App Router tiêu thụ toàn bộ các API trên qua một tầng
truy cập phía máy chủ, quản lý phiên bằng cookie httpOnly với cơ chế xoay token
tự động, và điều phối phiên học bằng một bộ rút gọn thuần tách khỏi React. Tất
cả được củng cố bằng các chỉ số tiến độ tự động và một chiến lược kiểm thử dựa
trên việc tách logic thành các hàm thuần. Những kết quả hiện thực hoá này là cơ
sở để chương tiếp theo đánh giá hệ thống và rút ra kết luận.
